\section{Geometry, motives, and the Grothendieck programme}
\label{sec:geometry}

The letters A and M have the dependency meaning fixed in the introduction.
In particular, accepting the stipulated endpoint A is not a certification of
every intermediate statement in the note, and assuming M does not certify
the functional analysis in its proposed proof. The findings below are new
formulations proved in this run; no claim of priority over the literature is
intended. Earlier findings are used with their original identifiers.

\subsection{The transfer results already established}
\label{ag:transfer}

\paragraph{AG-AT (the exact Artin--Tate domain).}
For a finite-type scheme over $\mathbb Z$ with exact all-prime point counts
\[
 \#X(\mathbb F_{p^r})=
 \sum_j p^{jr}\operatorname{Tr}(F_p^r\mid V_j^{I_p}),
\]
where the $V_j$ are finitely many virtual finite-image representations,
the identity
\[
 \zeta_X(s)=\prod_j L(V_j,s-j)
\]
is unconditional. This is ZETA\_TWO ZT01, in
\texttt{RETURN\_ZETA\_TWO\_final\_IMPLICATIONS.md}, Section 1,
extending BRIDGES F1. If the representations have abelian image, A gives
noninteger global divisor support on the lines $\Re s=j+1/2$.
AG-ARTIN records the arbitrary finite-image case: M gives the corresponding meromorphic
divisor-support assertion through integer Brauer induction (ZT02).
Negative exponents can produce poles; Artin holomorphy is an extra question.

\paragraph{AG-DED and AG-LOCAL (ramification and models).}
The permutation representation on the embeddings of a number field has
inertia invariants giving the Euler factors of $\operatorname{Spec}\mathcal
O_K$ at ramified primes as well. Thus every Dedekind zeta function is covered
by M, and abelian fields over $\mathbb Q$ are covered by A (ZT03).
If a factorization holds only away from a finite set $S$, the actual factors
at $S$ must be restored (ZT04). The counterexample ZT05, catalogued as AG-MODEL,
\[
 \zeta_{\operatorname{Spec}\mathbb Z[1/p]}(s)
       =(1-p^{-s})\zeta(s),
\]
has additional zeros $2\pi i k/\log p$, $k\ne0$. Such local divisor lines
are not shifted global critical lines. This is a statement about the
integral model, not a failure of the generic-fibre motivic normalization.

\paragraph{AG-CM (the potentially CM domain).}
ZETA\_TWO CM01 supplies the all-weight extension to pure objects in the
absolute-Hodge category generated by Artin motives, Tate motives, and
potentially CM abelian varieties. Its cited CM descent input gives a signed
product of Hecke functions of one weight; norm twisting makes these
characters unitary. The required M therefore includes infinite-order
unitary characters. The result is meromorphic divisor RH, and does not
declare absolute-Hodge projectors algebraic. The exact prior locator is
\texttt{RETURN\_ZETA\_TWO\_final\_IMPLICATIONS.md}, Section 2, CM01,
with its reported Anderson Section 5.7, Corollary 5.7.2 and Remark 5.7.3.
The original source has now been recovered and checked: G.~W.~Anderson,
\emph{Cyclotomy and an extension of the Taniyama group}, Compositio
Math.~57 (1986), 153--217. The category is defined in Section 5.3,
printed p.~178 (PDF p.~27); Theorem 6, Section 5.6, printed pp.~179--180
(PDF pp.~28--29), identifies its one-dimensional characters with algebraic
Hecke characters, including the inertia and Frobenius comparisons.
Corollary 5.7.2 and Remark 5.7.3 are both on printed p.~181 (PDF p.~30).
The former gives integer exponents, with no positivity claim; the latter
gives the common weight. Its internal reference to ``Corollary 5.8.2'' is
a printed cross-reference error, not a new result used here.

\paragraph{AG-CM-REC and AG-CM-PERIOD (classical conclusions already
stronger than zero location).}
In this precise potentially CM absolute-Hodge category, equality of the
whole $L$-functions implies isomorphism of objects: Anderson,
Proposition 5.7.5, printed pp.~181--182 (PDF pp.~30--31). This is
unconditional. The proof uses semisimplicity and Corollary 5.7.4's
identification, for a weight-zero object, of the pole order at $1$ with
the dimension of its invariant part. Equality of the local factors makes
the tensor products with a dual agree away from finitely many primes;
their local factors are nonzero and nonpolar at $1$. A pole of
$L(M\otimes M^\vee,s)$ therefore gives a nonzero morphism
$M\longrightarrow N$, and induction on the common rank proves the
claim. The proof was read in full, rather than inferred from GRH.
This is a recognition theorem for semisimple motives; it supplies no
effective finite bound for general infinite CM monodromy.

Anderson's Section 5.8, Theorem 7, printed p.~182 (PDF p.~31), also
proves Deligne's period conclusion
$L(M,0)/c^+(M)\in\mathbb Q$ for the stated \emph{tractible, critical}
objects: those for which the $L$-function is a single algebraic Hecke
$L$-function over a totally real or CM field. This is a separate
classical special-value theorem, using Siegel and Blasius, and is not
enlarged to arbitrary motives by A or M. It complements, rather than
replaces, the cycle/height and integral comparison gates below.

\paragraph{AG-WEIGHT and AG-NONCM (the residual geometric factors).}
With geometric Frobenius,
\[
 \operatorname{wt}N(n)=\operatorname{wt}N-2n,
 \qquad L(N(n),s)=L(N,s+n).
\]
The line for weight $w$ is $\Re s=(w+1)/2$. Positive weight is not the
remaining obstruction: CM objects of positive weight are covered by M,
whereas AG-NONCM-W0 records that $\operatorname{Sym}^2h^1(E)(1)$ for a non-CM elliptic curve still
has a non-CM factor despite having weight zero. AG-SCHEMES carries ZT06: it reduces a finite-type
scheme to finitely many smooth projective geometric factors and actual
exceptional fibres. ZT07 proves that, under A, an off-line elliptic zero
in $1/2<\Re s<3/2$ cannot be cancelled by the two Tate factors in the
good-model formula
\[
 \zeta_{\mathcal E}(s)
   =\zeta^S(s)\zeta^S(s-1)/L^S(E,s).
\]
These are retained reductions, not a proof of RH for the non-CM factor.
Framework locators carried from ZETA\_TWO are SGA 4 XVII Sections 5.1--5.3;
SGA 5 XV Section 3, no.~2, unnumbered theorem, equation (1);
EGA IV$_3$ (8.8.2); and Deligne, Weil I, Theorem 1.6, and Weil II,
Theorem 3.3.1.

\subsection{AG-FROB: a finite recognition theorem for Artin--Tate objects}
\label{ag:frob}

Let $L/\mathbb Q$ be a nontrivial finite Galois extension, with group $G$,
absolute discriminant $D_L$, and degree $n_L$. Put
\[
 B_L=(4\log D_L+\tfrac52 n_L+5)^2.
\]
Let
\[
 U=\bigoplus_j V_j(-j),\qquad W=\bigoplus_j W_j(-j),
\]
where all sums are finite and $V_j,W_j$ are actual finite-dimensional
characteristic-zero representations of $G$. Their Euler polynomials at
unramified primes are understood as those of the associated compatible
Artin--Tate systems; a single auxiliary prime may instead be chosen greater
than $B_L$.

\textbf{Statement.} Assume RH for $\zeta_L$ and $\zeta_{\mathbb Q}$.
If
\[
 \det(1-TF_p\mid U)=\det(1-TF_p\mid W)
 \quad\text{for every }p\le B_L\text{ unramified in }L,
\]
then $V_j\simeq W_j$ for every $j$, and hence $U\simeq W$.
For rational coefficients this is also an isomorphism in the rational
Artin--Tate category. A supplies the hypothesis when $L/\mathbb Q$ is
abelian; M supplies it in general. When $L=\mathbb Q$, the one prime $p=2$
suffices.

\begin{proof}
At an unramified prime, every eigenvalue in $V_j(-j)$ has the form
$p^j\omega$, where $\omega$ is a root of unity. Therefore its complex
absolute value is $p^j$. Distinct integers $j$ give distinct absolute
values, so equality of the total Euler polynomials separates uniquely into
equality of the polynomials in each weight. It follows that
$\operatorname{Tr}(F_p\mid V_j)=\operatorname{Tr}(F_p\mid W_j)$ for every
$j$ and every prime in the hypothesis.

Bach--Sorenson, \emph{Explicit bounds for primes in residue classes},
Math. Comp. 65 (1996), Theorem 5.1, printed p.~1729 (PDF p.~13), gives
an unramified prime $p\le B_L$ in every conjugacy class of $G$.
The precise analytic assumptions are stated in their Section 2,
printed p.~1720: zero-freeness of $\zeta_L$ and $\zeta_{\mathbb Q}$ in
$\Re s>1/2$. The reduction from cyclic to general Galois extensions is
Corollary 3.3, p.~1721. Thus the characters of $V_j$ and $W_j$ agree on
every conjugacy class. In characteristic zero finite-group representations
are semisimple, and equal characters give isomorphic representations.

For completeness, this last assertion also descends to a non-algebraically
closed characteristic-zero coefficient field $E$. The space of
$G$-equivariant maps is the solution space over $E$ of finitely many linear
equations. If it contains an invertible map after scalar extension, its
determinant polynomial is nonzero on that $E$-vector space. Since $E$ is
infinite it has a nonzero value at an $E$-point. That point is an invertible
$G$-equivariant map over $E$. For $L=\mathbb Q$, the eigenvalues $2^j$
already recover all multiplicities from the polynomial at $2$.
\end{proof}

In Grothendieck's language this is effective recognition inside a specified
semisimple Artin--Tate realization category. It extends
ALPOGE\_BRIDGES AQ-FINITE, whose sharper abelian conductor bound remains
available. It does not construct algebraic projectors for an arbitrary
geometric realization. Moreover, recognition of an Artin motive is weaker
than recognition of a finite cover as a scheme; the connected Gassmann
example in CX-GEO-004's companion finding makes this distinction explicit.
The forbidden-prime variant is handled there by a quadratic compositum.

\paragraph{The elliptic $h^1$ recognition extension (NT-E16).}
There is also an effective recognition statement beyond finite
monodromy, with the geometry supplied by Faltings's theorem.
For any two elliptic curves over $\mathbb Q$, including CM
curves, given by nonsingular integral Weierstrass equations,
put $b=\operatorname{bitlength}(2|\Delta_1\Delta_2|)$.
NT-E16 proves under M that equality of the point counts at
all primes $p\le(1532b+10181)^2$ not dividing
$2\Delta_1\Delta_2$ is equivalent to $\mathbb Q$-isogeny.
It uses only the Dedekind RH instances for the finite
Galois detection fields of degree at most $384$, together
with $\zeta_{\mathbb Q}$; it assumes no RH for the standard
elliptic $L$-function.  The complete proof also gives a
deterministic polynomial bit-time decision from those
equations without factoring their discriminants.
Thus the closed points of the actual good-reduction open
of $\operatorname{Spec}\mathbb Z$ recognize the rational
$h^1$ isogeny motive in this family.  An isogeny and its
dual give inverse rational correspondences after division
by their degree.  Neither isomorphism of the integral
curves nor an integral inverse correspondence is asserted.
This is the precisely scoped elliptic extension of the
recognition discussion in AG-FROB and AG-CM-REC.

\paragraph{AG-COVER: equal Artin motives need not give equal covers.}
Put $L=\mathbb Q(\sqrt2,\sqrt3)$ and let $K_1,K_2,K_3$ be its three
quadratic subfields. The finite schemes
\[
 X=\operatorname{Spec}L\amalg\operatorname{Spec}\mathbb Q
       \amalg\operatorname{Spec}\mathbb Q,
 \qquad
 Y=\coprod_{i=1}^3\operatorname{Spec}K_i
\]
are not isomorphic, but their rational Artin motives, and all their local
Artin Euler factors, are equal. Consequently
\[
 \zeta_L(s)\zeta_{\mathbb Q}(s)^2
       =\zeta_{K_1}(s)\zeta_{K_2}(s)\zeta_{K_3}(s).
\]
This identity includes the factors at ramified primes.

\begin{proof}
Write $G=\operatorname{Gal}(L/\mathbb Q)=C_2\times C_2$, with its three
subgroups $H_i$ of order two. The character of
$\mathbb Q[G]\oplus\mathbb Q^2$ takes value $6$ at the identity and
$2$ at every other element. Each nonidentity element belongs to exactly
one $H_i$, so the character of $\bigoplus_i\mathbb Q[G/H_i]$ has the
same values. Characteristic-zero semisimplicity and the descent argument
in AG-FROB give an isomorphism of rational representations. Isomorphic
representations have isomorphic inertia invariants with their residual
Frobenius action, proving equality of every local factor and hence the
zeta identity in the initial half-plane and by continuation. On the other
hand the degrees of the connected components are $4,1,1$ for $X$ and
$2,2,2$ for $Y$, precluding a scheme isomorphism.
\end{proof}

The same example over $\mathbb Z$ uses the respective rings of integers;
over $\mathbb Z[1/6]$ these are finite \emph{\'etale} covers. Thus even
complete Euler data forget the multiplication on an Artin algebra. A
applies to all fields here because they are abelian over $\mathbb Q$,
but does not restore the forgotten multiplication. The independent
transitive seven-element example CX-GEO-005 treats the connected
$G$-set distinction without asserting an arithmetic realization.

\subsection{AG-LAMBDA: the precise closure supplied by tensor operations}
\label{ag:lambda}

\textbf{Statement.} The finite direct sums of abelian Artin characters with
integer Tate twists are closed under tensor products, duals, exterior and
symmetric powers, and direct summands in characteristic zero. All their
$L$-functions have the A divisor-support conclusion. Replacing abelian
finite image by arbitrary finite image gives the corresponding M conclusion.
The same operations preserve the potentially CM absolute-Hodge category
used in CM01, so its M conclusion applies there as well.

\begin{proof}
Over an algebraic closure of the coefficient field an abelian Artin--Tate
object is a sum of lines $\chi_a(-j_a)$. A tensor product or a monomial
appearing in a symmetric or exterior power is a line
$(\prod_a\chi_a^{e_a})(-\sum_a e_aj_a)$. The character still has finite
abelian image and the Tate index is an integer. Duality replaces the
character and index by their inverses and negatives. Direct summands remain
in the same semisimple category, and descent preserves the asserted
factorization. The nonabelian assertion follows because all these
operations still factor through a finite group and commute with Tate
twists; apply ZT02 to their resulting Artin representations. Finally, the
potentially CM category is the tensor, dual, and summand closure in CM01.
Symmetric and exterior powers are the images of the usual idempotents in
the symmetric-group action on tensor powers, so belong to it too.
\end{proof}

The placement is SGA 6 V Section 2 and Section 7, especially (7.1.1),
and the determinant/$\lambda$-ring discussion carried in ZETA\_TWO ZT08
and SV01. These closures do not generate an arbitrary non-CM $GL_2$
motive from CM objects; enlarging a list of tensor operations is not the
missing analytic theorem.

\subsection{AG-SIGN: central signs with the necessary normalization}
\label{ag:sign}

\textbf{Statement.} Let $F$ be a nonzero holomorphic function on a
neighbourhood of $[c,\infty)$, real-valued there. Suppose $F$ has no zeros
on $(c,\infty)$ and is positive at one sufficiently large real point.
If $r=\operatorname{ord}_{s=c}F$, then
\[
 \frac{F^{(r)}(c)}{r!}>0.
\]
If, moreover, $F(c+z)=\varepsilon F(c-z)$ near $0$, with
$\varepsilon\in\{1,-1\}$, then $(-1)^r=\varepsilon$.
For a real, standardly normalized, holomorphic Hecke or CM $L$-function,
the absence of real zeros to the right of the centre follows from M
together with its usual analytic continuation and Euler-product theory.
The analogous Dirichlet statement uses A whenever no pole is present.

\begin{proof}
A continuous real nonvanishing function has constant sign on the connected
interval $(c,\infty)$. The normalization makes that sign positive.
Write $F(c+t)=a_rt^r+O(t^{r+1})$, with $a_r\ne0$. For $t>0$ the
left-hand side is positive, so $a_r>0$ after division by $t^r$ and passage
to the limit. Comparing the leading terms in the functional equation gives
$a_r=\varepsilon(-1)^r a_r$, hence the parity assertion.
\end{proof}

\paragraph{AG-HEIGHT (what this proves for a supplied height formula).}
If $F(c)=0$, then $F'(c)\ge0$: it is positive when $r=1$ and zero when
$r>1$. Consequently, an independently established identity
$H(Z)=C F'(c)$ with a specified constant $C>0$ yields $H(Z)\ge0$.
The elementary deduction is complete; the existence, domain, and
normalization of the height identity are extra inputs. It does not assert
that every homologically trivial cycle has such a formula, or that its
$L$-function has simple central vanishing. For an identity involving the
first nonzero derivative, the sign is strictly positive. For a formula with
an additional sign, that sign must be retained.

\paragraph{AG-SIGN-NORM (a counterexample to the unnormalized inference).}
Zero locations and functional-equation parity alone do not determine the
sign: $F(s)=s-c$ and $-F(s)$ have exactly the same zeros and the same odd
functional equation, but their central derivatives have opposite signs.
This is a counterexample to the stated analytic inference, not an example
of two standardly normalized arithmetic $L$-functions.

These findings refine \texttt{OURS\_HODGE\_BRIDGE\_SKETCH\_20261121.md},
Section 3.5. Their archive placement is cote 44, scans 2, 4 and 6,
Theorems 4--5, and scans 8--9, the letter to Serre of 12 August 1964,
as reported in that attached sketch. No original scan has been inferred
from this sign argument. The remaining arithmetic Hodge-index gate is
the cycle/height comparison for the stated class of cycles, not a sign
convention suppressed in a formula.

\subsection{AG-CRITICAL: why off-centre critical integers gain little}
\label{ag:critical}

\textbf{Statement.} Suppose $N$ is a pure, integral-weight, explicitly
Hecke-factored motive of weight $w$, with its usual functional equation
and archimedean factors. At an endpoint of the unitary strip, require each
factor in the specified signed Hecke factorization (or its dual, when the
functional equation is used) to be nonpolar. This factorwise requirement
prevents a denominator pole from producing a zero. Let $m$ be critical in
Deligne's sense: $L_\infty(N,s)$ and $L_\infty(N^\vee,1-s)$
are both regular at $s=m$. If $m\ne(w+1)/2$, then the standard
zero-free half-plane and boundary for unitary Hecke $L$-functions,
together with the functional equation, already give $L(N,m)\ne0$.
M is not needed for this conclusion.

The precise unitary boundary theorem used here is
D.~Ramakrishnan--R.~J.~Valenza, \emph{Fourier Analysis on Number Fields}
(1999), Section 7.7, Theorem 7-28, printed p.~289: for every number
field and every unitary idele class character, there is no zero at
$1+it$. Section 7.4, Theorem 7-19, printed p.~279, supplies the completed
continuation and functional equation, with its pure norm-character pole
exceptions retained. These are classical inputs independent of M.

\begin{proof}
After unitary normalization set $z=m-w/2$. If $w$ is even, $z$ is an
integer; if $w$ is odd, $z$ is a half-integer. An integer $m$ with
$0<z<1$ therefore has $z=1/2$, and is precisely the excluded centre.
For every other $m$ one has $z\ge1$ or $z\le0$. In the first case
standard unitary Hecke nonvanishing in $\Re z\ge1$, away from a pole,
gives the result for each factor and hence its finite signed product.
In the second case the functional equation relates the completed value
to its dual at $1-z\ge1$. Criticality ensures that both archimedean
factors are finite; gamma functions have no zeros. The nonzero dual
value therefore gives a nonzero original value. The factorwise nonpolar
condition excludes a Tate pole or its inverse being misread as a finite
nonzero special value. For a product with nonnegative exponents, the
usual no-pole condition implies this condition.
\end{proof}

Thus the strongest new real-axis information from M lies between the
off-centre integer points and at signs of central leading coefficients;
it does not evaluate regulators, periods, central ranks, or integral
torsion. This is the special-value distinction in ZETA\_TWO Section 4
and in the Hodge sketch's edge-versus-centre table. The relevant motivic
conventions are Deligne, \emph{Valeurs de fonctions $L$ et p\'eriodes
d'int\'egrales}, Sections 1.1--1.3; the period conjectures numbered 1.8
and 2.8 are separate inputs.

\subsection{AG-DET: Hecke versions of the square-variable determinant}
\label{ag:det}

The following formulation records exactly which analytic data are needed.
Let $\Lambda$ be a nonzero entire function of order at most one, with
\[
 \Lambda(c+z)=\varepsilon\Lambda(c-z),\qquad
 \varepsilon\in\{1,-1\},
\]
and let $r=\operatorname{ord}_{s=c}\Lambda$. Assume that every other zero
is $c+i\gamma$, with $\gamma$ real and nonzero. Choose one positive
$\gamma$ from each pair, with multiplicity, and write the sequence as
$\gamma_1,\gamma_2,\ldots$ in nondecreasing order, so $\gamma_1$
denotes the least positive zero when the sequence is nonempty.
Define the removable entire function
\[
 H(t)=t^{-r}\Lambda(c+it),\qquad
 D(w)=H(\sqrt w)/H(0).
\]
No choice of square-root branch affects $D$.

\textbf{Statement.} The diagonal operator
$B=\operatorname{diag}(\gamma_j^{-2})$ is positive and trace class, and
\[
 D(w)=\prod_j(1-w\gamma_j^{-2})=\det(I-wB).
\]
In particular this applies under M to standard entire selfdual completed
Hecke functions and to the entire completed $L$-function of a CM elliptic
curve, with its centre $c=1$. The usual order, functional equation,
and completion are additional classical inputs, not assertions supplied
by the word ``GRH''. Poles must first be removed in a separately specified
entire completion; arbitrary meromorphic Artin quotients are not included.

\begin{proof}
Comparison of the first nonzero Taylor coefficients shows
$\varepsilon=(-1)^r$. Hence $H(-t)=H(t)$, so $D$ is entire by the even
power series for $H$. Its order is at most one half: a bound of order
$1+\delta$ for $H(t)$ gives a bound of order $(1+\delta)/2$ for $D(w)$.
Its zeros, with multiplicity, are exactly $\gamma_j^2$.
The convergence exponent of the zeros of an entire function of order at
most one is at most one. In particular
$\sum_j|\gamma_j|^{-2}<\infty$. Equivalently the diagonal positive
operator $B$ has finite trace norm. Hadamard factorization for $D$, whose
order is strictly less than one, has no nonconstant exponential factor
and uses genus-zero factors. Since $D(0)=1$, its constant factor is one.
This gives the displayed product. The Fredholm determinant of a diagonal
trace-class operator is that same absolutely locally uniformly convergent
product, proving the identity. An exact analytic locator for the
factorization used here is C.~McMullen, \emph{Advanced Complex Analysis},
Math 213a course notes, version 16 November 2026, Theorem 3.16,
printed p.~92 (PDF p.~93); Theorems 3.10--3.12 give the zero-exponent
bound used just above. This is classical background, not a new claim
dependent on that revision date.
\end{proof}

This extends the prior zeta realization BR-TR-01 to a precisely specified
Hecke domain. Its relationship to cote 1 scans 79--80 and 111--112 is
the trace/Fredholm theorem reported in
\texttt{OURS\_INTAKE\_RESEARCH\_SA23\_20261120.md}, Section 4.1,
geo-H-lidskii and geo-H-genus. Those sheets are not themselves an
arithmetic RH claim. The diagonal presentation here is defined from
the zeros. The further comparison MT-M29 identifies this very operator
with the inverse square of the positive-frequency norm generator on
the associated graded of the canonical locally finite module in the
actual adelic summation quotient. Its image-power filtration retains
every zero multiplicity while removing the nilpotent extensions.
The resulting metric uses the fixed Mellin coordinate and its
positivity still uses the stated A or M zero-location input.
Thus the particular arithmetic-module identification is now supplied;
an independently positive cohomology or GGP-period realization remains
a separate construction. The ordinary trace-GNS completion has only
one dimension per distinct zero and does not replace that associated
graded when multiplicities occur.

\subsection{AG-MOMENT and AG-HANKEL: recoverable positive spectral data}
\label{ag:moment}

Assume the hypotheses of AG-DET and infinitely many distinct positive
zeros. Put $\lambda_j=\gamma_j^{-2}$ and
\[
 s_k=\operatorname{Tr}(B^k)=\sum_j\lambda_j^k\quad(k\ge1).
\]
The logarithmic determinant expansion is
\[
 -\log D(w)=\sum_{k\ge1}\frac{s_k}{k}w^k
 \qquad (|w|<\gamma_1^2).
\]

\textbf{AG-MOMENT.} For every $k\ge1$,
\[
 s_k^{-1/k}\le\gamma_1^2\le\frac{s_k}{s_{k+1}}.
\]
The left bounds increase to $\gamma_1^2$; the right bounds decrease to it.
If $m_1$ is the multiplicity of the least positive zero, then
$\lim_{k\to\infty}s_k\gamma_1^{2k}=m_1$.

\begin{proof}
Let $a=\max_j\lambda_j=\gamma_1^{-2}$. Then
$s_k\ge a^k$ and $s_{k+1}\le a s_k$, proving the two inequalities.
The $\ell^p$ norm monotonicity gives
$s_{k+1}^{1/(k+1)}\le s_k^{1/k}$. One direct proof is to divide by
$s_k^{1/k}$: the resulting nonnegative numbers have $k$-th powers
summing to one, are each at most one, and thus have $(k+1)$-th powers
summing to at most one. Hence the left bounds increase.
Cauchy--Schwarz applied to
$\lambda_j^{k/2}$ and $\lambda_j^{(k+2)/2}$ gives
$s_{k+1}^2\le s_ks_{k+2}$, so the right bounds decrease.

There are finitely many indices with $\lambda_j=a$, since $\sum\lambda_j$
is finite. For $k\ge1$,
$0\le(\lambda_j/a)^k\le\lambda_j/a$, whose sum is finite.
Dominated convergence shows
$s_k/a^k\to m_1$. It follows that both $s_k^{1/k}$ and
$s_{k+1}/s_k$ tend to $a$, which proves both limits and the
multiplicity formula.
\end{proof}

\textbf{AG-HANKEL.} For all integers $h\ge1$ and $d\ge0$, the matrix
\[
 (s_{h+i+j})_{0\le i,j\le d}
\]
is positive definite. Its determinant is therefore strictly positive.

\begin{proof}
For a nonzero vector $a=(a_0,\ldots,a_d)$ and
$P(X)=\sum_{i=0}^d a_iX^i$, its Hermitian quadratic form equals
\[
 \sum_{j\ge1}\lambda_j^h|P(\lambda_j)|^2.
\]
This converges because the $\lambda_j$ are bounded and
$\sum\lambda_j^h<\infty$. All terms are nonnegative. A nonzero
polynomial has only finitely many distinct zeros, whereas the positive
$\lambda_j$ have infinitely many distinct values. At least one term
is positive. This proves strict positive definiteness.
\end{proof}

These consequences strengthen the two-by-two moment inequalities carried
in BR-TR-02. They describe actual determinant coefficients available in
the cote 1 analytic framework, while leaving the arithmetic/geometric
realization gate unchanged. They assert no simplicity and no
random-matrix statistics.

\subsection{AG-SIGNED: a finite moment witness for a surviving pole}
\label{ag:signed}

Consider a finite signed product of specified entire selfdual functions
of order at most one, with a common centre $c$, each satisfying the
zero-location hypotheses of AG-DET. Let $r\in\mathbb Z$ be the net
central order. Divide out the central Laurent factor and normalize its
nonzero constant term. Its square-variable quotient has the form
\[
 D(w)=\prod_{\lambda\in\mathcal L}(1-w\lambda)^{m_\lambda},
 \qquad m_\lambda\in\mathbb Z\setminus\{0\},\quad
 \lambda>0,\quad \sum_\lambda |m_\lambda|\lambda<\infty,
 \qquad D(0)=1.                                      \tag{AG-S1}
\]
Here $\mathcal L$ has no positive accumulation point, and the signed
multiplicities are taken \emph{after} cancellation. The form (AG-S1)
follows by applying AG-DET to each factor. More generally the proposition
below applies whenever (AG-S1) is independently supplied. Put
\[
 s_k=\sum_\lambda m_\lambda\lambda^k\quad(k\ge1),\qquad
 H_d=(s_{1+i+j})_{0\le i,j\le d}.
\]
These moments are the local logarithmic coefficients
$-\log D(w)=\sum_{k\ge1}s_kw^k/k$ near $0$.

\textbf{Statement.} All the matrices $H_d$ are positive semidefinite
if and only if $m_\lambda>0$ for every $\lambda\in\mathcal L$.
Equivalently, the normalized quotient has no noncentral poles. The
central order condition $r\ge0$ must be checked separately. If a
negative multiplicity survives, a finite matrix $H_d$ already detects it,
with the explicit witness described below.

\begin{proof}
Set $a=\sup\mathcal L$ and introduce the finite signed measure
\[
 \mu=\sum_\lambda m_\lambda\lambda\,\delta_\lambda
 \quad\hbox{on }[0,a].
\]
The empty support case is immediate. For every complex polynomial
$P(X)=\sum_{i=0}^d b_iX^i$, absolute convergence gives
\[
 \sum_{i,j=0}^d\overline b_i s_{1+i+j}b_j
       =\int |P(x)|^2\,d\mu(x).
\]
Positive multiplicities therefore imply all the required inequalities.

Conversely, suppose $m_{\lambda_0}<0$. The point $\lambda_0>0$ is
isolated in the compact set $\{0\}\cup\mathcal L$. Let
\[
 \delta=\operatorname{dist}
   (\lambda_0,(\{0\}\cup\mathcal L)\setminus\{\lambda_0\})>0,
 \qquad R>a+\lambda_0,
 \qquad q=1-\delta^2/R^2<1.
\]
For each nonnegative integer $N$, the polynomial
\[
 P_N(x)=\left(1-\frac{(x-\lambda_0)^2}{R^2}\right)^N
\]
has degree at most $2N$, takes the value $1$ at $\lambda_0$, and has
$|P_N(\lambda)|\le q^N$ at every other point of $\mathcal L$.
Writing $W=\sum_{\lambda\ne\lambda_0}|m_\lambda|\lambda$, we obtain
\[
 \int |P_N|^2\,d\mu
   \le -|m_{\lambda_0}|\lambda_0+Wq^{2N}.
\]
Choose $N$ with $Wq^{2N}<|m_{\lambda_0}|\lambda_0$; take $N=0$
if $W=0$. Then the coefficient vector of $P_N$ is a negative
quadratic-form witness for the matrix $H_{2N}$, of size $2N+1$.
This proves the converse and its finite-witness assertion. In (AG-S1)
a negative net multiplicity is exactly a pole at $w=\lambda^{-1}$,
which proves the analytic equivalence. Removing the central Laurent
factor erased its order, explaining the separate condition on $r$.
\end{proof}

Geometrically this strengthens the cote 1 determinant bridge in the
meromorphic setting. A or M can first localize an identified divisor;
positivity of these finite moment matrices is then an exact additional
condition preventing cancellation from leaving poles. It is not supplied
by zero location alone. In particular, universal Hecke GRH plus Brauer
induction does not prove this positivity, and hence does not prove Artin
holomorphy. The criterion is for the specified quotient: symmetrizing
an arbitrary nonselfdual function with its dual can cancel a pole between
the two factors, so holomorphy of that product must not be attributed to
each unsymmetrized factor. The local character-matrix test in MT-ART02
is a complementary criterion at an individual Artin divisor point.

\subsection{AG-EXT: the precise information that the determinant forgets}
\label{ag:ext}

The earlier findings BR-MIX-01 and ZETA\_TWO SV01--SV02 are retained.
At a good prime an adapted matrix for a mixed object is block triangular,
so its determinant equals that of its weight-graded object. The nonsplit
Kummer extension from the unit $2$ over $\mathbb Z[1/2]$ and its split
Tate object therefore have the same incomplete $L$-function. The omission
of ramified factors is essential: inertia invariants need not preserve a
short exact sequence. These findings refute recovery of extensions from
semisimplified good Euler data; they do not assert that all richer
arithmetic data agree.

AG-EXT-K records the independent rational mixed-Tate background from Deligne--Goncharov,
\emph{Groupes fondamentaux motiviques de Tate mixte},
arXiv:math/0302267v2 (27 February 2003), (2.1.3), printed p.~14
(PDF p.~16):
\[
 \operatorname{Ext}^1(\mathbb Q(0),\mathbb Q(n))=
 \begin{cases}
 0&n\le0,\\
 \mathcal O_S^\times\otimes\mathbb Q&n=1,\\
 K_{2n-1}(k)\otimes\mathbb Q&n\ge2.
 \end{cases}
\]
Propositions 2.2--2.3 on the same page identify the corresponding graded
free motivic Lie algebra, while Section 2.4 explicitly distinguishes a
free presentation from a canonical set of generators. These statements
are unconditional background. They do not calculate integral torsion or
the pro-$\ell$ kernels in the prompt's cote 149 programme.

\subsection{An indexed integral $K$-group statement and its exact Tate-twist consequences}
\label{ag:k8}

\paragraph{AG-K8: the inherited theorem now has a primary proof locator.}
The assertion $K_8(\mathbb Z)=0$ in RETURN\_ALPOGE\_BRIDGES,
R-KGROUP/F-SM14--16, is a separate unconditional theorem.
It is Theorem~6.4 of Dutour Sikiri\'c--Elbaz-Vincent--Kupers--Martinet,
\emph{Voronoi complexes in higher dimensions, cohomology of
$GL_N(\mathbb Z)$ for $N\ge8$ and the triviality of $K_8(\mathbb Z)$},
JIMJ~25 (2026), 565--590; the primary author version
arXiv:1910.11598v1 has Theorem~6.4 and its proof on p.~19.
The result and its short proof were already available in 2019;
2026 is the journal issue year, not a consequence of the adopted A.

The exact dependency chain also explains the integral input.
With $R=\mathbb Z[1/7!]$, their Proposition~6.1 (author p.~18)
gives $H_9(BQ;R)=0$, and Proposition~6.3 (p.~19) makes the
Hurewicz map $\pi_9(BQ)\otimes R\to H_9(BQ;R)$ injective.
Since $\pi_9(BQ)=K_8(\mathbb Z)$, this proves vanishing after
inverting $2,3,5,7$.  Weibel VI, Corollary~9.8 (chapter p.~55)
removes the 2-primary part; Proposition~10.5 (p.~63) removes
the $3,5,7$ parts because those primes are regular.
These steps prove the stated integral vanishing from the cited
homological computations and standard integral comparisons.
Neither a rational special value nor a zero-location assertion
supplies the missing integral information by itself.

Kupers's independent \emph{A short proof that $K_8(\mathbb Z)=0$},
dated 21 October 2019, Theorem~A and pp.~1--4, supplies a
second short route.  Its rank-filtration spectral sequence has
zero entries on total degree nine after inverting $7!$;
the rest of the proof uses the low $K$-groups and the sphere
spectrum to pass from homology to homotopy.  The complete
four-page proof was read, as were the complete final reductions
in the joint paper.  The underlying earlier Voronoi enumeration,
Church--Putman proof and Quillen--Lichtenbaum theorem remain
cited classical inputs; no fresh audit of their computations
or full proofs is represented here.  In the usual homological
Serre spectral-sequence convention, the short proof's displayed
transgression $(9,0)\to(0,8)$ is $d_9$, not its printed $d_8$;
this indexing correction does not affect the argument.

\begin{proposition}[AG-K8-TWIST: exact integral and finite-twist vanishing]
For every odd prime $p$, let $U_p=\operatorname{Spec}\mathbb Z[1/p]$,
let $C_p$ be the $p$-primary part of
$\operatorname{Cl}(\mathbb Q(\zeta_p))$, and let
$\omega:\operatorname{Gal}(\mathbb Q(\zeta_p)/\mathbb Q)
\to\mathbb Z_p^\times$ be the Teichmüller character.
Then, unconditionally,
\begin{align*}
 H^2_{\text{ét}}(U_p,\mathbb Z_p(5))&=0,
 &H^2_{\text{ét}}(U_p,\mu_p^{\otimes5})&=0,\\
 C_p^{(-4)}&=0,
 &K_{8+2a(p-1)}(\mathbb Z)\{p\}&=0\qquad(a\ge0).
\end{align*}
Here the integral cohomology is the standard continuous
$p$-adic group in the cited comparison, and braces denote
the entire $p$-primary subgroup.  The class-group assertion
is the component $C_p^{(p-5)}$ in the joint paper's
Corollary~7.1, author p.~20, for $p>3$; the same statement
at $p=3$ follows from the proof below.  These are precisely
indexed classical consequences, not new applications of A or M.
\end{proposition}
\begin{proof}
Weibel VI, Theorem~8.2 and Corollary~8.3 (chapter pp.~47--48)
give, for odd $p$ and $i>0$,
\begin{align*}
 K_{2i}(\mathbb Z)\{p\}&\simeq
 H^2_{\text{ét}}(U_p,\mathbb Z_p(i+1)),\\
 K_{2i}(\mathbb Z)/p&\simeq
 H^2_{\text{ét}}(U_p,\mu_p^{\otimes(i+1)}).
\end{align*}
Putting $i=4$ and using AG-K8 proves the first two assertions.
We give the class-group comparison explicitly to fix its
character and the passage from mod-$p$ to integral vanishing.
Put $R_p=\mathbb Z[\zeta_p,1/p]$ and
$G=\operatorname{Gal}(\mathbb Q(\zeta_p)/\mathbb Q)$.
The finite etale cover $\operatorname{Spec}R_p\to U_p$
has degree $p-1$, so its Hochschild--Serre sequence,
with coefficients of exponent $p$, reduces to taking
$G$-invariants.  There is one prime of $\mathbb Q(\zeta_p)$
over $p$, and it is principal.  Consequently
$\operatorname{Pic}(R_p)=\operatorname{Cl}(\mathbb Q(\zeta_p))$.
The field is totally imaginary and only that one finite
place is removed.  The global Brauer invariant sequence
therefore gives $\operatorname{Br}(R_p)=0$; this is exactly
Weibel (8.1.1) and Example~8.3.2.  The Kummer sequence now
identifies $H^2_{\text{ét}}(R_p,\mu_p)$ with
$\operatorname{Pic}(R_p)/p$.
Tensoring by the rank-one module $\mu_p^{\otimes4}$,
and then taking invariants, yields
\[
 H^2_{\text{ét}}(U_p,\mu_p^{\otimes5})
 \simeq\bigl((C_p/p)\otimes\mu_p^{\otimes4}\bigr)^G
 \simeq (C_p/p)^{(-4)}.
\]
Thus the character is $\omega^{-4}$, not $\omega^{-5}$.
The order $p-1$ of $G$ is a $p$-adic unit, so the
Teichmüller idempotents split $C_p$ integrally and commute
with reduction modulo $p$.  Hence
$C_p^{(-4)}/p=0$.  A finite abelian $p$-group equal to its
$p$-multiple is zero, proving the full-component vanishing.

Finally set $i=4+a(p-1)$.  The mod-$p$ cyclotomic
character has order dividing $p-1$, whence
$\mu_p^{\otimes(i+1)}\simeq\mu_p^{\otimes5}$.
The same comparison gives
$K_{8+2a(p-1)}(\mathbb Z)/p=0$.
Positive even $K$-groups of $\mathbb Z$ are finite;
thus their $p$-primary parts vanish.  This proves the
entire displayed family.  It is also the $k=2$
instance of the periodic criterion in Weibel VI,
Theorem~10.9 (p.~65), with the regular-prime case
provided by Proposition~10.5.
\end{proof}

Geometrically, this kills one precisely specified
Tate-twisted $H^2$ group on each $U_p$ and the matching
cyclotomic unramified abelian class-group component.
It does not kill all positive-weight relation modules,
choose normalized Soule generators, or identify the
full pro-$p$ kernel in the LM6/Ihara programme.
The attached Alpoge footnote and step~17 still do not
export an indexed theorem on their own; the theorem
and comparison just supplied are separate proved inputs.

\subsection{AG-KUMMER: an integral obstruction on the actual cote 149 sheets}
\label{ag:kummer}

The original Montpellier PDF was inspected here, not merely the intake.
It has a cover page: PDF p.~80 has repository footer 79, and so on.
At PDF p.~80/footer 79, equations (13)--(17) relate an inertia-normalizer
extension, its $H^2$ class, and Kummer cohomology. PDF p.~81/footer 80,
(18)--(19), identifies the relevant line bundle as the normal bundle of
the section. PDF p.~82/footer 81, (20)--(21), and PDF p.~85/footer 84,
(22), describe the $H^1$ and completed-unit ambiguity. These are the
actual geometric objects on the cited sheets. Their link to LM6
A.44--A.45 is supplied by ORDER item 31; no theorem about an irregular
prime or $K_8(\mathbb Z)$ is printed on these pages.

Here is a precise integral consequence, independent of A and M. Let $S$
be a scheme on which $\ell$ is invertible, put $P=\operatorname{Pic}(S)$,
and let $\operatorname{Br}'(S)=H^2_{\text{ét}}(S,\mathbb G_m)$.
Write $\widehat P_\ell=\varprojlim_nP/\ell^nP$ and
$T_\ell A=\varprojlim_n A[\ell^n]$, with transition maps multiplication
by $\ell$. There are exact sequences
\begin{align*}
 0&\longrightarrow\widehat P_\ell
   \longrightarrow\varprojlim_nH^2_{\text{ét}}(S,\mu_{\ell^n})
   \longrightarrow T_\ell\operatorname{Br}'(S)\longrightarrow0,\\
 0&\longrightarrow\Gamma(S,\mathcal O_S^*)^{\wedge}_\ell
   \longrightarrow\varprojlim_nH^1_{\text{ét}}(S,\mu_{\ell^n})
   \longrightarrow T_\ell P\longrightarrow0.
\end{align*}
We deliberately use inverse-limit cohomology: identifying it with a
chosen definition of continuous cohomology requires its own
$\varprojlim{}^1$ check. If $P$ is finitely generated, the second right
term is zero. A line bundle $\mathcal L$ has zero compatible integral
$\ell$-adic Chern class if and only if its class in $P$ is torsion of
order prime to $\ell$.

\begin{proof}
The Kummer exact sequence gives
\[
 0\longrightarrow P/\ell^nP
 \longrightarrow H^2(S,\mu_{\ell^n})
 \longrightarrow\operatorname{Br}'(S)[\ell^n]\longrightarrow0,
\]
and the analogous degree-one sequence with units modulo $\ell^n$th
powers and $P[\ell^n]$. The exact source is Stacks Project,
Lemma 59.28.1, tag 03PK, and its immediately following cohomology
sequence. The transition maps on the left terms are surjective.
Consequently a compatible sequence on the right lifts compatibly in
the middle: after choosing a lift at level $n$, correct any lift at
level $n+1$ by a lift of its discrepancy in the left term. This recursive
construction proves right-surjectivity after taking inverse limits;
left-exactness is immediate. It proves both displayed sequences.

Now write the finitely generated group $P$ as
$\mathbb Z^a\oplus P[\ell^\infty]\oplus P_{\ell'}$, with finite torsion
summands. Its $\ell$-adic completion is
$\mathbb Z_\ell^a\oplus P[\ell^\infty]$, and the kernel of completion
is exactly $P_{\ell'}$. Also no nonzero compatible infinitely
$\ell$-divisible torsion sequence can occur in a finitely generated
group, so $T_\ell P=0$. The injection in degree two then proves the
Chern-class criterion.
\end{proof}

\paragraph{AG-KUMMER-2: a fully explicit integral/rational separation.}
Let $K=\mathbb Q(\sqrt{-35})$, $\omega=(1+\sqrt{-35})/2$, and
$S=\operatorname{Spec}\mathbb Z[\omega,1/2]$. The invertible ideal
$\mathfrak p=(3,\omega)$ determines a line bundle $\mathcal L$ whose
compatible integral $2$-adic Chern class is nonzero, although its image
after tensoring with $\mathbb Q_2$ is zero.

\begin{proof}
The ring of integers is $\mathbb Z[\omega]$ and its discriminant is
$-35$. The minimal polynomial $X^2-X+9$ is irreducible modulo $2$,
so $(2)$ is the only prime over $2$, and is principal. Modulo $3$ it is
$X(X-1)$; thus the primes over $3$ are
$\mathfrak p=(3,\omega)$ and $\overline{\mathfrak p}=(3,\omega-1)$.
Since $N(\omega)=9$ and $\omega$ is a unit modulo
$\overline{\mathfrak p}$, unique ideal factorization gives
$\mathfrak p^2=(\omega)$. The ideal $\mathfrak p$ is not principal:
an integral generator would have norm $3$, whereas
\[
 N(a+b\omega)=a^2+ab+9b^2=(a+b/2)^2+35b^2/4
\]
cannot be $3$ for integers $a,b$.

Minkowski's bound is $2\sqrt{35}/\pi<4$, so every ideal class has a
representative of norm $1$, $2$, or $3$. There is no ideal of norm $2$;
the two of norm $3$ have the same class of order two. Hence the class
group is $\mathbb Z/2$. The precise classical bound is Milne,
\emph{Algebraic Number Theory}, v3.08 (19 July 2020), Theorem 4.3,
PDF p.~72. Inverting $2$ quotients the class group by the classes of
the primes over $2$, which here are principal, so
$\operatorname{Pic}(S)=\mathbb Z/2$ as well. To see this last assertion
directly, fractional ideals of the localization are obtained by forgetting
the valuation at $(2)$; a relation introduced in the localization differs
from a principal ideal only by a power of that principal prime.

AG-KUMMER now injects this nonzero class into the inverse limit of
$H^2(S,\mu_{2^n})$. It has order two there, so rationalization kills it.
The total space of $\mathcal L$ has a zero section with normal bundle
$\mathcal L$, making the example a literal normal-bundle obstruction of
the kind appearing in cote 149. No claim identifying this particular
scheme's group-cohomology extension class is needed.
\end{proof}

This proves a concrete limit of rational information: a rationally zero
normal Chern class need not be integrally zero. A supplies Dirichlet GRH
for the quadratic field here, but changes none of this calculation.
The finite-level Kummer test is therefore a genuine extra gate for an
integral inertia-splitting argument. It does not refute the original
LM6 statements, whose comparison maps must still be supplied.

\subsection{AG-C34-TYPE and AG-C34-REPAIR: resolving the commutator notation}
\label{ag:c34repair}

The official cote 34 PDF has a cover: PDF p.~53 is manuscript/footer
p.~52, and PDF p.~64 is p.~63. Theorem 1.1 on the former page fixes a
field $k$ finitely generated over its prime field, a normal geometrically
connected $k$-scheme $X$ locally of finite type, a prime
$\ell\ne\operatorname{char}k$, and an open subgroup
$\pi\subseteq\pi_1(X)$. It defines $\pi_0$ to be the image in
$G_k$ and $I=\pi\cap\pi_1(X_{\bar k})$, in
\[
 1\longrightarrow I\longrightarrow\pi\longrightarrow\pi_0
 \longrightarrow1.
\]
Its $I(\ell)$ is the maximal abelian pro-$\ell$ quotient of $I$.
The theorem states that $I(\ell)$ is finitely generated over
$\mathbb Z_\ell$ and that its $\pi_0$-coinvariants are finite.
The proof was located and read in Sections 2--3, PDF pp.~70--78
(manuscript pp.~69--77), rather than guessed from the intervening
applications. Section 2.4, PDF p.~72, makes the dual $H^1$-invariant
reduction; Section 2.5, PDF pp.~73--74, uses a curve compactification,
its cyclotomic boundary, and the finite-field curve theorem after
spreading out. Sections 3.1--3.3, PDF pp.~75--78, reduce normal $X$ to
a generic hyperplane curve using surjectivity on fundamental groups.
These are classical geometric inputs; A and M do not enter them.

In the crossed-out Remark 1.17 on PDF p.~64, $H$ is a quotient of $I$
whose kernel is $\pi$-stable, and the readable final hypothesis is that
$H$ is a solvable pro-$\ell$ group. The preceding proposed
$\ell$-adic Lie-group language is itself struck out. The remark writes
\[
 1\longrightarrow H\longrightarrow\pi'\longrightarrow\pi_0
 \longrightarrow1
\]
and places the closure of $[\pi_0,H]$ as an open subgroup of $\pi'$.
This is a tentative, struck manuscript assertion, not a theorem.

\paragraph{AG-C34-TYPE (the literal ambient-group obstruction).}
For an infinite $\pi_0$, every commutator of a lift of an element of
$\pi_0$ with an element of $H$ lies in $H$, and its closure cannot be
open in $\pi'$. Indeed the quotient by $H$ is the infinite group
$\pi_0$, whereas a subgroup open in a profinite group has finite index.
For a concrete example satisfying the manuscript's geometric hypotheses,
take $X=\mathbb G_m/\mathbb F_q$ and $\ell\nmid q$. Its geometric
pro-$\ell$ quotient $H=\mathbb Z_\ell(1)$ is abelian and hence
solvable; the rational point $1$ gives a section. Arithmetic Frobenius
acts as multiplication by $q$, so the action commutators generate
$(q-1)\mathbb Z_\ell(1)$. This is open in $H$ but is not open in
$\pi'$, whose quotient is $G_{\mathbb F_q}\simeq\widehat{\mathbb Z}$.
Thus the literal ambient group in the inherited Q03 cannot be retained.
For nonabelian $H$ there is a second issue: $\pi_0$ acts canonically
only by outer automorphisms unless a section or lifts are specified.

\paragraph{AG-C34-REPAIR (the intrinsic corrected statement).}
Let $H$ be \emph{any} pro-$\ell$ quotient of $I$ with $\pi$-stable
kernel, and let $P$ be the resulting quotient extension of $\pi_0$ by
$H$. Define $K$ to be the closed subgroup generated by all
$[p,h]$, $p\in P$, $h\in H$. Then $K$ is open in $H$, and
\[
 H/K\simeq (H^{\mathrm{ab}})_{\pi_0},
 \qquad [H:K]\le |I(\ell)_{\pi_0}|<\infty.       \tag{AG-C34.1}
\]
In particular the solvability and Lie-group hypotheses are unnecessary
for this precise corrected conclusion.

\begin{proof}
Because $H$ is normal in $P$, conjugating a commutator $[p,h]$ by
$a\in P$ gives $[apa^{-1},aha^{-1}]$. Hence $K$ is normal in $P$.
It contains the closed commutator subgroup of $H$, so $H/K$ is abelian;
all conjugation by $P$ is trivial on it. Conjugation on
$H^{\mathrm{ab}}$ factors canonically through $\pi_0$, irrespective
of any section. Killing its action differences is exactly the additional
quotient by the images of the commutators $[p,h]$. This proves the
isomorphism in (AG-C34.1).

The surjection $I\twoheadrightarrow H$ induces an equivariant
surjection $I(\ell)\twoheadrightarrow H^{\mathrm{ab}}$ by the
universal property of the maximal abelian pro-$\ell$ quotient.
Taking coinvariants remains surjective: it is simply quotienting both
compact abelian groups by the closed subgroups generated by their
action differences. Theorem 1.1 therefore supplies a surjection from
the finite group $I(\ell)_{\pi_0}$ to $H/K$. This proves the bound.
Since $K$ is closed, finite index in the compact group $H$ implies
openness (the finite Hausdorff quotient is discrete).
\end{proof}

This repairs the archived Q03 as a theorem already implied by its
stated abelian predecessor. It does not identify the unqualified
symbol $[\pi_0,H]$ in every possible intended meaning. A smaller
commutator subgroup formed from a chosen section, or openness in a
different ambient group, needs a separately stated question. In
particular the repair is not described as an automatic proof of all
interpretations of the crossed-out remark.

\paragraph{AG-C34-INDEX (an explicit Frobenius refinement).}
More generally, let $M=H^{\mathrm{ab}}$ be a finitely generated
$\mathbb Z_\ell$-module with continuous $\Gamma$-action, let $T$ be
its finite torsion subgroup, and let $\Lambda=M/T$. Suppose some
$F\in\Gamma$ has no eigenvalue $1$ on
$\Lambda\otimes\mathbb Q_\ell$. Then
\[
 |\operatorname{coker}(F-1\mid M)|
   =|T^F|\ell^{v_\ell\det(1-F\mid\Lambda\otimes\mathbb Q_\ell)},
 \qquad
 |M_\Gamma|\le
 |T^F|\ell^{v_\ell\det(1-F\mid\Lambda\otimes\mathbb Q_\ell)}.
                                                               \tag{AG-C34.2}
\]
Equality in the second assertion holds if $\Gamma$ is topologically
generated by $F$.

\begin{proof}
The assumption gives $\Lambda^F=0$. Apply the snake lemma to
$0\to T\to M\to\Lambda\to0$ with vertical maps $F-1$.
It yields
\[
 0\longrightarrow T/(F-1)T\longrightarrow M/(F-1)M
 \longrightarrow\Lambda/(F-1)\Lambda\longrightarrow0.
\]
For an endomorphism of the finite group $T$, kernel and cokernel have
the same cardinality. Smith normal form over $\mathbb Z_\ell$ gives
the order of the last group as the displayed power of $\ell$.
This proves the equality. The continuous image $(F-1)M$ is compact
and hence closed; the full coinvariant quotient is a quotient of it.
If the powers of $F$ are dense in $\Gamma$, every action difference
is a limit of differences of powers, all of which belong to
$(F-1)M$, so the two quotients agree.
\end{proof}

For a smooth proper geometrically connected curve $C/\mathbb F_q$,
let $H$ be its full geometric pro-$\ell$ fundamental group and $J$ its
Jacobian. Then $H^{\mathrm{ab}}=T_\ell J$, compatibly with Galois
action, and the corrected commutator quotient is
\[
 H/\overline{[P,H]}
       \simeq T_\ell J/(F-1)T_\ell J
       \simeq J(\mathbb F_q)[\ell^\infty].             \tag{AG-C34.3}
\]
For the fundamental-group comparison, an exact modern primary locator
is J.~S.~Milne, \emph{Jacobian Varieties}, Proposition 9.1,
pp.~30--32, together with \emph{Abelian Varieties}, Remarks 15.3
and 15.5, pp.~28--29. To check the last isomorphism directly, use
$0\to T_\ell J\to V_\ell J\to J(\overline{\mathbb F}_q)
[\ell^\infty]\to0$. The fixed subgroup of $T_\ell J$ is zero:
$J(\mathbb F_q)$ is finite, so a compatible indefinitely divisible
sequence of its $\ell$-power torsion points vanishes. Thus $F-1$ is
invertible on $V_\ell J$, and the snake lemma gives (AG-C34.3).
Consequently the index is exactly $\ell^{v_\ell\#J(\mathbb F_q)}$.
This finite-field calculation even uses less than the full Weil bound.
In higher-dimensional smooth proper situations purity supplies the
corresponding no-eigenvalue-$1$ test, but finite-field purity is already
classical and no new global GRH input is involved.

\subsection{AG-C34-SECTION: the solvable quotient with a chosen section}
\label{ag:c34section}

There is a second precise interpretation of the struck Remark 1.17 that
can now be settled. Keep the exact geometric hypotheses and notation of
AG-C34-REPAIR. Suppose the quotient extension
$1\to H\to P\to\pi_0\to1$ has a specified continuous section, so
$P=H\rtimes\pi_0$. Assume that $H$ is solvable of finite derived
length, using closed derived subgroups. This is stronger than merely
being an inverse limit of finite solvable groups; every pro-$\ell$
group already satisfies that latter condition. For the action $\alpha$
given by the section put
\[
 N=\overline{\langle h^{-1}\alpha_g(h):h\in H,
                                      \ g\in\pi_0\rangle}.
\]
\textbf{Statement.} The subgroup $N$ is normal in $P$ and open in
$H$. Thus the chosen-section action commutators have finite index in
the geometric group. They need not equal the intrinsic subgroup
$\overline{[P,H]}$ of AG-C34-REPAIR.

\begin{proof}
For a fixed automorphism $\alpha$, write
$c_\alpha(h)=h^{-1}\alpha(h)$. The identity
\[
 c_\alpha(hx)=x^{-1}c_\alpha(h)x\,c_\alpha(x)
\]
shows that the subgroup generated by all the $c_\alpha(h)$ is normal
in $H$. Applying $\alpha_g$ to $c_{\alpha_f}(h)$ gives
$c_{\alpha_{gfg^{-1}}}(\alpha_g(h))$, so it is also stable under
$\pi_0$. Its closure $N$ is consequently normal in $P$. Set
$Q=H/N$. The induced action of $\pi_0$ on $Q$ is trivial by
construction, whence $P/N=Q\times\pi_0$.

We claim that every open subgroup $U$ of $Q$ has finite abelianization.
Take the inverse image of $U\times\pi_0$ first in $P$ and then in
the original open arithmetic subgroup $\pi$. This gives an open
subgroup $\pi_U$ of $\pi_1(X)$ with the same arithmetic image
$\pi_0$. Its geometric kernel $I_U$ maps onto $U$. Conjugation on
$U^{\mathrm{ab}}$ is trivial: the $U$-part acts by inner
automorphisms and the $\pi_0$-part acts trivially. Therefore
\[
 I_U(\ell)_{\pi_0}\twoheadrightarrow U^{\mathrm{ab}}.
\]
The left group is finite by Theorem 1.1 applied to the open subgroup
$\pi_U$, proving the claim. This use of the theorem on \emph{all
the relevant finite covers}, rather than on $H$ alone, is essential.

Let $Q^{(0)}=Q$ and
$Q^{(i+1)}=\overline{[Q^{(i)},Q^{(i)}]}$. Starting with the open
subgroup $Q$, the claim implies inductively that $Q^{(i+1)}$ has
finite index in the open subgroup $Q^{(i)}$, and hence is open in
$Q$. Since $Q$ is a quotient of the solvable group $H$, this finite
derived series ends at $1$. Thus $1$ is open in $Q$, so the compact
group $Q$ is finite. Equivalently $N$ is open in $H$.

For the non-equality assertion, take a finite nonabelian $\ell$-group
$H$ and let $\pi_0$ act trivially. Then $N=1$ whereas
$\overline{[P,H]}=[H,H]\ne1$. Both are open in $H$. This last
example concerns the subgroup identity, and asserts no realization of
that abstract example by a particular arithmetic variety.
\end{proof}

This is additional movement on BRIDGES Q03: under the explicitly stated
section hypothesis, its solvability condition is sufficient for the
action-commutator conclusion in the corrected ambient group $H$.
No section is inferred from purity or from A/M. Without a section the
intrinsic assertion AG-C34-REPAIR remains the canonical formulation.

\subsection{AG-C34-NONSOLV: why compact analytic groups alone do not suffice}
\label{ag:c34nonsolv}

The finite-derived-length condition just used cannot be replaced, at
the level of this group argument, by finite generation or compact
$\ell$-adic analyticity. The following counterexample is deliberately
abstract; no geometric realization is asserted.

Fix an odd prime $\ell$ and put
\[
 H_n=\ker\bigl(\mathrm{SL}_2(\mathbb Z_\ell)
            \longrightarrow\mathrm{SL}_2(\mathbb Z/\ell^n)\bigr),
 \qquad H=H_1.
\]
Then $H$ is an infinite, topologically finitely generated compact
pro-$\ell$ group, and every open subgroup of $H$ has finite
abelianization. Nevertheless, for the trivial action of any profinite
group $\Gamma$, the action-commutator subgroup is $1$, which is not
open in $H$. In particular the absence of eigenvalue $1$ on
$H^{\mathrm{ab}}\otimes\mathbb Q_\ell$ does not repair that
inference: here this vector space is zero.

\begin{proof}
For each $n\ge1$, the map
$1+\ell^n A\mapsto A\bmod\ell$ identifies
$H_n/H_{n+1}$ with the additive group
$\mathfrak{sl}_2(\mathbb F_\ell)$. The determinant equation gives
trace zero. Surjectivity follows from the upper and lower unipotent
matrices and the diagonal matrix
$\operatorname{diag}(1+\ell^n,(1+\ell^n)^{-1})$, whose leading terms
are the basis
\[
 E=\begin{pmatrix}0&1\\0&0\end{pmatrix},\quad
 F=\begin{pmatrix}0&0\\1&0\end{pmatrix},\quad
 T=\begin{pmatrix}1&0\\0&-1\end{pmatrix}.
\]
Thus $H$ is an inverse limit of finite $\ell$-groups and is compact;
the upper unipotent matrices with entry in $\ell\mathbb Z_\ell$
already show that it is infinite.
It is explicitly a three-dimensional $\ell$-adic analytic group:
write its matrices as
$\bigl(\begin{smallmatrix}a&b\\c&(1+bc)/a\end{smallmatrix}\bigr)$
with $a\in1+\ell\mathbb Z_\ell$ and
$b,c\in\ell\mathbb Z_\ell$. These three parameters are analytic
coordinates, and multiplication and inversion are analytic in them.

The three displayed lifts at level $1$ topologically generate $H$.
Indeed, for $n\ge1$ and odd $\ell$,
\[
 (1+\ell^n A)^\ell\equiv1+\ell^{n+1}A
                       \pmod{\ell^{n+2}}.
\]
Their successive $\ell$-power iterates therefore give all three basis
vectors in every $H_n/H_{n+1}$. Any element of $H$ can be matched
successively modulo $H_2,H_3,\ldots$ by products of elements in the
closed group they generate. The products converge to the given element,
so that closed group is all of $H$.

For $n,m\ge1$, direct matrix multiplication gives
\[
 [1+\ell^n A,1+\ell^m B]
       \equiv1+\ell^{n+m}(AB-BA)
                      \pmod{\ell^{n+m+1}}.
\]
The Lie brackets $[E,F]=T$, $[T,E]=2E$, and $[T,F]=-2F$
span $\mathfrak{sl}_2(\mathbb F_\ell)$ because $\ell$ is odd.
Let $C_n=\overline{[H_n,H_n]}$. For every $k\ge2n$, commutators
of $H_n$ with $H_{k-n}\subseteq H_n$ therefore show that
$C_n\cap H_k$ surjects onto $H_k/H_{k+1}$. Successively lifting a
given element of $H_{2n}$ through these quotients and taking the
convergent product proves
\[
 H_{2n}\subseteq C_n.
\]
Every open subgroup $U\subseteq H$ contains some $H_n$. Its closed
commutator subgroup contains $C_n$, and hence $H_{2n}$, so
$U^{\mathrm{ab}}$ is finite. Taking $U=H$ also proves
$H^{\mathrm{ab}}\otimes\mathbb Q_\ell=0$.
For a trivial $\Gamma$-action, each action commutator is $1$,
which has infinite index in $H$. This proves the claimed failure.
\end{proof}

The example pinpoints what AG-C34-SECTION uses beyond abelian
coinvariant control: its induction must terminate. It does not challenge
the intrinsic commutator theorem, since $\overline{[H,H]}$ is open
here, nor any additional geometric constraints that could exclude this
trivial-action example in a proposed realization.

\subsection{Cote 149: the fundamental-group class and the section obstruction}
\label{ag:c149-comparison}

The comparison omitted by a bare Kummer calculation is available under
precise homotopy hypotheses. Borne--Emsalem--Stix,
\emph{Lifting Galois sections along torsors}, Math. Nachr.~288 (2015),
1939--1951, arXiv:1301.4429v2, Proposition 3.1 and Proposition 3.4
(PDF pp.~5--6), and Proposition 4.4 (PDF pp.~8--9), supply the
following classical statement. The complete proofs and the continuous
cohomology argument in Proposition 4.2 (PDF p.~7) have been read.

\paragraph{AG-C149-COMP (the exact comparison).}
Let $S$ be a geometrically connected, geometrically unibranch variety
over a characteristic-zero field $k$, with
$\pi^{\text{ét}}_2(S,\bar s)=0$. Let $\mathcal L$ be a line
bundle and $\mathcal L^\times$ its complement of the zero section.
For compatible base points there is an exact sequence
\[
 1\longrightarrow\widehat{\mathbb Z}(1)
 \longrightarrow\pi_1(\mathcal L^\times)
 \longrightarrow\pi_1(S)\longrightarrow1.
\]
The conjugation action on its kernel factors through $G_k$ and is
cyclotomic. Quotienting the kernel by its prime-to-$\ell$ factors gives
an extension
\begin{equation}\label{eq:ag-c149-extension}
 1\longrightarrow\mathbb Z_\ell(1)
 \longrightarrow\mathcal E_\ell(\mathcal L)
 \longrightarrow\pi_1(S)\longrightarrow1.
\end{equation}
The comparison isomorphism
\[
 H^2_{\mathrm{cont}}(\pi_1(S),\mathbb Z_\ell(1))
 \simeq H^2_{\mathrm{cont}}(S,\mathbb Z_\ell(1))
\]
sends its extension class $e_\ell(\mathcal L)$ to the arithmetic
Kummer class $c_{1,\ell}(\mathcal L)$. Continuous cohomology here is
the cohomology of the inverse system $(\mu_{\ell^n})$; it is not
silently identified with an inverse limit in arbitrary situations.

The source proves exactness from the homotopy sequence of the torus
torsor and proves the class comparison by compatible transgressions:
the identity of the Tate module maps to the negative of each of the
two classes. The hypothesis on $\pi_2$ is substantive: over an
algebraically closed field, the torsor
$\mathbb A^{d+1}\setminus\{0\}\to\mathbb P^d$, $d\ge1$, has trivial
total-space fundamental group, so its fibre inertia does not inject
(source Remark 3.2, PDF p.~5). Smooth geometrically connected curves
other than the geometric projective line are algebraic $K(\pi,1)$
spaces (source Section 3.1, PDF p.~4).

Geometrically, $\mathcal L$ is the normal bundle of the zero section in
$\operatorname{Tot}(\mathcal L)$, so this is an actual punctured
normal-bundle extension. It instantiates the $K(\pi,1)$ comparison
explicitly contemplated on cote 149, original PDF p.~80, manuscript
p.~79, equations (13)--(17), followed by the normal-bundle discussion
on PDF p.~81, manuscript p.~80. For the torsion-line-bundle hyperbolic family below,
AG-C149-NORMAL supplies the corresponding inertia-normalizer map.
A different family still requires its own comparison; none of these
statements identifies the global tripod kernel in LM6 A.44--A.45.

\begin{proposition}[AG-C149-LIFT: torsion classes vanish on every section]
Under the hypotheses of AG-C149-COMP, a continuous section
$s:G_k\to\pi_1(S)$ of the arithmetic projection lifts to
$\mathcal E_\ell(\mathcal L)$ if and only if
$s^*e_\ell(\mathcal L)=0$. If $\mathcal L$ is torsion in
$\operatorname{Pic}(S)$, every such section lifts, although the
extension \eqref{eq:ag-c149-extension} need not split over
$\pi_1(S)$. When lifts exist, their classes modulo conjugation by
$\mathbb Z_\ell(1)$ form a torsor under
\[
 H^1_{\mathrm{cont}}(k,\mathbb Z_\ell(1))
   \simeq\varprojlim_n k^\times/(k^\times)^{\ell^n}.
\]
\end{proposition}
\begin{proof}
Pulling back \eqref{eq:ag-c149-extension} by $s$ gives an extension of
$G_k$ with class $s^*e_\ell(\mathcal L)$. A lift is precisely a
continuous splitting of this pullback, so the class criterion follows
from extension classification, equivalently source Lemma 4.5
(PDF p.~9) and its preceding comparison.

We verify the continuous-cohomology issue directly for the field.
Hilbert 90 and the Kummer sequence give
\[
 H^1(k,\mu_{\ell^n})=k^\times/(k^\times)^{\ell^n},\qquad
 H^2(k,\mu_{\ell^n})=\operatorname{Br}(k)[\ell^n].
\]
The maps on the first groups are the surjective quotient maps, and
those on the second groups are multiplication by $\ell$. Thus
$\varprojlim^1 H^1(k,\mu_{\ell^n})=0$, and the defining exact
sequence for continuous cohomology gives
\[
 H^2_{\mathrm{cont}}(k,\mathbb Z_\ell(1))
   \simeq T_\ell\operatorname{Br}(k).
\]
For any abelian group $A$, its Tate module
$T_\ell A=\varprojlim(A[\ell^n],\ell)$ is torsion-free over
$\mathbb Z_\ell$: if $\ell t=0$, then compatibility gives
$t_n=\ell t_{n+1}=0$ for every $n$, and induction excludes all
$\ell$-power torsion. Multiplication by an integer prime to $\ell$
is invertible. If $m[\mathcal L]=0$, additivity of the Kummer class
therefore gives
$m s^*e_\ell(\mathcal L)=0$, whence $s^*e_\ell(\mathcal L)=0$.
This is the $\mathbb G_m$, $\ell$-primary specialization of the
source's Proposition 4.2 and Lemma 4.6.

Finally, once one splitting is fixed, the difference of a second
splitting from it is a continuous $1$-cocycle. Conjugation by the
kernel changes that cocycle by a coboundary. This proves the torsor
claim. The finite groups $H^0(k,\mu_{\ell^n})$ form a
Mittag--Leffler inverse system; hence their $\varprojlim^1$ vanishes
and the degree-one continuous-cohomology sequence gives the displayed
Kummer identification.
\end{proof}

\begin{proposition}[AG-C149-ELLIPTIC: non-splitting does not obstruct all sections]
Let $E/\mathbb Q$ be the elliptic curve
\[
 y^2=x^3-x,
 \qquad O=\text{the point at infinity},\qquad P=(0,0),
\]
and let $\mathcal L=\mathcal O_E(P-O)$.
The class of the extension $\mathcal E_2(\mathcal L)\to\pi_1(E)$
is nonzero of order $2$, so this extension does not split. Its image
with $\mathbb Q_2(1)$ coefficients is zero. Nevertheless, every
Galois section $G_{\mathbb Q}\to\pi_1(E)$ lifts through this
extension. Thus vanishing of the obstruction after pullback by
every Galois section does not force the global integral class to
vanish.
\end{proposition}
\begin{proof}
The point $P$ is a nonidentity point of order $2$, since inversion on
this Weierstrass curve sends $(x,y)$ to $(x,-y)$. The standard
identification $E(\mathbb Q)=\operatorname{Pic}^0(E)$ sends $P$ to
$[\mathcal L]$, so $\mathcal L$ has order $2$.

Its discriminant is $64$, so the displayed integral model has good
reduction at $3$. Direct enumeration gives
\[
 E(\mathbb F_3)=\{O,(0,0),(1,0),(2,0)\}
                 \simeq(\mathbb Z/2\mathbb Z)^2.
\]
Properness and good reduction give a group homomorphism
$E(\mathbb Q)\to E(\mathbb F_3)$. If $P=2Q$ for a rational point
$Q$, reduction would yield $\bar P=2\bar Q=0$, whereas
$\bar P=(0,0)\ne O$. Therefore $P\notin2E(\mathbb Q)$.
If $\mathcal L=\mathcal M^{\otimes2}$ in $\operatorname{Pic}(E)$,
degrees give $\deg\mathcal M=0$, and the same Picard identification
would express $P$ as twice a rational point. Thus
$[\mathcal L]\ne0$ in $\operatorname{Pic}(E)/2$.

The Kummer injection
\[
 \operatorname{Pic}(E)/2\hookrightarrow H^2(E,\mu_2)
\]
shows that the reduction of $c_{1,2}(\mathcal L)$ modulo $2$ is
nonzero. Consequently the continuous integral class is nonzero;
its order is exactly $2$ by additivity and $\mathcal L^{\otimes2}
\simeq\mathcal O_E$. The curve has genus $1$ and is an algebraic
$K(\pi,1)$, so AG-C149-COMP identifies this with the nonzero
extension class. A torsion class maps to zero in cohomology with
$\mathbb Q_2(1)$ coefficients. Finally AG-C149-LIFT applies to
every Galois section, and such sections do exist because
$O\in E(\mathbb Q)$. This proves all the claims.
\end{proof}

This supplies the group/cohomology comparison missing from the purely
arithmetic example AG-KUMMER-2, under explicit hypotheses and for a
literal normal bundle. It also separates three questions on the cote:
nonvanishing of a global integral extension class, its rational image,
and the obstruction for a specified section. None of these distinctions
uses A or M, and none identifies an irregular-prime class-group
component or proves the remaining LM6 tripod-kernel assertion.

\subsection{An actual hyperbolic-family cusp normalizer}
\label{ag:c149-normalizer}

We now make the remaining normalizer map explicit for a family
constructed from a torsion line bundle. The only local normalizer
input is Stix, \emph{On cuspidal sections of algebraic fundamental
groups}, author manuscript dated 24 February 2012, Section 3,
Lemma 14 (PDF p.~7) and Proposition 16(1)--(2) (PDF p.~8).
For a characteristic-zero hyperbolic curve, two distinct prolonged
cusps have inertia groups with trivial intersection, and the
arithmetic normalizer of a cusp inertia group is its decomposition
group. The complete proofs of these statements have been read.
The product input is SGA~1, Expos\'e X, Corollary 1.7, recomposed
pp.~203--204 (arXiv PDF pp.~219--220), whose proof and Corollary 1.8
have also been read. Its proper-factor hypothesis will be retained.

\begin{proposition}[AG-C149-NORMAL: a cyclic-cover normalizer comparison]
Let $k$ have characteristic zero, let $S/k$ be smooth, proper and
geometrically connected, and fix a line bundle $\mathcal L$ with
an isomorphism $\tau:\mathcal L^{\otimes n}\simeq\mathcal O_S$,
where $n\ge2$. Set
\[
 B=\mathbb P^1_k\setminus\{0,1,\infty\},\qquad
 U=\{v\in\mathcal L^\times:\tau(v^{\otimes n})\ne1\}.
\]
Then $U\to S$ is a smooth family of hyperbolic curves of type
$(0,n+2)$. Compactify the line bundle by adjoining its infinity
section, so that the boundary consists of the zero section, the
infinity section, and the degree-$n$ finite \'etale divisor
$\tau(v^{\otimes n})=1$. Choose compatible geometric base points,
and let $I\simeq\widehat{\mathbb Z}(1)$ be the zero-cusp inertia
subgroup of a geometric fibre, viewed inside $\pi_1(U)$.
The open immersion $U\to\mathcal L^\times$ restricts to an
isomorphism
\begin{equation}\label{eq:ag-cusp-normalizer}
 N_{\pi_1(U)}(I)\ \xrightarrow{\ \sim\ }
 \pi_1(\mathcal L^\times)
\end{equation}
over $\pi_1(S)$, with the usual compatibility up to choices of
paths. In particular, it gives an exact sequence
\[
 1\longrightarrow I\longrightarrow N_{\pi_1(U)}(I)
 \longrightarrow\pi_1(S)\longrightarrow1,
\]
whose action on $I$ is cyclotomic. No $\pi_2(S)=0$ hypothesis is
needed for this explicit torsion-line-bundle construction.
\end{proposition}
\begin{proof}
The maps $v\mapsto\tau(v^{\otimes n})$ give finite \'etale
$\mu_n$-torsors
\[
 U\longrightarrow S\times B,\qquad
 \mathcal L^\times\longrightarrow S\times\mathbb G_m.
\]
This follows locally by writing the map as $z\mapsto u z^n$ with
$u$ a unit: its derivative is invertible when $z\ne0$.
After the finite \'etale cover
$T=\{e\in\mathcal L^\times:\tau(e^{\otimes n})=1\}\to S$,
the family is
$T\times(\mathbb P^1\setminus(\{0,\infty\}\cup\mu_n))$.
Thus it has the stated smooth hyperbolic fibres, and $U$ and
$\mathcal L^\times$ are geometrically connected.

Write $\Pi_X=\pi_1(X)$ and $G=G_k$.
Proper K\"unneth over $\bar k$, followed by the geometric/arithmetic
exact sequences, gives
\[
 \Pi_{S\times B}=\Pi_S\times_G\Pi_B,\qquad
 \Pi_{S\times\mathbb G_m}=\Pi_S\times_G\Pi_{\mathbb G_m}.
\]
Indeed the maps induced by projections are isomorphisms on their
geometric kernels and on the quotient $G$, and hence are
isomorphisms. This verifies the arithmetic product assertion even
when $S$ has no rational point.

Use additive notation for $\mu_n$. Choose a point in the geometric
fibre of $T$, so that $T$ is represented by a continuous crossed
cocycle $a:\Pi_S\to\mu_n$. Let
$b:\Pi_B\to\mu_n$ represent the Kummer torsor $z^n=t$, where
$t$ is the coordinate on $B$. The two actions on $\mu_n$ factor
through the same cyclotomic action of $G$. The first displayed
cover has cocycle $a+b$: locally write $v=e z$ with
$\tau(e^{\otimes n})=1$ and $z^n=t$. Consequently its fundamental
group is the stabilizer
\begin{equation}\label{eq:ag-cyclic-stabilizer}
 H=\{(\sigma,\beta)\in\Pi_S\times_G\Pi_B:
                 a(\sigma)+b(\beta)=0\}.
\end{equation}
The zero set of a crossed cocycle is a subgroup. Unless the
$\mu_n$ action is trivial, it need not be a normal subgroup; no
ordinary homomorphism-kernel assertion is being made here.
The same description gives a subgroup $H'$ for
$\Pi_{\mathcal L^\times}$, using the Kummer cocycle on
$\Pi_{\mathbb G_m}$.

Choose the zero cusp of $B$, its inertia $I_0$ and its
decomposition group $D_0$. In the product group the relative
inertia is $\{1\}\times I_0$. Since $t$ has ramification index
$n$ in the local parameter $v$ upstairs, the inertia in $H$ is
\[
 I=\{1\}\times nI_0.
\]
This also verifies its injection: multiplication by $n$ on the
free profinite procyclic group $I_0$ is injective. Equivalently,
the Kummer cocycle restricted to $I_0$ is its reduction modulo
$n$, so its intersection with $H$ is exactly $nI_0$.

If $\beta\in\Pi_B$ normalizes $nI_0$, then
$I_0\cap\beta I_0\beta^{-1}$ contains the nontrivial subgroup
$nI_0$. Stix's Proposition 16(1) forces the two prolonged cusps
to coincide, so $\beta\in D_0$. Conversely, $D_0$ normalizes
$I_0$ and its characteristic subgroup $nI_0$. Therefore
\[
 N_H(I)=H\cap(\Pi_S\times_G D_0).
\]
The filling map $B\to\mathbb G_m$ identifies
$D_0\xrightarrow{\sim}\Pi_{\mathbb G_m}$.
To see this without a choice of splitting, compare their exact
sequences with quotient $G$: on the quotient the map is the
identity, and on geometric inertia it is an isomorphism, since
both groups are $\widehat{\mathbb Z}(1)$ with the same local
parameter $t$. The injections of the punctured local group for
$B$ and for $\mathbb G_m$ are supplied by Stix's Lemma 14,
whose hypothesis is $\chi\le0$ and whose proof explicitly includes
$\mathbb G_m$. For the latter curve this also follows directly
from the Kummer covers $t\mapsto t^a$: geometrically, its entire
fundamental group is this inertia group.

The resulting isomorphism
$\Pi_S\times_GD_0\simeq\Pi_S\times_G\Pi_{\mathbb G_m}$
preserves the Kummer cocycles. It therefore identifies $N_H(I)$
with $H'$, proving \eqref{eq:ag-cusp-normalizer}.
Finally the rational point $1\in\mathbb G_m(k)$ gives
$\Pi_{\mathbb G_m}=\widehat{\mathbb Z}(1)\rtimes G$; its Kummer
cocycle is reduction of the first coordinate modulo $n$.
Thus the projection $H'\to\Pi_S$ is surjective, and its kernel
is $n\widehat{\mathbb Z}(1)\simeq\widehat{\mathbb Z}(1)$.
This proves exactness and the cyclotomic action, with the kernel
identified with $I$ by the same local parameter computation.
\end{proof}

\begin{proposition}[AG-C149-FULLLIFT: every arithmetic section lifts]
Under the hypotheses of AG-C149-NORMAL, every continuous section
$s:G_k\to\Pi_S$ lifts to $N_{\Pi_U}(I)$, and hence to $\Pi_U$.
This is a statement for the full profinite inertia group,
not just one $\ell$-primary quotient.
\end{proposition}
\begin{proof}
Using the last coordinates of the previous proof, the group in
question consists of pairs $(\sigma,z)$ with
\[
 \sigma\in\Pi_S,\quad z\in\widehat{\mathbb Z}(1),\qquad
 a(\sigma)+z\bmod n=0,
\]
and multiplication
$(\sigma,z)(\sigma',z')=(\sigma\sigma',z+\chi(\sigma)z')$,
where $\chi$ is the cyclotomic action. It suffices to lift the
crossed cocycle $-a\circ s:G_k\to\mu_n$ to a continuous
crossed cocycle with values in $\widehat{\mathbb Z}(1)$.

Kummer theory gives $H^1(k,\mu_n)=k^\times/(k^\times)^n$.
Choose $c\in k^\times$ representing $[-a\circ s]$ and a
compatible collection of roots $c^{1/m}\in\bar k$ for all
positive $m$, ordered by divisibility. Such a collection exists
because the finite nonempty root sets have surjective transition
maps. Then
\[
 \kappa(g)=\bigl(g(c^{1/m})/c^{1/m}\bigr)_m
       \in\varprojlim_m\mu_m=\widehat{\mathbb Z}(1)
\]
is a continuous crossed cocycle, each finite coordinate being
continuous. Its reduction modulo $n$ has the desired cohomology
class. If this reduction differs from $-a\circ s$ by the
coboundary $du$ of $u\in\mu_n$, lift $u$ to
$\tilde u\in\widehat{\mathbb Z}(1)$ and replace $\kappa$ by
$\kappa-d\tilde u$. Its reduction is then exactly $-a\circ s$.
The map $g\mapsto(s(g),\kappa(g))$ is the required continuous
group homomorphism and is a lift of $s$. Its projection to $G_k$
is the identity, so its composition with the inclusion in
$\Pi_U$ is an arithmetic section as well.
\end{proof}

\begin{proposition}[AG-C149-CUSP: a nonsplit normalizer with every section liftable]
Take $E/\mathbb Q$ and $\mathcal L=\mathcal O_E((0,0)-O)$
from AG-C149-ELLIPTIC, fix any rational trivialization of
$\mathcal L^{\otimes2}$, and form $U$ as above with $n=2$.
Then $U\to E$ is a family of four-punctured projective lines,
with two boundary sections and a degree-two \'etale boundary
divisor. Its zero-cusp normalizer extension over $\pi_1(E)$
does not split. Its $2$-primary class is nonzero of order $2$
and has zero rational image. Nevertheless every section
$G_{\mathbb Q}\to\pi_1(E)$ lifts through the full cusp normalizer.
\end{proposition}
\begin{proof}
The normal bundle of the zero section in the ruled compactification
is $\mathcal L$, because its neighbourhood is the total space of
$\mathcal L$. AG-C149-NORMAL identifies the actual cusp
normalizer with $\pi_1(\mathcal L^\times)$, over $\pi_1(E)$
and on inertia. The curve $E$ is an algebraic $K(\pi,1)$, so
AG-C149-COMP identifies the $2$-primary extension class with
$c_{1,2}(\mathcal L)$; AG-C149-ELLIPTIC proved that this class
is nonzero of order $2$. Thus its $2$-primary extension is
nonsplit, which also excludes a splitting of the full extension.
The same earlier computation gives its zero rational image.
AG-C149-FULLLIFT proves the assertion for every arithmetic
section, and the section associated with $O\in E(\mathbb Q)$
makes the assertion nonvacuous.
\end{proof}

Thus the normalizer comparison has been supplied for an explicit
hyperbolic family, rather than merely posited for a punctured line
bundle. This advances the actual integral-versus-rational and
section-lifting questions at cote 149, PDF pp.~80--82,
equations (13)--(21). The proof uses neither A nor M. Its scope is
the specified isotrivial family; identifying a different inertia
normalizer still needs its own map, and the global tripod-kernel
assertion LM6 A.45 remains separate.

\subsection{LM6 A.44--A.45: the exact integral generator gate}
\label{ag:lm6}

\paragraph{AG-LM6-REG (the known kernel-field theorem).}
For a prime $p$, let $\Omega_p^*$ be the fixed field of the kernel
of the outer Galois action on
$\pi_1^{\mathrm{pro}\text{-}p}
(\mathbb P^1_{\overline{\mathbb Q}}\setminus\{0,1,\infty\})$.
Let $\Omega_p$ be the maximal pro-$p$ extension of
$\mathbb Q(\mu_p)$ unramified away from $p$. Then
$\Omega_p^*\subseteq\Omega_p$, and equality is classical for odd
regular $p$. An exact primary statement is Hoshi,
\emph{On the Kernels of the Pro-$l$ Outer Galois Representations
Associated to Hyperbolic Curves over Number Fields}, revised March
2014 manuscript, Definition 1.8, Remark 1.8.1, and Theorem 1.9,
PDF p.~11. Its proof cites Sharifi's Theorem 1.1 and Brown's theorem.
This formulation corrects the informal phrase ``the kernel is the
Galois group of $\Omega_p$'': the kernel is
$\operatorname{Gal}(\overline{\mathbb Q}/\Omega_p^*)$.

Brown, \emph{Mixed Tate motives over $\mathbb Z$}, Ann. Math.~175
(2012), 949--976, Theorem 1.1 and the following paragraph, printed
p.~950, prove the motivic-generation assertion formulated as
Conjecture 1 on p.~949. Its Deligne--Ihara consequence concerns the
graded Lie algebra after tensoring with $\mathbb Q_p$. It does not
discard integral indices. Brown's theorem is a classical input here;
its entire proof has not been re-refereed in this run.

\paragraph{AG-LM6-IRREG (where the regular-prime argument uses more).}
Write $K=\mathbb Q(\mu_{p^\infty})$ and
$G=\operatorname{Gal}(\Omega_p^*/K)$, with its outer-action
filtration $F^mG$ and graded Lie algebra $\mathfrak g$.
Choose the usual maximal-image Soulé elements $\sigma_m\in F^mG$
for odd $m\ge3$. They define maps $\Psi:S\to G$ from the free
pro-$p$ group on these generators and
$\psi:\mathfrak s\to\mathfrak g$ from the free graded
$\mathbb Z_p$-Lie algebra. Sharifi,
\emph{Relationships between conjectures on the structure of pro-$p$
Galois groups unramified outside $p$}, Proc. Sympos. Pure Math.~70
(2002), 275--284, author PDF Theorem 1.1 p.~2 and proof pp.~8--9,
uses the freeness of
$\operatorname{Gal}(\Omega_p/\mathbb Q(\mu_p))$ at regular primes.
The full proof and its preceding group lemmas have been read.

For irregular $p$, his Theorem 1.3 (author PDF p.~3; proof p.~10),
together with Brown's rational freeness, says that $\Psi$ and
$\psi$ are not surjective if Greenberg's generalized conjecture
holds for $\mathbb Q(\mu_p)$. Here that conjecture says that the
unramified abelian pro-$p$ Iwasawa module over the compositum of
all $\mathbb Z_p$-extensions is pseudo-null. This is a conditional
obstruction to the proposed generator argument, not a disproof of
$\Omega_p^*=\Omega_p$: the latter could hold with additional
integral generators and relations.

\begin{proposition}[AG-LM6-NORM: the filtered Soulé normalization]
For odd $p$ and odd $m\ge3$, in the integral normalization of
Sharifi's equation (2.10), author PDF p.~7,
\[
 \kappa_m(F^mG)\subseteq
 p^{\sum_{j\ge1}\lfloor(m-1)/p^j\rfloor}\mathbb Z_p(m).
\]
In particular $\kappa_5(F^5G)\subseteq3\mathbb Z_3(5)$,
although $3$ is regular. Thus surjectivity of these filtered
characters onto the full displayed lattice is false in every
degree with $m-1\ge p$.
\end{proposition}
\begin{proof}
The cited identity is
\[
 \kappa_m=(p^{m-1}-1)(m-1)!\lambda_m^{(p)}
       \quad\hbox{on }F^mG,
\]
where $\lambda_m^{(p)}$ is integral, obtained from the integral
stable-derivation algebra. The first factor is a $p$-adic unit.
Counting the multiples of $p^j$ among $1,\ldots,m-1$ gives
$v_p((m-1)!)=\sum_{j\ge1}\lfloor(m-1)/p^j\rfloor$.
The inclusion follows, and for $(p,m)=(3,5)$ this valuation is $1$.
\end{proof}

This is a precise test for the suggestion in
\texttt{OURS\_FABLE\_PROGRAM\_SUGGESTIONS\_20261120\_EXPLAINED.md},
Section 5, linked there to LM6 A.44--A.45, bridge B3, and cote 149
scans 79--85. It does not assert nonsurjectivity on a different,
larger Galois domain, and it does not replace a normalized character
by an unnormalized one. A proposed extension of the regular-prime
proof must identify the exact integral generator quotient, its
relations, and the cohomological comparison producing them.
Neither endpoint A nor endpoint M supplies these data. The
$p=2$ and irregular-prime kernel-field questions in the intake
therefore remain unresolved by this argument; they are not listed
as refuted.

\subsection{The formal-height mechanism in Alpöge step 2}
\label{ag:formal-height}

The supplied Fable suggestion, Section 2, gives the height-dependent
derivative estimate but passes through a Lubin--Tate normal form.
The following proof works directly with the given formal group.
It strengthens the estimate to an explicit bound and requires no
identification of a general deformation with a particular
Lubin--Tate group.

\begin{proposition}[AG-FG-EXP: an exact exponential coefficient bound]
Let $R$ be a complete unramified discrete valuation ring of mixed
characteristic $(0,p)$, with fraction field $K$, perfect residue
field, and valuation $v(p)=1$, so $v(K^\times)=\mathbb Z$;
set $v(0)=+\infty$.
Let $F(X,Y)$ be a one-dimensional commutative formal group over
$R$. Suppose its reduction has finite height $h\ge1$, and put
$q=p^h$. Write
\[
 [p]_F(t)=pt+\sum_{j\ge2}a_jt^j.
\]
Thus $a_j\in pR$ for $2\le j<q$, $a_q\in R^\times$, and
all $a_j\in R$. Normalize the formal exponential by
\[
 e(z)=\exp_F(z)=z+\sum_{n\ge2}e_nz^n.
\]
Then
\[
 v(e_n)\ge-\frac{n-1}{q-1},\qquad
 v(e_n)\ge-\left\lfloor\frac{n-1}{q-1}\right\rfloor.
\]
The second inequality records that the valuation group is integral.
\end{proposition}
\begin{proof}
We include the logarithm construction to specify all normalizations.
Put $b(t)=\partial_2F(t,0)\in1+tR[[t]]$ and
\[
 \omega_F=\frac{dt}{b(t)},\qquad D=b(t)\frac{d}{dt}.
\]
The differential is integral, and $D$ preserves $R[[t]]$ with
$(Dt)(0)=1$. Associativity, differentiated in the third variable
at zero, gives
\[
 b(F(x,y))=\partial_2F(x,y)b(y).
\]
Define $\lambda(t)=\int_0^t b(u)^{-1}\,du\in K[[t]]$.
The displayed identity implies
$\partial_y\lambda(F(x,y))=\lambda'(y)$.
Evaluating the resulting constant at $y=0$ yields
$\lambda(F(x,y))=\lambda(x)+\lambda(y)$.
Its linear coefficient is $1$, so it has a formal inverse $e$,
with $e(0)=0$ and $e'(0)=1$. In particular
\[
 e(pz)=[p]_F(e(z)).
\]

We prove the first inequality by induction; it is equality for
$n=1$. Comparing the coefficient of $z^n$ in this identity gives
\[
 (p^n-p)e_n=
 \sum_{j=2}^n a_j
  \sum_{\substack{n_1+\cdots+n_j=n\\n_i\ge1}}
       e_{n_1}\cdots e_{n_j}.
\]
Every $n_i<n$, and the induction hypothesis bounds the valuation
of each product below by $-(n-j)/(q-1)$.
For $2\le j<q$, multiplication by $a_j$ adds at least $1$;
division by $p^n-p$, whose valuation is $1$, therefore leaves
a bound $-(n-j)/(q-1)\ge-(n-1)/(q-1)$.
For $j\ge q$, multiplication adds at least $0$, and division
leaves
\[
 -1-\frac{n-j}{q-1}\ge-\frac{n-1}{q-1}.
\]
The ultrametric inequality proves the induction. Since each
$e_n\in K$ has integral valuation, rounding up gives the second
bound. No convergence assertion or special normal form was used.
\end{proof}

\begin{proposition}[AG-FG-DER: integral derivatives with the exact error]
With $R,F,D,h,q$ as above, for every $f\in R[[t]]$ and integer
$a\ge1$,
\begin{equation}\label{eq:ag-height-derivative}
 v\bigl((D^af)(0)\bigr)\ge
 \max\left\{0,\ v(a!)-
                 \left\lfloor\frac{a-1}{q-1}\right\rfloor\right\}.
\end{equation}
In particular, writing $s_p(a)$ for the base-$p$ digit sum,
\[
 v\bigl((D^af)(0)\bigr)\ge
 a\left(\frac1{p-1}-\frac1{p^h-1}\right)
 -\frac{s_p(a)}{p-1}+\frac1{p^h-1}.
\]
For height $2$ the coefficient of $a$ is
\[
 \frac{p}{p^2-1}=\frac1p+\frac1{p^3}+\frac1{p^5}+\cdots.
\]
\end{proposition}
\begin{proof}
Differentiating $\lambda(e(z))=z$ gives $e'(z)=b(e(z))$.
Consequently, for each $a$,
\[
 (D^af)(0)=a!\,[z^a]f(e(z)).
\]
In a term of $[z^a]e(z)^j$, with $1\le j\le a$, the
coefficient is a product of $j$ coefficients whose indices sum
to $a$. AG-FG-EXP bounds its valuation by $-(a-j)/(q-1)$,
which is at least $-(a-1)/(q-1)$. The integral coefficients of
$f$ and the ultrametric inequality give the same bound for
$[z^a]f(e(z))$. Rounding to integral valuations and multiplying
by $a!$ gives the second entry in the maximum. The first follows
because $D$ preserves $R[[t]]$.

Finally, counting multiples of $p^j$ among $1,\ldots,a$ gives
$v(a!)=\sum_{j\ge1}\lfloor a/p^j\rfloor
=(a-s_p(a))/(p-1)$. Subtracting $(a-1)/(q-1)$ gives the
displayed digit-sum bound. The height-$2$ identity is the
geometric-series sum. Since
$s_p(a)\le(p-1)(1+\lfloor\log_pa\rfloor)$, it also gives the
claimed linear gain with an explicit logarithmic error.
\end{proof}

\begin{proposition}[AG-FG-PPOWER: the exact restricted derivation]
With the same coordinate and normalization, reduction modulo $p$
satisfies
\[
 \overline D^{p}=\overline{a_p}\,\overline D
             \quad\hbox{on }(R/p)[[t]].
\]
Thus the $p$-th power is zero for $h\ge2$ and nonzero for $h=1$.
\end{proposition}
\begin{proof}
In characteristic $p$, the iterated Leibniz formula shows that
the $p$-th power of a derivation is a derivation. It is invariant
under the formal-group translations because $\overline D$ is.
Every invariant derivation is a scalar multiple of
$\overline D$: an invariant vector field $c(t)\partial_t$
satisfies $c(F(t,u))=\partial_1F(t,u)c(t)$, and setting $t=0$
gives $c(u)=b(u)c(0)$. It remains to compute the scalar
$(D^pt)(0)$ modulo $p$.

The coefficients $e_i$ are integral for $1\le i<p$ by
AG-FG-EXP, since $q\ge p$. In the coefficient recursion for
$n=p$, every summand with $2\le j<p$ lies in $pR$, and the
$j=p$ term is $a_p$. Hence
\[
 pe_p\equiv-a_p\pmod p.
\]
As $(D^pt)(0)=p!e_p$ and $(p-1)!\equiv-1\pmod p$,
the scalar is $a_p$ modulo $p$. The factorial congruence follows
by pairing inverse elements in $\mathbb F_p^\times$ and holds
also at $p=2$. The height definition gives the last assertions.
\end{proof}

\begin{proposition}[AG-FG-PRODUCT: independent formal factors]
Suppose the integral formal group is a product of one-dimensional
groups as above, of heights $h_1,\ldots,h_r$, with commuting
normalized invariant derivations $D_1,\ldots,D_r$. For
$f\in R[[t_1,\ldots,t_r]]$ and integers $a_i\ge0$, put
\[
 B_i(0)=0,\qquad B_i(a)=
 \max\left\{0,v(a!)-
       \left\lfloor\frac{a-1}{p^{h_i}-1}\right\rfloor\right\}
          \quad(a\ge1).
\]
Then
\[
 v\bigl((D_1^{a_1}\cdots D_r^{a_r}f)(0)\bigr)
       \ge\sum_{i=1}^r B_i(a_i).
\]
\end{proposition}
\begin{proof}
Regard $f$ as a power series in all variables other than $t_1$,
with coefficients in $R[[t_1]]$. Apply AG-FG-DER to each
coefficient and evaluate at $t_1=0$. The result lies in
$p^{B_1(a_1)}R[[t_2,\ldots,t_r]]$. Divide out this scalar and
repeat for each remaining coordinate. Evaluation and the other
coordinate derivations commute, proving the bound. If $a_i=0$,
the corresponding operation is simply evaluation and preserves
integrality. An isogeny decomposition over $K$ alone is not an
integral product decomposition and does not meet this hypothesis.
\end{proof}

\paragraph{AG-FG-THETA (what transfers to a theta calculation).}
Alpöge's note, step 2, PDF p.~3, asserts the height-$2$ exponent
for the normalized derivatives of a CM theta value. AG-FG-DER
supplies that exponent whenever the theta section, in a specified
local trivialization and integral parameter, is an element of
$R[[t]]$, and the operator is the integral invariant derivation
above. The period and trivialization comparison must be supplied
for that theta object. More generally, for constants
$u,b\in K^\times$ and $\widetilde D=uD$,
\[
 v\bigl((\widetilde D^{a}(bf))(0)\bigr)
  =v(b)+a\,v(u)+v\bigl((D^af)(0)\bigr),
\]
so their effect on the bound is explicit. A nonconstant change of
trivialization or a connection requires its own comparison.
The supplied Lean map, step-2 row, reports only a low-confidence
match to CM-curve/invariant-derivation declarations; no new exact
Lean lemma identifier is asserted here.

This proves the formal-group mechanism in Fable Section 2 and
shrinks its general-height gate to the actual theta integrality
and normalization data. Its archive map is inherited explicitly
from that record: cote 54 scans 5--7, the restricted $p$-Lie
algebra and Cartier-duality discussion; cotes 4--9 for crystals;
cotes 136--137, cote 16 scan 14, and cote 47 scan 9 for the
divided-power line. These additional sheets have not been newly
image-checked here. The theorem does not by itself place a theta
section on the crystalline site or replace a theta-lift chain.
It uses neither A nor M.

\subsection{Integral elliptic functions and actual special-value divisibility}
\label{ag:elliptic-integrality}

There is a concrete sector in which the formal integrality required
by AG-FG-THETA can be proved directly. The normalizations below agree
with the regularized Weierstrass functions used by Bannai--Kobayashi,
\emph{Integral structures on $p$-adic Fourier theory}, Ann. Inst.
Fourier 66 (2016), 521--550, Section 5, especially printed pp.~543--546
(Numdam PDF pp.~24--27). Their Theorems 5.1 and 5.8 locate the stronger
classical Katz congruences. The following explicit floor bound is a
direct application of the general formal-height argument proved here;
it does not assume a Lubin--Tate normal-form coordinate.

\begin{proposition}[AG-THETA-REG: integral regularizations]
Let $R$ be a complete unramified discrete valuation ring of mixed
characteristic $(0,p)$, $p\ge5$, with $v(p)=1$. Let $E/R$ have good
reduction and equation
\[
 y^2=4x^3-g_2x-g_3,\qquad g_2,g_3\in R,
\]
with unit discriminant. Use the parameter $t=-2x/y$, invariant
differential $\omega=dx/y$, and invariant derivation $D$ dual to
$\omega$. Let the formal group have height $h\in\{1,2\}$, and put
\[
 d_h(a)=\max\left\{0,\ v(a!)-
                   \left\lfloor\frac{a-1}{p^h-1}\right\rfloor\right\}
 \quad(a\ge1).
\]
Suppose $[b]\in\operatorname{End}_R(E)$ satisfies
$[b]^*\omega=b\omega$ with $b\in R^\times$. Then
\[
 f_b(t)=x(t)-b^2x([b](t))\in R[[t]],\qquad
 v(D^a f_b(0))\ge d_h(a).
\]
Define the normalized Weierstrass coefficients $\beta_m$ by
\[
 x(\exp_F z)=z^{-2}+\sum_{a\ge0}\beta_{a+2}\frac{z^a}{a!}.
\]
For $a\ge1$ one has the exact identity and bound
\begin{equation}\label{eq:ag-bh-bound}
 D^a f_b(0)=(1-b^{a+2})\beta_{a+2},\qquad
 v((1-b^{a+2})\beta_{a+2})\ge d_h(a).
\end{equation}
In particular
$v(\beta_m)\ge d_h(m-2)-1-v_p(m)$ for $m\ge3$.
If $E$ has an action of the full ring of integers of an imaginary
quadratic CM field over $R$, $p$ is
inert in that CM field, and $h=2$, then the stronger conclusion
$v(\beta_m)\ge d_2(m-2)$ holds whenever $p^2-1\nmid m$.
\end{proposition}
\begin{proof}
Write $x=t^{-2}u$. The equation of the curve gives
\[
 u=1+\frac{g_2t^4}{4u}+\frac{g_3t^6}{4u^2}.
\]
Its formal solution has $u\in1+t^4R[[t^2]]$, by coefficient
recursion, since the derivative of the defining equation with
respect to $u$ at $(t,u)=(0,1)$ is a unit. Thus $x(t)$ has integral
Laurent coefficients, leading term $t^{-2}$, and no other negative
power. The endomorphism has integral odd expansion
$[b](t)=bt+O(t^3)$: oddness follows from its compatibility with
inversion, which is $t\mapsto-t$. Since $b$ is a unit,
$b^2x([b](t))$ has the same sole negative term $t^{-2}$.
Their difference is therefore integral and regular in the formal
parameter. The derivation is normalized by $\omega/dt=1$ at zero,
so AG-FG-DER applies and gives the displayed derivative bound.

In characteristic zero the formal logarithm linearizes every
endomorphism: $[b](\exp_F z)=\exp_F(bz)$. Substitution into the
Laurent expansion cancels the pole and gives
\eqref{eq:ag-bh-bound} coefficient by coefficient.
Take now the integer multiplication endomorphism $b=1+p$.
For odd $p$, binomial expansion gives
$v_p((1+p)^m-1)=1+v_p(m)$: for $p\nmid m$ the first term
has valuation one and all later terms have larger valuation;
each further raising to the $p$th power increases the valuation
by one. This proves the uniform bound for $\beta_m$.

In the inert CM case, the residues of the CM integers contain
$\mathbb F_{p^2}$, so choose $b$ whose residue generates
$\mathbb F_{p^2}^{\times}$. It acts as the stipulated
endomorphism and is a unit. When $p^2-1\nmid m$, the factor
$1-b^m$ is a unit, proving the stronger bound.
\end{proof}

\begin{proposition}[AG-THETA-TORS: torsion translations close another integrality gate]
Keep the elliptic model and normalization above. If $P\in E(R)$
reduces away from the origin, then
\[
 x(P+_E t)\in R[[t]],\qquad
 v(D^a x(P))\ge d_h(a)\quad(a\ge1).
\]
Every nonzero torsion point of order prime to $p$ satisfies the
reduction condition after a finite unramified extension defining
it. For a model defined over a number field and a compatible
complex uniformization $E(\mathbb C)=\mathbb C/\Lambda$ with
$\omega=dz$, write $z_P$ for a lift of $P$. Then
\begin{equation}\label{eq:ag-eisenstein-derivative}
 D^a x(P)=(-1)^a(a+1)!
     \sum_{\gamma\in\Lambda}(z_P+\gamma)^{-a-2}.
\end{equation}
The valuation bound therefore applies to these actual algebraic
normalized Eisenstein values.
\end{proposition}
\begin{proof}
The rational function $x$ is regular on $E\setminus\{O\}$.
Translation by $P$ takes the formal neighbourhood of $O$ into
this open subscheme when $\bar P\ne O$. Pullback of $x$ is
therefore an element of $R[[t]]$. Translation commutes with the
invariant derivation, so AG-FG-DER proves the asserted bound at $P$.
For an integer $n$ prime to $p$, multiplication by $n$ on the
formal group has unit linear term and hence an inverse formal
series over $R$. It has no nonzero $n$-torsion in that
neighbourhood. Thus a nonzero prime-to-$p$ torsion point cannot
reduce to $O$. Moreover $E[n]/R$ is finite \'etale, so its
points become rational over a finite unramified extension;
the unramified hypothesis needed in AG-FG-DER is preserved.

In complex coordinates $x=\wp_\Lambda(z)$ and $D=d/dz$.
Differentiating the defining series for $\wp$ $a\ge1$ times
gives \eqref{eq:ag-eisenstein-derivative}; the resulting series
converges absolutely since $a+2>2$. The left side is algebraic
and fixes the period and differential normalization of the
right side. Its valuation is thus well defined through the
chosen $p$-adic embedding; no valuation of a free-standing
complex period is being asserted.
\end{proof}

\begin{proposition}[AG-THETA-HECKE: an explicitly normalized Hecke-value family]
Let the global elliptic curve have CM by $\mathcal O_K$, where
$K$ is imaginary quadratic of class number one. Suppose its
period lattice for the chosen algebraic differential is
$\Lambda=\Omega\mathcal O_K$. Put
$w_K=|\mathcal O_K^\times|$. For $m\ge3$ divisible by $w_K$,
the conductor-one algebraic Hecke character
$\psi_m((\alpha))=\alpha^m$ is well defined, and
\[
 \beta_m=w_K(m-1)!\,\Omega^{-m} L(m,\psi_m).
\]
Consequently all bounds for $\beta_m$ in AG-THETA-REG apply
to this displayed algebraic normalization. These values are
already in the region of absolute convergence, so the
divisibility conclusion uses neither A nor M.
\end{proposition}
\begin{proof}
Every nonzero integral ideal of $\mathcal O_K$ is principal,
and changing its generator multiplies $\alpha^m$ by $1$.
Thus the formula defines the stated Hecke character. Its value
on an ideal has complex absolute value $N\mathfrak a^{m/2}$,
so the series at $s=m$ is absolutely convergent for $m>2$.
Counting the $w_K$ generators of each ideal gives
\[
 \begin{aligned}
 L(m,\psi_m)&=\frac1{w_K}
   \sum_{\alpha\in\mathcal O_K\setminus\{0\}}
          \frac{\alpha^m}{(\alpha\bar\alpha)^m}\\
  &=\frac1{w_K}
   \sum_{\alpha\in\mathcal O_K\setminus\{0\}}\alpha^{-m}.
 \end{aligned}
\]
where the final equality uses the bijection
$\alpha\mapsto\bar\alpha$ of $\mathcal O_K$.
The Laurent expansion of $\wp_\Lambda$ has coefficient
$(m-1)\sum_{\gamma\in\Lambda\setminus\{0\}}\gamma^{-m}$
at $z^{m-2}$; this follows by expanding each summand in its
defining series, absolutely for $m\ge3$.
Comparison with the definition of $\beta_m$ gives
$\beta_m=(m-1)!\sum_{\gamma\ne0}\gamma^{-m}$, and scaling
the lattice by $\Omega$ proves the identity.
\end{proof}

For example, for $y^2=4x^3-4x$ over $\mathbb Q(i)$ at $p=7$,
take $m=52$. Then $48\nmid52$ and
$d_2(50)=v_7(50!)-\lfloor49/48\rfloor=8-1=7$.
Thus the normalized value
$4\cdot51!\,\Omega^{-52}L(52,\psi_{52})$ is divisible by
$7^7$ in the specified algebraic normalization. This is a
concrete special-value consequence of the formal method.
The larger theta-lift programme still needs its own section,
period and operator comparison; these proofs close the
regularized elliptic-function and prime-to-$p$ torsion sectors.

\begin{proposition}[AG-CM-INTEGRAL-MODEL: existence of the required CM models]
For every imaginary quadratic field $K$, with class number $h_K$
and $w_K$ roots of unity, there are a number field $F\supset K$
and an elliptic scheme with full $\mathcal O_K$-action over
$\mathcal O_F[1/3]$ such that
\[
 [F:K]\le\frac{48h_K}{w_K},\qquad
 F/K\text{ is unramified away from }3.
\]
A complex embedding can be chosen so that its generic curve is
isomorphic to $\mathbb C/\mathcal O_K$. At any prescribed finite
set of places above rational primes $p\ge5$ unramified in $K$,
one may choose one algebraic invariant differential which is
an integral generator at every chosen place. Its complex period
lattice is then $\Omega\mathcal O_K$.

In particular the model and differential hypotheses of
AG-THETA-CLASS/RAY/PCOND can be realized for every $K$ and every
finite set of inert primes $p\ge5$; at those places the completed
base is unramified over $\mathbb Q_p$ and the formal height is two.
This is an existence statement with a specified integral
normalization, not a claim about every twist or every preselected
differential of a given CM curve.
\end{proposition}
\begin{proof}
Use the CM moduli stack $\mathcal M$ of pairs $(E,\iota)$ with
$\iota:\mathcal O_K\to\operatorname{End}(E)$ inducing the
structure map on the Lie algebra. Kudla--Rapoport--Yang,
\emph{A peculiar modular form of weight one}, arXiv:math/9808143,
Section 5, Proposition 5.1 and its two deformation cases,
PDF pp.~23--25, prove the requisite finite étale CM moduli
statement over $\mathcal O_K$. This section explicitly allows
arbitrary imaginary quadratic $K$. Proposition 5.3 and
Corollary 5.4 identify its geometric characteristic-zero
isomorphism classes with $\operatorname{Cl}(K)$ and its coarse
scheme with the Hilbert class field's integer scheme.

Over $\mathcal O_K[1/3]$, add an ordered basis of $E[3]$.
The choices form a finite étale torsor with group
$\operatorname{GL}_2(\mathbb F_3)$, of order $48$.
No nontrivial automorphism of a CM pair fixes such a basis.
Indeed that automorphism is a unit $u\in\mathcal O_K^\times$;
if $u\ne1$, the isogeny $[u-1]$ has degree
$N_{K/\mathbb Q}(u-1)=|u-1|^2\le4$ and cannot kill the
subgroup $E[3]$ of degree $9$. The rigidified moduli object
therefore has no stabilizers and is a finite étale scheme
$Y$ over $\mathcal O_K[1/3]$, with its universal elliptic curve
and CM action. There are exactly $48/w_K$ level classes over
each characteristic-zero CM isomorphism class, so
$\deg(Y/\mathcal O_K[1/3])=48h_K/w_K$.

Choose the component containing the complex point
$\mathbb C/\mathcal O_K$ with any ordered basis of its
$3$-torsion. This component is
$\operatorname{Spec}\mathcal O_F[1/3]$ for a finite extension
$F/K$, since a connected finite étale scheme over this
normal Dedekind base is the corresponding normalized integer
scheme. Its degree is at most the degree of $Y$ and it is
unramified away from $3$. Pull back the universal curve.
At a prime unramified in $K/\mathbb Q$, its local base is
therefore unramified over $\mathbb Q_p$. The height-two
assertion at inert primes is exactly Case 2 of the cited
deformation analysis.

The invariant differential module is an invertible module on
$\mathcal O_F[1/3]$. Choose any nonzero differential in its
one-dimensional generic fibre. At the finitely many chosen
places it differs from a local integral generator by a scalar
of specified integral valuation. Weak approximation supplies
a single element of $F^\times$ with the opposite valuations.
Multiplying the differential by that element makes it a
generator at every chosen place. The chosen complex curve
has principal CM lattice, so scaling its differential changes
the common period $\Omega$ and leaves the lattice in the form
$\Omega\mathcal O_K$. Any two local integral generators
differ by a local unit; the valuation conclusions are invariant
under that remaining unit ambiguity.
\end{proof}

\begin{proposition}[AG-CM-MODEL-DEGREE: an A-only quantitative CM model]
Let $-D$ be a negative fundamental discriminant with
$D\ge10^{10}$, and put $K=\mathbb Q(\sqrt{-D})$.
Under adopted A, the field and CM model in
AG-CM-INTEGRAL-MODEL can be chosen with
\[
 [F:K]\le\frac{48e^\gamma}{\pi}\sqrt D
    \left(\log\log D-\log2+\frac12+
                         \frac1{\log\log D}\right).
\]
The extension remains unramified away from $3$ over $K$,
and the simultaneous local differential normalizations in
that proposition require no increase of its degree.
\end{proposition}
\begin{proof}
Here $w_K=2$. NT-H02, from Lamzouri--Li--Soundararajan,
corrected arXiv:1309.3595v3, Corollary 1.3 and (1.3),
author p.~4, gives
\[
 h_K\le\frac{2e^\gamma}{\pi}\sqrt D
    \left(\log\log D-\log2+\frac12+
                         \frac1{\log\log D}\right)
\]
from Dirichlet GRH, now supplied by A.
Multiply this by $48/w_K=24$ in the proved model-degree
bound. The model construction and weak approximation do
not otherwise change. This is a quantitative statement about
a field carrying a universal level-three CM elliptic curve,
or equivalently a point of this zero-dimensional CM moduli
scheme; it asserts neither a small defining polynomial nor
an algorithm for constructing it within the degree bound.
\end{proof}

\begin{proposition}[AG-THETA-CLASS: the class-number restriction is unnecessary]
Let $K$ be any imaginary quadratic field, with ring of integers
$\mathcal O$ and $w=|\mathcal O^\times|$. Fix a CM elliptic curve
$E_0$ over a number field, with its full $\mathcal O$-action and an
algebraic differential $\omega$ whose complex period lattice is
$\Omega\mathcal O$. At the chosen $p$-adic embedding, suppose that
$p\ge5$, the curve has a good model over the integers $R$ of a
finite unramified extension of $\mathbb Q_p$, the action extends
over $R$, and $\omega$ generates its invariant differentials.
Write $h$ for its formal height and normalize $v(p)=1$.
Use short Weierstrass models normalized by these differentials;
such models exist integrally since $6\in R^\times$.

Let $\psi$ be any conductor-one algebraic Hecke character of
infinity type $(m,0)$, where $m\ge3$; necessarily $w\mid m$.
Choose integral ideal-class representatives $\mathfrak a$ prime
to $p$. The quotient $E_{\mathfrak a}=E_0/E_0[\mathfrak a]$ has
a normalized differential $\omega_{\mathfrak a}$ with pullback
$\omega$ and lattice $\Omega\mathfrak a^{-1}$. With the
coefficients $\beta_m$ defined in AG-THETA-REG for these models,
\begin{equation}\label{eq:ag-cm-class-sum}
 w(m-1)!\,\Omega^{-m}L(m,\psi)
   =\sum_{[\mathfrak a]\in\operatorname{Cl}(K)}
      \psi(\mathfrak a)^{-1}
      \beta_m(E_{\mathfrak a},\omega_{\mathfrak a}).
\end{equation}
All the weights $\psi(\mathfrak a)^{-1}$ are $p$-adic units.
Consequently the valuation of the displayed normalized Hecke
value is at least
\[
 d_h(m-2)-1-v_p(m),
\]
and it is at least $d_2(m-2)$ when $p$ is inert, $h=2$, and
$p^2-1\nmid m$. There is no denominator equal to the class number.
\end{proposition}
\begin{proof}
First we check the integral normalization under ideal isogenies.
Since $N\mathfrak a$ is a unit in $R$, the subgroup
$E_0[\mathfrak a]\subset E_0[N\mathfrak a]$ is finite
\'etale. After a finite unramified extension it is constant;
its nonzero points have integral coordinates and remain distinct
from the origin after reduction, by AG-THETA-TORS. Apply the
normalized quotient construction of V\'elu, \emph{Isog\'enies
entre courbes elliptiques}, C.~R.~Acad.~Sci.~Paris 273 (1971),
A238--A241, formulas (7), (11), and Remark 2; in Ghitza's
translation these occur on PDF pp.~2, 4. Its coefficient formulas
are integral polynomials in the kernel coordinates. Their reductions
construct the quotient by the reduced, still separable kernel,
so the resulting cubic remains nonsingular. The construction
descends to $R$. Its pullback on invariant differentials is the
identity. Thus the quotient has good reduction, its differential
is an integral generator, and its formal group is isomorphic to
that of $E_0$, so has height $h$. The $\mathcal O$-action descends
because it preserves the kernel. Complex uniformization identifies
this quotient with $\mathbb C/(\Omega\mathfrak a^{-1})$ and
identifies its normalized differential with $dz$.

Every nonzero integral ideal in the class $[\mathfrak a]^{-1}$
has the form $(\alpha)\mathfrak a^{-1}$ for exactly $w$
elements $\alpha\in\mathfrak a\setminus\{0\}$, modulo this
unit multiplicity. Its character value and norm are
\[
 \psi((\alpha)\mathfrak a^{-1})
       =\alpha^m\psi(\mathfrak a)^{-1},\qquad
 N((\alpha)\mathfrak a^{-1})
       =\alpha\bar\alpha/N\mathfrak a.
\]
Since $|\psi(\mathfrak b)|=(N\mathfrak b)^{m/2}$, all series
below converge absolutely when $m>2$. Partitioning the Hecke
series by ideal classes gives
\[
 L(m,\psi)=\frac1w\sum_{[\mathfrak a]}
  \frac{(N\mathfrak a)^m}{\psi(\mathfrak a)}
      \sum_{0\ne\alpha\in\mathfrak a}\bar\alpha^{-m}
 =\frac1w\sum_{[\mathfrak a]}\psi(\mathfrak a)^{-1}
      \sum_{0\ne\gamma\in\mathfrak a^{-1}}\gamma^{-m}.
\]
The second equality uses the exact ideal identity
$\bar{\mathfrak a}=(N\mathfrak a)\mathfrak a^{-1}$.
The Laurent-coefficient calculation in AG-THETA-HECKE now proves
\eqref{eq:ag-cm-class-sum} for each lattice separately.

If the order of $[\mathfrak a]$ is $e$ and
$\mathfrak a^e=(\alpha)$, then
$\psi(\mathfrak a)^e=\alpha^m$. Because $\mathfrak a$ is prime
to $p$, $\alpha$ is a unit at every prime above $p$, hence
$v(\psi(\mathfrak a))=0$. Apply AG-THETA-REG to each normalized
quotient model and then use the ultrametric inequality in
\eqref{eq:ag-cm-class-sum}. The field containing the character
values may be ramified at $p$: only their unit valuations are
used. The formal derivative estimates were made on the unramified
models before extending the valuation. Neither averaging over
ideal classes nor dividing by their number occurs.
\end{proof}

\begin{proposition}[AG-THETA-RAY: proper conductor gives a torsion-value formula]
Keep the global CM model, embeddings, differential, and $p$-adic
hypotheses of AG-THETA-CLASS. Let $\psi$ have infinity type
$(m,0)$, $m\ge3$, and a proper conductor ideal
$\mathfrak f\ne\mathcal O$ prime to $p$. Fix the convention
\[
 \psi((\alpha))=\chi_{\mathfrak f}(\alpha)\alpha^m
 \quad ((\alpha,\mathfrak f)=1),
\]
where $\chi_{\mathfrak f}$ is a finite character of
$(\mathcal O/\mathfrak f)^\times$. In particular
$\chi_{\mathfrak f}(u)u^m=1$ for every unit $u$.
Choose integral class representatives $\mathfrak a$ prime to
$p\mathfrak f$, and for every
$c\in(\mathcal O/\mathfrak f)^\times$ choose
$\alpha_{\mathfrak a,c}\in\mathfrak a$ with residue $c$.
Let $E_{\mathfrak a,\mathfrak f}$ have differential-normalized
lattice and point
\[
 \Lambda_{\mathfrak a,\mathfrak f}
    =\Omega\bar{\mathfrak f}\mathfrak a^{-1},\qquad
 P_{\mathfrak a,c}
    =\frac{\Omega\bar\alpha_{\mathfrak a,c}}{N\mathfrak a}
          \pmod{\Lambda_{\mathfrak a,\mathfrak f}}.
\]
These curves have good unramified models with normalized integral
differentials and height $h$; the points are nonzero torsion of
order prime to $p$. Writing $D_{\mathfrak a,\mathfrak f}$ for
the corresponding invariant derivation, one has
\begin{equation}\label{eq:ag-cm-ray-sum}
 \begin{aligned}
 w(m-1)!\,\Omega^{-m}L(m,\psi)
   ={}&(-1)^{m-2}\sum_{[\mathfrak a]}
         \psi(\mathfrak a)^{-1}
       \sum_{c\in(\mathcal O/\mathfrak f)^\times}
          \chi_{\mathfrak f}(c)\\[-2pt]
     &\hspace{25mm}\cdot
       D_{\mathfrak a,\mathfrak f}^{m-2}
          x_{\mathfrak a,\mathfrak f}(P_{\mathfrak a,c}).
 \end{aligned}
\end{equation}
It follows that
\[
 v\bigl(w(m-1)!\,\Omega^{-m}L(m,\psi)\bigr)
       \ge d_h(m-2).
\]
No condition $w\mid m$ or $p^2-1\nmid m$ is needed here.
The expression whose valuation is taken is algebraic by the
displayed identity, with the period fixed by $\omega$.
\end{proposition}
\begin{proof}
Because $\mathfrak a$ is prime to $\mathfrak f$, its residue
map onto $\mathcal O/\mathfrak f$ is surjective, with kernel
$\mathfrak f\mathfrak a$. Thus the choices in the statement
exist, and their cosets are precisely the elements of
$\mathfrak a$ prime to $\mathfrak f$.

The normalized quotient attached to the ideal
$\mathfrak f\mathfrak a$ in the preceding proof has lattice
$\Omega\mathfrak f^{-1}\mathfrak a^{-1}$ and good model of
height $h$. Scale its lattice coordinate by the integer
$N\mathfrak f$. Since
$(N\mathfrak f)\mathfrak f^{-1}=\bar{\mathfrak f}$, this
gives the asserted lattice. On the Weierstrass model the change
is $x'=(N\mathfrak f)^{-2}x$,
$y'=(N\mathfrak f)^{-3}y$, and
$\omega'=(N\mathfrak f)\omega_{\mathfrak f\mathfrak a}$.
The scaling factor is a $p$-adic unit, so integrality, good
reduction, height and the generator property of the differential
are preserved.

Changing $\alpha_{\mathfrak a,c}$ by
$\mathfrak f\mathfrak a$ changes its displayed lift by
$\Omega\bar{\mathfrak f}\mathfrak a^{-1}$, so the point is
well defined. Multiplication by any element of
$\bar{\mathfrak f}$ takes it into that lattice; in particular
its order divides $N\mathfrak f$. It is nonzero, since being
zero would imply
$\alpha_{\mathfrak a,c}\in\mathfrak f\mathfrak a$, contrary
to $c$ being a unit modulo the proper ideal $\mathfrak f$.
These prime-to-$p$ points are defined after a finite unramified
extension and satisfy AG-THETA-TORS.

Now partition the absolutely convergent ideal series first by
class and then by residue. The same $w$-to-one generator count
as before gives
\[
 L(m,\psi)=\frac1w\sum_{[\mathfrak a]}
   \frac{(N\mathfrak a)^m}{\psi(\mathfrak a)}
   \sum_c\chi_{\mathfrak f}(c)
      \sum_{\alpha\in\alpha_{\mathfrak a,c}
                       +\mathfrak f\mathfrak a}\bar\alpha^{-m}.
\]
The unit compatibility in the statement guarantees that each
of the $w$ generators contributes the same summand even when
$w\nmid m$. For the lattice and point just constructed,
AG-THETA-TORS gives
\[
 D_{\mathfrak a,\mathfrak f}^{m-2}
     x_{\mathfrak a,\mathfrak f}(P_{\mathfrak a,c})
  =(-1)^{m-2}(m-1)!\,\Omega^{-m}(N\mathfrak a)^m
       \sum_{\alpha\in\alpha_{\mathfrak a,c}
                        +\mathfrak f\mathfrak a}\bar\alpha^{-m}.
\]
This proves \eqref{eq:ag-cm-ray-sum}. All coefficients in its
finite sum are $p$-adic units. For $\psi(\mathfrak a)$ this
follows by raising to the class order as before; the additional
factor $\chi_{\mathfrak f}(\alpha)$ is a root of unity and
has valuation zero. Each derivative has valuation at least
$d_h(m-2)$, proving the bound. A ramified field of finite-character
values causes no loss: the estimates already hold over unramified
fields of definition of the models and torsion points, and
the scalar weights remain units after any extension.
\end{proof}

These two formulas close the class-number and prime-to-$p$
conductor gates for the elliptic-function sector of Fable
Section 2 and Alp\"oge Step 2. Geometrically the class sum
runs over the prime-to-$p$ CM isogeny orbit, while a proper
ray conductor replaces an origin expansion by evaluation on
finite \'etale torsion sections. The resulting exact periods
belong to these algebraic differentials; no comparison with
an arbitrary higher-dimensional theta section or a theta-lift
lattice is asserted. All special values used here lie in the
region of absolute convergence and require neither A nor M.

\begin{proposition}[AG-FG-POLE: formal torsion and an explicit pole loss]
Let $R,F,D,h$ satisfy AG-FG-EXP and put $q=p^h$, $Q=q-1$.
Extend $v$, still normalized by $v(p)=1$, to an algebraic closure
of the fraction field. The following estimates are valid over
every finite, possibly ramified, extension $R'/R$.
For $g\in R'[[t]]$ one still has
\[
 v((D^ag)(0))\ge d_h(a)\qquad(a\ge1).
\]
Let $P$ be a nonzero formal point with $\delta=v(t(P))>0$,
and let $f(t)\in t^{-d}R[[t]]$, $d\ge1$. If
$\delta\le1/Q$, write $a-1=\kappa Q+\eta$,
$0\le\eta<Q$. Then
\begin{equation}\label{eq:ag-formal-pole-bound}
 v(D^af(P))\ge
 \max\left\{-(a+d)\delta,\quad
       v_p(a!)-\kappa-(\eta+d+1)\delta\right\}.
\end{equation}
If $P$ is formal torsion of exact order $p^r$, $r\ge1$,
then its parameter has the exact valuation
\[
 \delta_r=\frac1{q^{r-1}(q-1)}.
\]
Thus \eqref{eq:ag-formal-pole-bound} applies with $\delta_r$.
\end{proposition}
\begin{proof}
For the first assertion, each coefficient of $[z^a]\exp_F(z)^j$
is a sum of products with indices $n_1+\cdots+n_j=a$.
AG-FG-EXP gives the integer lower bound
\[
 -\sum_i\left\lfloor\frac{n_i-1}{Q}\right\rfloor
 \ge-\left\lfloor\frac{a-j}{Q}\right\rfloor
 \ge-\left\lfloor\frac{a-1}{Q}\right\rfloor.
\]
Multiplication by an integral coefficient of $g$ cannot lower
this bound. Multiplying by $a!$ and also using that $D$ preserves
$R'[[t]]$ proves $d_h(a)$. This proof rounds the exponential
coefficients over the original unramified ring; it does not round
arbitrary valuations in $R'$.

Put $u=t(P)$. The translated series satisfies
$F(u,T)=u+T A(T)$ with $A(T)\in R'[[T]]$.
All coefficients are well defined by convergence in $u$.
Expansion of $(u+TA(T))^{-d}$ shows that the coefficient of
$T^j$ has valuation at least $-(j+d)\delta$.
The remaining factor in $f(F(u,T))$ is integral, so if
$f(F(u,T))=\sum_{j\ge0}g_jT^j$, then
\[
 v(g_j)\ge-(j+d)\delta.
\]
Translation commutes with the invariant derivation. Therefore
\[
 D^af(P)=a!\,[z^a]\sum_{j\ge0}g_j\exp_F(z)^j,
\]
and the product estimate above yields
\[
 v(D^af(P))\ge v_p(a!)+
   \min_{1\le j\le a}
       \left\{-(j+d)\delta-
           \left\lfloor\frac{a-j}{Q}\right\rfloor\right\}.
\]
To compute this minimum, put $a-j=kQ+b$, $0\le b<Q$.
The expression equals
$-(a+d)\delta+k(Q\delta-1)+b\delta$.
Because $Q\delta\le1$, its minimum occurs at
$k=\kappa$, $b=0$, equivalently $j=\eta+1$.
It is $-\kappa-(\eta+d+1)\delta$.
Independently, $D$ preserves integral Laurent series and
increases their pole order by at most one. Thus
$D^af\in t^{-d-a}R[[t]]$ and its value has valuation at least
$-(a+d)\delta$. Combining the two bounds proves
\eqref{eq:ag-formal-pole-bound}.

Finally write
$[p](t)=pt+\sum_{j\ge2}a_jt^j$ as in AG-FG-EXP.
For a point of positive parameter valuation $\delta$, the
terms with $2\le j<q$ have valuation strictly greater than
$1+\delta$, and those with $j>q$ have valuation strictly
greater than $q\delta$. Unless $\delta=1/(q-1)$, the
valuation of $[p](t(P))$ is therefore
$\min\{1+\delta,q\delta\}$. A nonzero point killed by
$p$ must have $\delta=1/(q-1)$, since otherwise a unique
term of least valuation cannot sum to zero.
For exact order $p^r$, $r>1$, induction gives the parameter
valuation of $[p]P$ as $\delta_{r-1}\le1/(q-1)$.
If $\delta\ge1/(q-1)$, the parameter of $[p]P$ would have
valuation at least $1+1/(q-1)$, a contradiction.
Hence $\delta<1/(q-1)$ and $q\delta=\delta_{r-1}$.
This proves the formula for $\delta_r$.
\end{proof}

\begin{proposition}[AG-THETA-PCOND: the inert $p$-conductor case with its denominator]
Keep the hypotheses and normalization of AG-THETA-CLASS, and
suppose $p$ is inert and the formal height is $2$.
Let $\psi$ have infinity type $(m,0)$, $m\ge3$, and conductor
\[
 \mathfrak f=(p)^r\mathfrak f_0,\qquad r\ge1,
 \quad(\mathfrak f_0,p)=1.
\]
Write $a=m-2$, $a-1=\kappa(p^2-1)+\eta$,
$0\le\eta<p^2-1$, and
\[
 \delta_r=\frac1{p^{2(r-1)}(p^2-1)},\qquad
 B_2(a,r)=\max\{-(a+2)\delta_r,\,
       v_p(a!)-\kappa-(\eta+3)\delta_r\}.
\]
Then the exactly normalized algebraic value satisfies
\begin{equation}\label{eq:ag-cm-p-conductor-bound}
 v\bigl(p^{rm}w(m-1)!\,\Omega^{-m}L(m,\psi)\bigr)
 \ge
 \begin{cases}
 d_2(m-2),&\mathfrak f_0\ne\mathcal O,\\
 B_2(m-2,r),&\mathfrak f_0=\mathcal O.
 \end{cases}
\end{equation}
Equivalently, remove $p^{rm}$ from the left by subtracting
$rm$ from the corresponding lower bound. This is an explicit
denominator estimate, not an assertion of unrestricted integrality.
\end{proposition}
\begin{proof}
Choose class representatives $\mathfrak a$ prime to
$p\mathfrak f_0$, and representatives
$\alpha_{\mathfrak a,c}\in\mathfrak a$ of
$c\in(\mathcal O/\mathfrak f)^\times$ as in AG-THETA-RAY.
Use the good normalized models with lattices
$\Omega\bar{\mathfrak f}_0\mathfrak a^{-1}$; for
$\mathfrak f_0=\mathcal O$ these are $E_{\mathfrak a}$.
On them put
\[
 P^{(r)}_{\mathfrak a,c}
   =\frac{\Omega\bar\alpha_{\mathfrak a,c}}
                    {(N\mathfrak a)p^r}
       \pmod{\Omega\bar{\mathfrak f}_0\mathfrak a^{-1}}.
\]
The residue decomposition in AG-THETA-RAY is an identity of
absolutely convergent lattice sums even when the conductor has
a $p$-part. Its lattice there is now $p^r$ times the good
lattice just displayed. Under a lattice scaling by $p^r$,
$x$ scales by $p^{-2r}$ and the invariant derivation by
$p^{-r}$. The derivative of order $m-2$ therefore scales by
$p^{-rm}$, giving the exact identity
\[
 \begin{aligned}
 p^{rm}w(m-1)!\,\Omega^{-m}L(m,\psi)
 ={}&(-1)^{m-2}\sum_{[\mathfrak a]}\psi(\mathfrak a)^{-1}
    \sum_{c\in(\mathcal O/\mathfrak f)^\times}
       \chi_{\mathfrak f}(c)\\[-2pt]
 &\hspace{22mm}\cdot
       D_{\mathfrak a,\mathfrak f_0}^{m-2}
          x_{\mathfrak a,\mathfrak f_0}
                (P^{(r)}_{\mathfrak a,c}).
 \end{aligned}
\]
Each scalar coefficient is a $p$-adic unit, by the class-order
argument already proved. The curve models are still unramified;
the torsion points and finite-character values may require
ramified extensions.

If $\mathfrak f_0$ is proper, then
$[p^r]P^{(r)}_{\mathfrak a,c}$ is the nonzero prime-to-$p$
torsion point attached to the residue of
$\alpha_{\mathfrak a,c}$ modulo $\mathfrak f_0$.
It reduces away from zero. Consequently
$P^{(r)}_{\mathfrak a,c}$ itself reduces away from zero.
Translation of $x$ is integral over the possibly ramified
field defining this point, and the first part of AG-FG-POLE
gives $d_2(m-2)$ without changing that bound.

If $\mathfrak f_0=\mathcal O$, each point has exact order
$p^r$: it is killed by $p^r$, and being killed by $p^{r-1}$
would mean $\bar\alpha_{\mathfrak a,c}\in p\bar{\mathfrak a}$,
impossible because the residue $c$ is a unit modulo $(p)^r$.
The height-two special fibre has no nonzero geometric
$p$-torsion point. Indeed, the formal multiplicity at the
origin of $[p]$ is $p^2$, the full degree of this elliptic
isogeny; hence its geometric kernel has only that point.
Iteration gives the same assertion for $p^r$.
Thus all the displayed points are formal, and AG-FG-POLE gives
their exact parameter valuation $\delta_r$ and the bound
$B_2(m-2,r)$ for $x$, whose pole order is $2$.
Applying the ultrametric inequality to the finite identity
proves both cases of \eqref{eq:ag-cm-p-conductor-bound}.
\end{proof}

The comparison now covers every finite conductor in the inert
CM elliptic sector, with the denominator displayed when the
conductor meets $p$. The ramification affects the torsion points
and scalar values, while the exponential estimate is still
computed on the original unramified formal group. General theta
sections retain their separate integral model, connection and
period comparison; these explicit functions do not supply them.

\begin{proposition}[AG-THETA-NONINT: a precise reduced-theta normalization obstruction]
On the good CM model
\[
 E:\ y^2=4x^3-4x,\qquad t=-2x/y,\qquad\omega=dx/y,
\]
over $R=\mathbb Z_7[i]$, let $\lambda(t)$ be the formal logarithm
for $\omega$. For its complex period lattice, let $\theta(z)$
be the reduced theta function for the divisor $[0]$, normalized
by $\theta'(0)=1$. Then
\begin{equation}\label{eq:ag-theta-nonintegral}
 \frac{\theta(\lambda(t))}{t}
   =1-\frac5{12}t^4+\frac{1411}{2016}t^8+O(t^{12}),
 \qquad v_7\left(\frac{1411}{2016}\right)=-1.
\end{equation}
In particular this normalized formal series is not in $R[[t]]$.
Multiplication by an integral formal unit or an integral change
of local parameter with unit linear coefficient cannot make it
integral.
\end{proposition}
\begin{proof}
The theta convention is the one in Bannai--Kobayashi,
\emph{Algebraic theta functions and the $p$-adic interpolation
of Eisenstein--Kronecker numbers}, arXiv:math/0610163v4,
Example 1.9, PDF pp.~10--11:
\[
 \theta(z)=\exp(-e_2^*z^2/2)\sigma(z).
\]
Here the discriminant $64$ is a $7$-adic unit, and $R$ is
unramified quadratic since $-1$ is not a square modulo $7$.
The automorphism $(x,y)\mapsto(-x,iy)$ acts on $\omega$ by
$i$, giving the asserted CM action and lattice symmetry.
The lattice is stable under multiplication by $i$.
For positive regularization parameter, the weight-two lattice
sum defining $e_2^*$ changes sign under $\gamma\mapsto i\gamma$
and converges absolutely, so is zero. Its continuation is
therefore zero also. Thus $\theta=\sigma$ for this model.
This uses the canonical reduced theta convention, not an
arbitrary exponential modification of a sigma function.

Write $x=t^{-2}u$. The curve equation gives
$u^2-u=t^4$, with $u(0)=1$. Since $y=-2t^{-3}u$,
\[
 \lambda'(t)=\frac{\omega}{dt}
       =1-\frac{tu'}{2u}
       =\frac1{2u-1}
       =(1+4t^4)^{-1/2}.
\]
Integrating the first three binomial terms gives
\[
 \lambda(t)=t-\frac25t^5+\frac23t^9+O(t^{13}).
\]
For completeness the sigma coefficients can be computed from
the curve itself. Put $x(z)=\wp(z)$ in the coordinate with
$\omega=dz$. The differential equation
$x''=6x^2-2$, and the $i$-symmetry, give
\[
 \wp(z)=z^{-2}+\frac15z^2+\frac1{75}z^6+O(z^{10}).
\]
For instance, comparison of the constant and $z^4$ terms
gives $2c_2=12c_2-2$ and
$30c_6=12c_6+6c_2^2$. The defining product for sigma gives
$-(\log\sigma)''=\wp$, with
$\sigma(z)/z=1+O(z^4)$. Consequently
\[
 \log\frac{\sigma(z)}z
      =-\frac{z^4}{60}-\frac{z^8}{4200}+O(z^{12}),
 \qquad
 \sigma(z)=z-\frac{z^5}{60}-\frac{z^9}{10080}+O(z^{13}).
\]
Substitution of the computed logarithm yields the coefficient
of $t^4$ as $-2/5-1/60=-5/12$ and that of $t^8$ as
\[
 \frac23+5\left(-\frac25\right)\left(-\frac1{60}\right)
           -\frac1{10080}=\frac{1411}{2016}.
\]
The denominator is $7\cdot288$ and the numerator is $4$
modulo $7$, proving \eqref{eq:ag-theta-nonintegral}.
All earlier coefficients are $7$-integral.

If this series is multiplied by $u(t)\in R[[t]]^\times$,
the first nonintegral coefficient becomes
$u(0)(1411/2016)$ plus an integral element, and still has
valuation $-1$. Under $t=a_1T+a_2T^2+\cdots$, with
$a_1\in R^\times$ and all $a_j\in R$, its first
nonintegral coefficient becomes
$a_1^8(1411/2016)$ plus an integral element. This also has
valuation $-1$. If the divisor equation is changed from $t$
to $T$, their ratio is itself an integral unit, so the same
conclusion applies to that normalization as well.
\end{proof}

This refutes the blanket transfer that a canonically normalized
reduced theta section on every good inert CM model is automatically
an integral formal function. It identifies the missing integral
trivialization in that transfer. It neither refutes A nor asserts
failure of any stated derivative-divisibility theorem in its
specified normalization. The stronger classical theta theory
may supply different normalizations or estimates. In particular,
the ordinary-prime hypothesis in Bannai--Kobayashi,
Propositions 2.14--2.15, PDF pp.~28--29, is an explicit extra
hypothesis for their formal integrality theorem; it is absent
at the inert prime in this example.

\subsection{A proved Fourier refinement of the supersingular divisibility}

The following argument uses the proved integral Fourier estimates of
Bannai--Kobayashi, not the announcement of a future inert-prime paper in
their earlier theta manuscript.  Its input is the Lubin--Tate distribution
estimate in \textbf{S-AG-BK-FOURIER}, Proposition 5.3(i),(ii), printed
pp.~544--545, together with the Fourier moment and torsion-averaging
identities used in Proposition 5.5, pp.~546--547.  These statements and
their complete displayed proofs have been read.  The exact period
valuation is given in their \S2, printed pp.~526--527.  No assertion
about the full sigma-function radius is substituted for these results.

\begin{proposition}[AG-FG-FOURIER: divisibility from an integral trace family]
Let $p\geq3$, $q=p^2$, $Q=q-1$ and $\alpha=p/Q$.  Let $G$ be a
Lubin--Tate formal group for the unramified quadratic extension of
$\mathbf Q_p$, with an integral invariant differential whose value at the
origin is a unit.  Normalize its derivation $D$ by the formal logarithm.
Let $R\subset\mathcal O_{\mathbf C_p}$ be a complete discrete valuation
ring containing the coefficients of $G$, with $v_p(R\setminus\{0\})
\subset e^{-1}\mathbf Z$.  Suppose that $f_j\in R[[t]]$, $j\geq0$,
satisfy the convergent formal identities
\[
 \frac1q\sum_{\tau\in G[p]} f_j(t\mathbin\oplus\tau)
       =f_{j+1}([p]t).                                      \tag{AG-Fourier-trace}
\]
For every $n\geq1$ put
\[
 b_e(n)=
 \begin{cases}
 \max\{0,e^{-1}\lceil e(n\alpha-1)\rceil\},&Q\nmid n,\\
 \max\{0,e^{-1}\lceil e(n-1)\alpha\rceil\},&Q\mid n.
 \end{cases}
\]
Then $v_p(D^nf_j(0))\geq b_e(n)$ for every $j$.  In particular, for
an unramified coefficient ring,
\[
             v_p(D^nf_j(0))\ \geq\ \left\lfloor\frac{pn}{p^2-1}\right\rfloor.
                                                               \tag{AG-Fourier-floor}
\]
An integral change of parameter or multiplication of the invariant
differential by a unit does not change these valuation assertions.
\end{proposition}

\begin{proof}
Choose a Lubin--Tate period $\varpi$ in the Fourier construction.  Its
valuation is
\[
 v_p(\varpi)=\frac1{p-1}-\frac1{p^2-1}=\alpha.
\]
For an integral series $f$, denote its Fourier distribution on
$\mathcal O_{\mathbf Q_{p^2}}$ by $\mu_f$.  This is a distribution;
it need not be a bounded measure.  Proposition 5.3 gives
\[
 \left|\int_{\mathcal O^\times}x^n\,d\mu_f\right|_p\leq p,
 \qquad
 \left|\int_{\mathcal O^\times}(x^m-x^n)\,d\mu_f\right|_p
       \leq p^{-l+\alpha}
 \quad(m\equiv n\bmod p^lQ).                              \tag{AG-Fourier-est}
\]
The indicator of $\mathcal O^\times$ corresponds under the transform to
subtracting the average over $G[p]$.  Differentiating its transform and
using (AG-Fourier-trace) therefore gives, for $a_{j,n}=D^nf_j(0)$,
\[
 c_{j,n}:=a_{j,n}-p^n a_{j+1,n}
       =\varpi^n\int_{\mathcal O^\times}x^n\,d\mu_{f_j}.    \tag{AG-Fourier-moment}
\]
In degree zero its mass is $f_j(0)-f_{j+1}(0)\in R$.  Consequently,
when $Q\mid n$, the second estimate with exponents $n,0$ and $l=0$
shows that the $n$th moment has valuation at least $-\alpha$:
the constant mass is integral, and the difference has that valuation
bound.  For $Q\nmid n$ the first estimate instead gives valuation
at least $-1$.  After multiplication by $\varpi^n$, these are respectively
$(n-1)\alpha$ and $n\alpha-1$.  Each $c_{j,n}$ belongs to $R$, so rounding
up in its actual value group gives the corresponding lower bound in the
definition of $b_e(n)$.  The additional bound zero follows directly from
$a_{j,n},a_{j+1,n}\in R$, since $D$ is an integral derivation.

Iterate (AG-Fourier-moment):
\[
 a_{j,n}=\sum_{i=0}^{N-1}p^{in}c_{j+i,n}+p^{Nn}a_{j+N,n}.
\]
The last term has valuation at least $Nn$, uniformly in $j$, while all
terms of the sum have valuation at least $b_e(n)$.  Letting $N$ tend to
infinity proves the claim.  For $e=1$ and $Q\nmid n$, $n\alpha$ is
nonintegral because $(p,Q)=1$, hence $\lceil n\alpha-1\rceil
=\lfloor n\alpha\rfloor$.  For $Q\mid n$, $n\alpha$ is an integer
and $0<\alpha<1$, hence $\lceil(n-1)\alpha\rceil=n\alpha$.
Finally an integral parameter isomorphism preserves the coefficient
lattice, and rescaling the differential by a unit multiplies an $n$th
derivative by a unit to the power $n$.
\end{proof}

\begin{proposition}[AG-THETA-FOURIER: the exact floor bound for the CM values]
Use the good full-CM models and the algebraic differential in
AG-CM-INTEGRAL-MODEL, at an inert prime $p\geq5$.  Write
$B_p(n)=\lfloor pn/(p^2-1)\rfloor$.
For a CM endomorphism $b$ with $p\nmid N(b)$ and every $m\geq3$,
\[
 v_p\bigl((1-b^m)\beta_m(E,\omega)\bigr)\geq B_p(m-2).
                                                               \tag{AG-Fourier-BH}
\]
For $P$ defined over an unramified extension with $\bar P\ne O$,
\[
                  v_p(D^n x(P))\geq B_p(n)\qquad(n\geq1).   \tag{AG-Fourier-point}
\]
Thus the conductor-one identity of AG-THETA-CLASS has the improved
lower bounds
\[
 v_p\bigl(w(m-1)!\Omega^{-m}L(m,\psi)\bigr)\geq
 \begin{cases}
 B_p(m-2),&p^2-1\nmid m,\\
 B_p(m-2)-1-v_p(m),&p^2-1\mid m.
 \end{cases}                                               \tag{AG-Fourier-class}
\]
For the proper prime-to-$p$ conductor identity of AG-THETA-RAY the
lower bound is $B_p(m-2)$, with no weight congruence exception.
For the mixed conductor $\mathfrak f=(p)^r\mathfrak f_0$ of
AG-THETA-PCOND, $\mathfrak f_0\ne\mathcal O_K$, let $e$ be the
ramification index of a field containing the displayed points.  The
valuation of its left side
$p^{rm}w(m-1)!\Omega^{-m}L(m,\psi)$ is at least
\[
                         \max\{d_2(m-2),b_e(m-2)\}.         \tag{AG-Fourier-mixed}
\]
The pure $p$-conductor pole loss remains the one in AG-THETA-PCOND.
\end{proposition}

\begin{proof}
Full CM is relevant here.  At an inert prime the formal
$\mathcal O_{K,p}$-module has relative height one.  Over the completed
maximal unramified extension it is isomorphic to a Lubin--Tate module,
with unit linear coefficient.  This is the unique-lifting statement
used in \textbf{S-AG-KRY-CM}, Proposition 5.1, Case 2: a fixed
height-one formal $\mathcal O_{K,p}$-module and its fixed action on the
Lie algebra have a unique lift.  Apply it to the special fibre of a
Lubin--Tate module and to that of our full-CM formal group after their
geometric identification.  The differential is a generator, so the
normalization discrepancy is a unit.  This step uses the actual CM
action; it asserts no such classification for arbitrary height-two
deformations.  The unramified extension still has integral value group,
so the rounding in AG-FG-FOURIER is preserved.

The classical multiplication identity, also Lemma 5.4 of
\textbf{S-AG-BK-FOURIER}, is
\[
          \sum_{T\in E[p]}x(Q+T)=p^2x([p]Q).
\]
It holds as an identity of rational functions and hence after every
formal substitution in question.  Since $[b]$ permutes $E[p]$, it
implies that the integral regularization
$f_b=x-b^2x\circ[b]$ of AG-THETA-REG satisfies
(AG-Fourier-trace) with the constant family $f_j=f_b$.
Its $m-2$ derivative is $(1-b^m)\beta_m$, proving
(AG-Fourier-BH).

For (AG-Fourier-point), use
$f_j(t)=x([p]^jP+_E t)$.  The supersingular multiplication map $[p]$
is purely inseparable on the special fibre, so its only geometric
point above $O$ is $O$.  Thus $[p]^j\bar P\ne O$ for all $j$.
Each $f_j$ is an integral formal function because $x$ is regular in
that formal neighbourhood, and the same multiplication identity is
exactly (AG-Fourier-trace).  Apply the proposition with $e=1$.
The argument over a ramified coefficient extension instead gives
$b_e(n)$; no rounding to an integer is made there.

If $Q\nmid m$, choose a CM element $b$ whose residue generates
$\mathbf F_{p^2}^{\times}$.  Then $1-b^m$ is a unit.  If $Q\mid m$,
take $b=1+p$ and use
$v_p((1+p)^m-1)=1+v_p(m)$.  These choices in (AG-Fourier-BH),
followed by the exact finite ideal-class identity already proved in
AG-THETA-CLASS, give (AG-Fourier-class).  All character coefficients
in that identity are units.  They may enlarge the coefficient field,
but do not weaken a valuation bound proved before forming the sum.

The points in AG-THETA-RAY have prime-to-$p$ order, are defined over
unramified extensions, and reduce away from $O$.  Apply
(AG-Fourier-point) to each term of its exact identity.  In the mixed
conductor case of AG-THETA-PCOND the points are possibly ramified,
but still reduce away from $O$ because $[p]^rP$ is nonzero
prime-to-$p$ torsion.  The ramified version just proved and the earlier
independent bound $d_2(m-2)$ both apply.  The exact lattice scaling
$p^{-rm}$ has already been accounted for on the displayed left side.
For pure $p$-conductor, the points reduce to the pole; the integral
trace-family hypothesis is then absent, which is why its earlier
explicit Laurent bound is retained.
\end{proof}

\paragraph{Geometric and method translation.}
The operator in (AG-Fourier-trace) is the normalized trace for the finite
flat CM isogeny $[p]$.  The Fourier identity turns its subtraction into
restriction of a locally analytic distribution to the units.  The
integral function lattice and the unramified value group then provide
the exact divisibility.  This is a concrete trace-to-arithmetic-lattice
comparison relevant to Alpöge step 2 and the Fable \S2 bridge; the
finite-flat isogeny and the integral formal functions are the additional
geometric inputs.  It supplies neither an integral canonical theta
trivialization contradicted by AG-THETA-NONINT nor a theorem about
every automorphic theta lift.  The statement is unconditional and
does not use A or M.  In the cote 1 framework it exemplifies a
geometric trace with a specified integral lattice, in addition to the
analytic trace/determinant statements already catalogued there.

\subsection{A geometric period at nonsplit CM primes}

Here is a proved nonsplit-prime result with its actual period, rather
than an unverified use of the supersingular sigma-radius announcement.
\textbf{S-AG-AI-CM} is Andreatta--Iovita, \emph{Katz type $p$-adic
$L$-functions for primes $p$ non-split in the CM field},
\emph{Commentarii Mathematici Helvetici} 99 (2024), 641--716.
We use its Proposition 5.9, printed pp.~689--690, and the analyticity
proved in Proposition 5.6 and Remark 5.8.  The algebraic normalization
below is the one in its \S2.4; its comparison with complex $L$-values
is the Damerell comparison for the Eisenstein case and the stated
Waldspurger comparison for the cuspidal case.

\paragraph{AG-SHIMURA-PERIOD: the precise classical input.}
Fix an imaginary quadratic field $K$, a level $N\geq5$ satisfying the
Heegner condition $\mathcal O_K/\mathfrak N\simeq\mathbf Z/N\mathbf Z$,
and a nebentype $\epsilon$.  Let $p\geq5$ be inert in $K$ and
$p\nmid N$.  Put $c=p^n d$, $n\geq2$, with $d$ odd and prime to
$pN\operatorname{disc}(K)$.  Let $F$ and $\chi$ satisfy the hypotheses
of \S2.4 of the source: $\chi$ is of its type $(c,\mathfrak N,\epsilon)$,
has infinity type $(k+m,-m)$, $m\geq0$, and has exact $p$-part of
the finite conductor $p^n$.  In the Eisenstein case $F=E_{k,\epsilon}$
is the series of \S2.3.1, $\epsilon$ has parity $k$ and is nontrivial
when $k=2$.  In the cuspidal case $k\geq2$, the discriminant of $K$
is odd, $\chi|_{\mathbf A_{\mathbf Q}^{\times}}=\epsilon\mathbf N^k$,
and the finite local epsilon conditions of its Theorem 2.4 are all
$+1$.  These are retained hypotheses, not supplied by A or M.

Choose the CM elliptic curve of conductor $d$, a generically cyclic
subgroup $C^{(n)}\subset E[p^n]$ of order $p^n$, and
$E^{(n)}=E/C^{(n)}$, as in \S5.1.  Its endomorphism order has
conductor $p^n d$.  Write $\omega^{(n)}$ for an integral generator of
its invariant differentials, and use the isogeny-normalized
differentials on the ideal-class translates.  The algebraic number is
\[
 L_{\rm alg}(F,\chi)=
 \sum_{[\mathfrak a]\in\operatorname{Pic}(\mathcal O_c)}
   \chi_m(\mathfrak a)^{-1}\,
   \delta_k^mF\bigl(\mathfrak a*(E^{(n)},\psi_N^{(n)},
                                      \omega^{(n)})\bigr),
 \qquad \chi_m=\chi\mathbf N^m.                         \tag{AG-AI-algebraic}
\]
The operator $\delta_k^m$ is the weight-$k+2m$ operator obtained from
the Gauss--Manin connection and the CM splitting of the Hodge
filtration.  The source constructs a locally analytic function
$L_p(F,-)$ on its character domain and a period $\Omega_{p,n}$ such
that
\[
 L_{\rm alg}(F,\chi)=\Omega_{p,n}^{k+2m}L_p(F,\chi),
 \qquad
 v_p(\Omega_{p,n})=\frac1{p^{n-1}(p^2-1)}.               \tag{AG-AI-period}
\]
For a varying modular weight, the source requires $n\geq3$ at inert
$p\geq5$; the threshold $n\geq2$ above is for a fixed classical
weight.  This distinction is Definition 4.24, pp.~681--682.

\begin{proof}[Normalization and the source proof]
Lemma 3.1 of the source gives
\[
 v_p\bigl(\operatorname{Hdg}(E^{(n)})\bigr)
       =\frac1{p^{n-1}(p+1)}.
\]
The canonical subgroup supplies a distinguished differential
$\omega_n$ in the canonical Hodge sublattice, with
$\omega_n=\Omega_{p,n}\omega^{(n)}$.  By the definition of that
sublattice, its valuation relative to an integral generator is
$v_p(\operatorname{Hdg})/(p-1)$; this proves the asserted period
valuation.  Lemma 5.1 and Definition 5.2 use the isogeny through
$E/C^{(n)}[p]$ to give the required integral Hodge projection before
evaluation.  Thus the period specifies a lattice comparison, not
merely a nonzero complex scaling constant.

The $p$-depleted modular form, with the source's precise operator
convention in Corollary 4.23, is
\[
 F^{[p]}=F-VUF
       =F-a_pVF+\epsilon(p)p^{k-1}V^2F,
\]
where $UV=1$.  Definition 5.7 defines $L_p$ by the finite CM sum
for $\delta_k^mF^{[p]}$ evaluated with $\omega_n$, divided by
$| (\mathcal O_K/\mathfrak N)^\times |$ when its larger adelic
class set $H(c,\mathfrak N)$ is used.  Keep this factor: it is what
identifies the adelic sum with (AG-AI-algebraic).
The $V$ and $V^2$ terms factor through conductor $p^{n-1}d$ or
$p^{n-2}d$.  Their $\chi$-weighted sums vanish, because the restriction
of the primitive conductor-$p^n$ character to the order-$p$ kernel
of the preceding class-group map is nontrivial.  This is the last
paragraph of the complete proof of Proposition 5.9.  The remaining
sum is the one for $F$.  Replacing $\omega_n$ by
$\omega^{(n)}$ multiplies that weight-$k+2m$ evaluation by
$\Omega_{p,n}^{k+2m}$, giving (AG-AI-period).
Finally, Proposition 5.6 proves analyticity in local weight
coordinates by applying the Hodge projection to the constant
coefficient in the marked-section expansion; Remark 5.8(1) gives
the fixed-weight version used here.  The underlying overconvergent
Gauss--Manin iteration is the source's cited construction, not a
new assertion proved from GRH.
\end{proof}

\begin{proposition}[AG-SHIMURA-SLOPE: a uniform divisibility and interpolation test]
Fix the data of AG-SHIMURA-PERIOD and a closed analytic weight disk
$\mathcal D$ contained in one of its character branches on which
$L_p(F,-)$ is represented by a convergent power series.  There is a
finite constant $C_{F,\mathcal D}\geq0$, independent of $m$, such
that every classical point $\chi\in\mathcal D$ of type $(k+m,-m)$
satisfies
\[
 v_p L_{\rm alg}(F,\chi)
       \geq\frac{k+2m}{p^{n-1}(p^2-1)}-C_{F,\mathcal D}.  \tag{AG-AI-slope}
\]
For any sequence of such points converging in the weight topology,
with ordinary integer weights $m\longrightarrow+\infty$, the
unnormalized values tend to zero $p$-adically.  If the limiting
classical value is nonzero, these unnormalized values have no
continuous interpolant agreeing at that limit.  In the Eisenstein
case every classical point has a sequence with unchanged local-unit character data
and exact finite conductor of this form.
\end{proposition}

\begin{proof}
Write the convergent series on the chosen closed disk as
$L_p(F,T)=\sum a_j(T-T_0)^j$, of radius $r$.  Its Gauss bound
$M=\sup_j|a_j|_p r^j$ is finite.  If it is nonzero, take
$C_{F,\mathcal D}=\max\{0,\log_p M\}$; if it is identically zero,
take $C_{F,\mathcal D}=0$ and interpret $v_p(0)=+\infty$.
Then $v_pL_p(F,\chi)\geq-C_{F,\mathcal D}$ everywhere on the
disk.  Taking valuations in (AG-AI-period) proves (AG-AI-slope).
Its right side tends to infinity with $m$.  Continuity at the
limit would force the limiting unnormalized value to be zero,
which proves the claimed obstruction when that value is nonzero.

For the last assertion let $h$ be the class number and $w$ the
number of roots of unity of $K$, and put $H=hw$.  For a fractional
ideal $\mathfrak a$ choose a generator $\alpha$ of $\mathfrak a^h$
and put $\psi(\mathfrak a)=\alpha^{-w}$.  Changing the generator
multiplies it by a $w$th root of unity, so this is well defined;
multiplication of generators also proves multiplicativity.
Moreover $\psi((z))=z^{-H}$.  In the source's adelic infinity-type
convention this is a conductor-one algebraic Hecke character of
type $(H,0)$.  The character
$\eta=\psi/\overline\psi$ has type $(H,-H)$ and conductor one.
On ideals prime to $p$ its values are units of the unramified field
$K_p$.  Since its residue field has order $p^2$,
\[
      \eta^{(p^2-1)p^j}\longrightarrow1
\]
uniformly in the unit-valued $p$-adic character coordinates.  This
also follows directly from $v_p(u^{p^2-1}-1)\geq1$ for
$u\in\mathcal O_{K_p}^{\times}$ and successive $p$th powers.
For a fixed classical $\chi_0$, set
$\chi_j=\chi_0\eta^{(p^2-1)p^j}$.  Its restrictions to the local unit groups and its exact
finite conductor are unchanged, while its weight is
$m_j=m_0+H(p^2-1)p^j$.  Hence $\chi_j\to\chi_0$ and
$m_j\to+\infty$.  In the Eisenstein case these characters retain
all the stated admissibility conditions.  They eventually belong
to any sufficiently small closed analytic weight disk about
$\chi_0$, proving the final assertion.  We impose no assertion here
that a cuspidal local epsilon condition is preserved by an
arbitrary deformation without checking that condition.
\end{proof}

\paragraph{Place in the programme.}
The geometric object is the Hodge bundle of a modular curve at a CM
point of nonmaximal order, with its canonical subgroup and the
Gauss--Manin connection.  The relevant comparison is between an
integral differential and a canonical Hodge sublattice after a cyclic
$p$-power isogeny.  Its positive, explicitly computed valuation
explains both the divisibility and the normalization required for
$p$-adic interpolation.  This gives an actual Shimura-geometric
continuation of the theta/divisibility discussion in Alpöge step 2
and Fable \S2.  It is compatible with the formal-trivialization
obstruction in AG-THETA-NONINT.  Neither the interpolation theorem
nor its slope consequence uses A or M.  In the cuspidal case the
complex value is the specific Rankin value in the source's
Waldspurger theorem; it is not thereby assigned Hecke GRH or
Artin holomorphy.

\begin{proposition}[AG-SHIMURA-RAM: the ramified period and its exact slope]
Keep the Heegner, nebentype, parity and, where applicable, cuspidal
epsilon hypotheses of AG-SHIMURA-PERIOD, but let $p\geq5$ be
ramified in $K$, with unique prime $\mathfrak p$ above $p$.
Let $c=p^n d$, where $n\geq2$ and $d$ is odd and prime to
$pN\operatorname{disc}(K)$.  Restrict to characters whose exact
$p$-part of the conductor ideal is
$(p)^n=\mathfrak p^{2n}$.  No assertion about an odd
$\mathfrak p$-exponent is implicit here.

Start with $E$ having CM by $\mathcal O_d$ over the integers of
a finite extension of $\mathbb Q_p$.  Choose a subgroup
$C^{(n)}\subset E[p^n]$ which is cyclic of order $p^n$ on the
generic fibre, and whose generic-fibre $p$-torsion line differs
from $E[\mathfrak p]$.  Put $E^{(n)}=E/C^{(n)}$.
Its CM order has conductor $p^n d$.  With an integral invariant
differential $\omega^{(n)}$, the canonical differential is
$\omega_n=\Omega_{p,n}\omega^{(n)}$, and
\[
 v_p(\Omega_{p,n})=\frac1{2p^n(p-1)},\qquad
 L_{\rm alg}(F,\chi)=\Omega_{p,n}^{k+2m}L_p(F,\chi).
                                                        \tag{AG-AI-ram}
\]
Here $L_{\rm alg}$ uses precisely (AG-AI-algebraic), and $L_p$
uses the source's ramified Hodge projection.  For every closed
analytic character disk $\mathcal D$ in its domain there is a
finite $C_{F,\mathcal D}\geq0$ such that
\[
 v_p L_{\rm alg}(F,\chi)
       \geq\frac{k+2m}{2p^n(p-1)}-C_{F,\mathcal D}.       \tag{AG-AI-ram-slope}
\]
The threshold $n\geq2$ works both for a fixed classical modular
weight and for the source's analytic finite-slope modular families
at $p\geq5$.  For a fixed Eisenstein series, every classical point
has an explicit sequence of admissible characters converging to
it while $m\to+\infty$, preserving its local-unit character data
and exact conductor.  Consequently a nonzero algebraic value at
the limit again obstructs continuous interpolation without the
period normalization.
\end{proposition}

\begin{proof}
We first spell out the transversality condition and the CM order.
Over an algebraic closure the $p$-adic CM Tate lattice of $E$ is
free of rank one over $\mathcal O_{K_{\mathfrak p}}$.
Thus identify its $p^n$-torsion with
$p^{-n}\mathcal O_{K_{\mathfrak p}}/\mathcal O_{K_{\mathfrak p}}$.
Write a generator of the chosen cyclic group as
$p^{-n}\beta$.  The line $E[\mathfrak p]$ in $E[p]$ is
the nonunit line $\mathfrak p/(p)$, after this identification.
Our transversality condition is therefore exactly that $\beta$
is a unit.  It can be achieved after a finite extension by
choosing any such generator and taking its finite flat closure.
The quotient lattice is
\[
 \mathcal O_{K_{\mathfrak p}}+\mathbb Z_p p^{-n}\beta
 \sim \mathcal O_{K_{\mathfrak p}}+p^{-n}\mathbb Z_p
 =p^{-n}(\mathbb Z_p+p^n\mathcal O_{K_{\mathfrak p}}),
\]
where the homothety uses $\beta^{-1}$.  The multiplier ring of
the last lattice is exactly the order
$\mathbb Z_p+p^n\mathcal O_{K_{\mathfrak p}}$: membership is
necessary by applying a multiplier to $1$, and sufficient since
this is a ring.  The prime-to-$p$ local orders are unchanged by
the isogeny.  This proves the claimed global conductor.
The transversality condition is on generic fibres; it must not
be read as disjointness of the special fibres of finite flat
group schemes at this supersingular point.

Andreatta--Iovita, Lemma 3.2, printed pp.~659--660, gives
\[
 E[\mathfrak p]\text{ canonical},\qquad
 v_p\operatorname{Hdg}(E)=\frac12,\qquad
 v_p\operatorname{Hdg}(E^{(n)})=\frac1{2p^n}.
\]
Its proof uses the $p$-divisible CM endomorphism given by a
uniformizer of $K_{\mathfrak p}$ and the Katz--Lubin canonical
subgroup theorem.  Lemma 6.1, printed p.~692, supplies the
integral Hodge projection needed for evaluation.  Unlike the
inert case it includes the different, whose $p$-valuation is
$1/2$.  Its reduced numerical inequality is
$3p^2-p-1\leq p^3$, which holds for the present $p\geq5$.
The canonical Hodge sublattice has valuation
$v_p\operatorname{Hdg}(E^{(n)})/(p-1)$ relative to an integral
generator.  This proves the period valuation.

Propositions 6.4 and 6.6, printed pp.~693--694, establish
analyticity and (AG-AI-ram), respectively; their proofs reduce
to the full arguments of Propositions 5.6 and 5.9, with the
ramified projection just described.  Retain the precise
depletion $F^{[p]}=F-VUF$, the adelic class-sum denominator,
and primitive-character cancellation as in our preceding
normalization proof.  Definition 4.24 gives $n\geq2$ in the
ramified case even for the varying modular weight.
The finite Gauss bound for $L_p$ on $\mathcal D$, exactly as
in AG-SHIMURA-SLOPE, now gives (AG-AI-ram-slope).

For completeness use the same conductor-one character
$\eta=\psi/\overline\psi$ of type $(H,-H)$, where $H=hw$,
constructed in AG-SHIMURA-SLOPE.  Its unit values now lie in
the ramified quadratic field $K_{\mathfrak p}$, whose residue
field is $\mathbb F_p$.  Every unit $u$ satisfies
$v_p(u^{p-1}-1)\geq1/2$.  If $v_p(a)\geq1/2$ and $p\geq5$,
the binomial expansion gives
\[
 v_p((1+a)^p-1)\geq v_p(a)+1:
\]
the intermediate binomial coefficients are divisible by $p$,
and $p v_p(a)\geq v_p(a)+1$.  Induction proves uniformly that
$v_p(u^{(p-1)p^j}-1)\geq j+1/2$.
Consequently
\[
 \chi_j=\chi_0\eta^{(p-1)p^j}\longrightarrow\chi_0,
 \qquad m_j=m_0+H(p-1)p^j\longrightarrow+\infty.
\]
Conductor-one twisting changes no local-unit restrictions or
finite conductor, and in the Eisenstein case preserves all the
stated admissibility conditions.  The sequence eventually lies
in the chosen disk, so (AG-AI-ram-slope) proves the final claim.
No integrality rounding is used: the differential and period
may lie over ramified extensions and have fractional valuation.
\end{proof}

The ramified construction concerns the same Hodge-bundle and
Gauss--Manin comparison as the inert construction, but the
different and the canonical CM $\mathfrak p$-subgroup change
its valuation.  It is an unconditional geometric contribution
to the same Alpöge-step-2 and Fable-\S2 normalization gate,
not a consequence of A or M.

\subsection{Arthur blocks, the Lefschetz action, and the integral lattice}
\label{ag:arthur}

The supplied Hodge bridge sketch, Route 1.4 and Route 4, compares
Alpöge's added blocks $\mu[t]$ with Lefschetz strings and with the
integral ring of cote 16, scans 14--18, no.~7. The following
statements locate the classical comparison and separate the
additional representation and lattice calculations. The archive
locators here are inherited from that sketch and from
\texttt{RETURN\_DP\_MULT\_cp01\_IMPLICATIONS.md}, Section 6,
DP15, rather than newly checked manuscript images.

\paragraph{AG-ARTHUR-ARCH (the exact classical comparison).}
Arthur, \emph{Unipotent automorphic representations: conjectures},
Astérisque 171--172 (1989), 13--71, Section 9,
Proposition 9.1, printed pp.~62--66 (Clay PDF pp.~50--54),
compares the complex cohomology of the cohomological real packet
of Section 5 with the representation $r_\mu\circ\psi_\infty$.
Here $r_\mu$ is the minuscule representation specified by the
Shimura cocharacter. The equivalence respects the component-group
action, $\mathbb C^\times$, and the Lefschetz $SL_2$.
The Lefschetz action is constructed on printed pp.~60--61.
This is a statement about the sum of relative Lie algebra
cohomologies in the specified packet, after applying $r_\mu$;
it is not a comparison of integral lattices. The global
decomposition (9.1), printed p.~58, explicitly assumes
Conjecture 8.1 in that source. These exact local and global
scopes replace the sketch's unqualified identification.
The entire proof of Proposition 9.1 has been read: it compares
the $\mathbb C^\times$ and component-group characters, then the
Lefschetz weights using Vogan--Zuckerman Proposition 6.19.
That underlying theorem is cited, not re-proved here.

\begin{proposition}[AG-ARTHUR-PARAM: the two $SL_2$ factors]
Let $k$ be nonarchimedean with residue cardinality $q$, and choose
geometric Frobenius $\Phi$ with $|\Phi|=q^{-1}$. Let $\mu$ be an
unramified unitary character with $\mu(\Phi)=\alpha$, and let
$t\ge1$. An Arthur parameter trivial on the Deligne $SL_2$
and equal to $\mu\boxtimes\operatorname{Sym}^{t-1}$ on the
Arthur factor has associated Weil--Deligne representation
\[
 r(\Phi)e_j=\alpha q^{j-(t-1)/2}e_j\quad(0\le j<t),
 \qquad N=0.
\]
Its local factor is
\begin{equation}\label{ag:arthur-factor}
 L(s,r,0)=\prod_{j=0}^{t-1}
       (1-\alpha q^{-(s+(t-1)/2-j)})^{-1}.
\end{equation}
On the same underlying Weil representation, the monodromy chain
$Ne_0=0$, $Ne_j=e_{j-1}$ for $j\ge1$ is a different
Weil--Deligne representation when $t>1$, and has factor
\[
 L(s,r,N)=(1-\alpha q^{-(s+(t-1)/2)})^{-1}.
\]
\end{proposition}
\begin{proof}
Gan--Gross--Prasad, \emph{Branching laws for classical groups:
the non-tempered case}, arXiv:1911.02783v2, Section 2,
PDF p.~8, define the separate Deligne and Arthur factors.
The associated Langlands parameter inserts
$\operatorname{diag}(|w|^{1/2},|w|^{-1/2})$ into the Arthur
factor. On the monomials $x^{t-1-j}y^j$ its eigencharacters
are $\mu|\cdot|^{(t-1)/2-j}$. Triviality of the Deligne
factor gives $N=0$, proving the first description.
For the second description a direct calculation gives
$r(\Phi)Nr(\Phi)^{-1}=q^{-1}N$, the required
Weil--Deligne relation; inertia is trivial.
The local factor is, by definition,
\[
 \det(1-q^{-s}r(\Phi)\mid(\ker N)^I)^{-1}.
\]
For $N=0$ the kernel is the whole space; for the chain it is
$\mathbb C e_0$. This gives the two displayed answers.
The definition is Deligne, \emph{Les constantes des équations
fonctionnelles des fonctions $L$}, LNM 349 (1973),
Section 8.12, printed p.~572 (IAS PDF p.~72);
Wedhorn, arXiv:math/0011210v2, (3.2.1), PDF p.~37,
states the same convention explicitly.
\end{proof}

If a global unitary Hecke character $\mu$ is given, the formal
product of the factors in \eqref{ag:arthur-factor}, at the places
where these hypotheses hold, is
\[
 \prod_{j=0}^{t-1}L^S(s+(t-1)/2-j,\mu).
\]
M therefore places the zeros inherited from nontrivial zeros of
these constituents on the shifted lines $\Re s=1-t/2+j$;
for a Dirichlet character over $\mathbb Q$, A supplies this
deduction. Known poles and trivial zeros are kept separate.
Omitted Euler factors have their explicit divisors as in AG-LOCAL.
This is a product of shifted constituents, not a new pure
weight assertion. In geometric terms, weights of an associated
graded object do not determine its monodromy operator or the
$\ker N$ factor at a degeneration. No occurrence of an arbitrary
formal parameter in a theta lift is asserted.

\begin{proposition}[AG-ARTHUR-MIN: a block can change after $r_\mu$]
In characteristic zero,
\[
 \bigwedge\nolimits^2\operatorname{Sym}^3(\mathbb Q^2)
       \simeq\operatorname{Sym}^4(\mathbb Q^2)\oplus\mathbb Q
\]
as $SL_2$ representations. Thus the second fundamental
representation of $GL_4$, applied to a standard block of length
$4$, gives strings of lengths $5$ and $1$.
\end{proposition}
\begin{proof}
Set $v_j=x^{3-j}y^j$ for $0\le j\le3$ and use
$E=x\partial_y$, $F=y\partial_x$,
$H=x\partial_x-y\partial_y$. Let $w=v_0\wedge v_1$.
It has $H$-weight $4$ and $Ew=0$. The successive nonzero
vectors generated by $F$ are
\[
 \begin{split}
 w,&\qquad 2v_0\wedge v_2,\\
 2v_0\wedge v_3+6v_1\wedge v_2,&\qquad
 12v_1\wedge v_3,\qquad24v_2\wedge v_3.
 \end{split}
\]
Their weights are $4,2,0,-2,-4$, so they are independent;
their span is stable under $E,F,H$ and is the irreducible
highest-weight-$4$ module. For example the commutator identities
give $EF^jw=j(5-j)F^{j-1}w$, and $F^5w=0$.
The vector
$u=v_0\wedge v_3-3v_1\wedge v_2$ is killed by $E,F,H$.
It is independent of the preceding weight-zero vector.
The two submodules therefore have dimensions $5$ and $1$,
with direct sum the six-dimensional exterior square. Finite
rational $SL_2$ representations are determined here by their
Lie algebra action, or directly by exponentiating the nilpotent
operators, so this is the asserted group decomposition.
\end{proof}

The example is an explicit representation calculation; it does
not claim a new automorphic occurrence. Its geometric use is
precise: a Shimura-cohomology calculation must first apply the
representation $r_\mu$ appropriate to that Shimura datum.
Even a correct equality of standard Arthur parameters does not
allow one to omit this step. For a signature where $r_\mu$ is
the standard representation, the obstruction illustrated here
is absent, but the integral comparison below is still required.

\begin{proposition}[AG-LEFSCH-LATTICE: the rational string loses indices]
There are two graded, integral $\mathfrak{sl}_2$-stable lattices
in the same rational irreducible module, with identical
Frobenius characteristic polynomials in every degree and every
power, for which the extreme Lefschetz square has cokernel
$\mathbb Z/2$ and $0$, respectively.
\end{proposition}
\begin{proof}
On $V=\operatorname{Sym}^2(\mathbb Q^2)$ use the operators
$E,F,H$ just defined and set $m_j=x^{2-j}y^j$. Consider
\[
 M=\mathbb Zm_0\oplus\mathbb Zm_1\oplus\mathbb Zm_2,
 \qquad
 N=\mathbb Zm_0\oplus\mathbb Z(2m_1)\oplus\mathbb Z(2m_2).
\]
Both are stable: on $M$,
$Fm_0=2m_1$, $Fm_1=m_2$,
$Em_1=m_0$, $Em_2=2m_1$.
Writing $n_0=m_0,n_1=2m_1,n_2=2m_2$, on $N$ one has
$Fn_0=n_1$, $Fn_1=n_2$,
$En_1=2n_0$, $En_2=2n_1$.
The $H$-weights in both bases are $2,0,-2$.
Consequently $F^2:M_2\to M_{-2}$ multiplies by $2$,
whereas $F^2:N_2\to N_{-2}$ is an isomorphism.
Here subscripts denote $H$-weight, not the basis index.
For any positive integer $q$, the endomorphism
$\Phi_qm_j=q^jm_j$ preserves both lattices, satisfies
$\Phi_qF\Phi_q^{-1}=qF$, and has the same eigenvalues
$1,q,q^2$ in both. The same holds for all powers and for
restrictions to the graded pieces. Finally $N$ is not stable
under $F^{(2)}=F^2/2!$, since $F^{(2)}n_0=n_2/2$.
Thus even integral Lie operators do not automatically supply
an integral divided-power structure.
\end{proof}

This is an abstract integral-lattice counterexample, not a claim
that either lattice has been realized as the cohomology of a
specified variety. It establishes why rational parameter
recognition and Euler factors cannot determine the theta-series
index or the integral correspondence ring sought in Route 4.

\begin{proposition}[AG-LEFSCH-SPLIT: an exact sufficient lattice criterion]
Let $V$ be a finite rational $\mathfrak{sl}_2$ module with
\[
 V=\bigoplus_{m\ge0}
       \operatorname{Sym}^m(\mathbb Q^2)\otimes W_m.
\]
Suppose a lattice $M\subset V$ splits as
$\bigoplus_m M^{(m)}$ in these isotypic summands, each
$M^{(m)}$ splits into its $H$-weight lattices, and all
$E^{(a)}=E^a/a!$, $F^{(a)}=F^a/a!$ preserve each
isotypic lattice $M^{(m)}$.
For $r\ge0$ write $M_{\pm r}$ for the weight lattices.
Then
\begin{equation}\label{ag:integral-lefschetz-test}
 \operatorname{coker}
  (E^{(r)}:M_{-r}\longrightarrow M_r)
 \simeq
 \bigoplus_{\substack{m\ge r\\m\equiv r\ (2)}}
 \left(\mathbb Z/
     \binom{(m+r)/2}{r}\mathbb Z\right)^{\dim W_m}.
\end{equation}
For $E^r$ the cyclic orders are multiplied by $r!$.
Thus the inverse of $E^{(r)}$ has denominator dividing the
least common multiple of the displayed binomial coefficients.
\end{proposition}
\begin{proof}
The finite sums $\exp(E)=\sum_aE^{(a)}$ and
$\exp(-F)=\sum_a(-1)^aF^{(a)}$ preserve the lattice,
as do their inverses. Hence
$w=\exp(E)\exp(-F)\exp(E)$ is an integral automorphism.
On the standard two-dimensional module it sends $x$ to $-y$
and $y$ to $x$. It therefore carries weight $-r$ to weight
$r$ isomorphically on each $M^{(m)}$.
For $j=(m+r)/2$ the weight-$-r$ subspace of
$\operatorname{Sym}^m\otimes W_m$ is
$x^{m-j}y^j\otimes W_m$, and
\[
 E^{(r)}(x^{m-j}y^j)
   =\binom{j}{r}x^jy^{m-j},\qquad
 w(x^{m-j}y^j)=(-1)^{m-j}x^jy^{m-j}.
\]
Thus $E^{(r)}$ on this subspace is the displayed binomial
coefficient times an integral isomorphism, up to sign.
Its cokernel is the corresponding cyclic group repeated
$\dim W_m$ times. The assumed integral splitting permits the
direct sum over $m$. Multiplication by $r!$ gives the ordinary
power statement. The denominator assertion follows by applying
the integer annihilating this finite cokernel.
\end{proof}

For an extreme string ($m=r$), the divided power is unimodular
under these hypotheses, while the ordinary power contributes
$r!$. For intermediate weights the binomial factor can remain
nontrivial. This agrees with the inherited DP15 obstruction:
in dimension $3$, exterior divided powers are integral, but
the $H^2\to H^4$ paired block $J-I$ has cokernel
$\mathbb Z/2$. The general H8/Wilson diagonal result and
H8(iv)'s separate order/index claim retain their reported
status; they are not re-proved or silently identified with an
automorphic lattice here.

\begin{proposition}[AG-LEFSCH-INDEX: the index does not require integral splitting]
Keep the rational module $V=\bigoplus_m\operatorname{Sym}^m
(\mathbb Q^2)\otimes W_m$ and the operators above. Let
$M\subset V$ be any full lattice and put $M_j=M\cap V_j$,
without assuming that $M$ splits into its weight or isotypic
lattices. Fix $r\ge0$, and suppose that
\[
 f=E^{(r)}\text{ sends }M_{-r}\text{ into }M_r,
 \qquad w(M_{-r})=M_r,
 \quad w=\exp(E)\exp(-F)\exp(E).
\]
These two conditions hold, in particular, when all divided
powers of $E$ and $F$ preserve $M$.
Write
\[
 \begin{gathered}
 \mathcal D_r=\{m\ge r:m\equiv r\pmod2,\ W_m\ne0\},\\
 a_m=\binom{(m+r)/2}{r},\qquad d_m=\dim W_m,
 \qquad h_r=\sum_{m\in\mathcal D_r}d_m.
 \end{gathered}
\]
Then
\begin{equation}\label{eq:ag-lefschetz-index}
 [M_r:fM_{-r}]=\prod_{m\in\mathcal D_r}a_m^{d_m}.
\end{equation}
For the ordinary power $E^r$ the index is multiplied by
$(r!)^{h_r}$. Moreover, the exponent of $\operatorname{coker}f$,
equivalently the least denominator of $f^{-1}$ between these
lattices, divides $\prod_{m\in\mathcal D_r}a_m$.
Empty products are $1$.
\end{proposition}
\begin{proof}
The intersection of a full lattice with a rational subspace
is a full lattice in that subspace. Thus the two weight lattices
have rank $h_r$. The stated Weyl condition makes
$A=w^{-1}f$ an integral endomorphism of $M_{-r}$.
The monomial calculation in AG-LEFSCH-SPLIT gives its action on
the weight-$-r$ part of the $m$th rational isotypic summand:
it is multiplication by the nonzero integer
\[
 b_m=(-1)^{(m-r)/2}a_m.
\]
Consequently $f$ is injective with finite cokernel, and the
absolute determinant of $A$ is the right side of
\eqref{eq:ag-lefschetz-index}. The determinant-index formula
for an injective endomorphism of a free abelian group proves
the assertion, because $w$ is a lattice isomorphism.
Replacing $f$ by $r!f$ proves the ordinary-power formula.

For the exponent claim, the integral polynomial
$P(T)=\prod_{m\in\mathcal D_r}(T-b_m)$ annihilates $A$:
this is checked on each rational isotypic summand, where $A$
is scalar. Write $P(T)=P(0)+TQ(T)$ with $Q\in\mathbb Z[T]$.
Then
$P(0)M_{-r}=-A Q(A)M_{-r}\subseteq A M_{-r}$.
The absolute value of $P(0)$ is $\prod_m a_m$, proving the
annihilator and denominator statements. If the weight space
is zero the claims hold directly.
\end{proof}

\begin{proposition}[AG-LEFSCH-DEFECT: exact gluing and its possible primes]
Under the hypotheses of AG-LEFSCH-INDEX, put
\[
 N_{\pm r}=\bigoplus_{m\in\mathcal D_r}
       (M_{\pm r}\cap V^{(m)}),\qquad
 Q_{\pm r}=M_{\pm r}/N_{\pm r},
\]
where $V^{(m)}=\operatorname{Sym}^m(\mathbb Q^2)\otimes W_m$.
The groups $Q_{\pm r}$ are finite. Let
$\bar f:Q_{-r}\to Q_r$ be the induced map. There is an exact
sequence of finite groups
\begin{equation}\label{eq:ag-lefschetz-gluing}
 0\longrightarrow\ker\bar f
 \longrightarrow\bigoplus_{m\in\mathcal D_r}
        (\mathbb Z/a_m\mathbb Z)^{d_m}
 \longrightarrow\operatorname{coker}f
 \longrightarrow\operatorname{coker}\bar f
 \longrightarrow0.
\end{equation}
The first and last groups have equal orders. If a prime $\ell$
divides at most one of the integers $a_m$, the $\ell$-primary
cokernel is already forced: it is zero if no $a_m$ is divisible
by $\ell$, and otherwise it is
\[
 (\mathbb Z/\ell^{v_\ell(a_m)}\mathbb Z)^{d_m}
\]
for that unique $m$. Thus ambiguity beyond the split cyclic
formula can occur only at primes dividing two of its binomial
factors.
\end{proposition}
\begin{proof}
Each isotypic intersection is a full lattice in its rational
weight subspace, so their direct sum has finite index in the
whole weight lattice. The Weyl element preserves the rational
isotypic summands and identifies $N_{-r}$ with $N_r$, as well
as $M_{-r}$ with $M_r$. The operator $f$ acts on the $m$th
intersection as $b_m w$. It follows that it preserves these
two exact lattice sequences and has the split cokernel on
$N_{-r}\to N_r$ displayed in
\eqref{eq:ag-lefschetz-gluing}. The snake lemma, using
injectivity on both $N_{-r}$ and $M_{-r}$, gives that exact
sequence. Since $w$ identifies $Q_{-r}$ and $Q_r$, they have
the same finite order, so $\ker\bar f$ and
$\operatorname{coker}\bar f$ have the same order as well.

For the primary assertion use $A=w^{-1}f$ on
$M_{-r}\otimes\mathbb Z_\ell$. If all $b_m$ are units,
its determinant is a unit and its cokernel is zero.
If only $b_{m_0}$ is a nonunit, then every
$b_{m_0}-b_m$ for $m\ne m_0$ is a unit in $\mathbb Z_\ell$.
The polynomial projector
\[
 P_{m_0}(A)=\prod_{\substack{m\in\mathcal D_r\\m\ne m_0}}
       \frac{A-b_m}{b_{m_0}-b_m}
\]
therefore preserves this lattice. On the rational weight
space it projects onto $V^{(m_0)}_{-r}$, so it splits the
lattice into that intersection and a complementary summand.
On the first summand $A=b_{m_0}$ and the rank is $d_{m_0}$;
on the complement the determinant is a unit, so $A$ is an
integral automorphism. This gives exactly the asserted
$\ell$-primary group. For $r=0$ all $a_m=1$ and the first
case applies. This also avoids any issue with repeated unit
eigenvalues in that case.
\end{proof}

\paragraph{Sharpness of the denominator under the stated weaker hypotheses.}
Take $V=\operatorname{Sym}^3(\mathbb Q^2)\oplus
\operatorname{Sym}^7(\mathbb Q^2)$ and $r=1$.
Let $u=xy^2$ and $v=x^3y^4$ be the two weight-$-1$ monomials,
one in each summand. Then $A=w^{-1}E$ acts by $-2$ on $u$
and $-4$ on $v$. Choose
\[
 M_{-1}=\mathbb Zu+\mathbb Z\frac{u+v}{2},\qquad
 M_1=wM_{-1},
\]
and complete these to a full lattice using monomial lattices
in all other weights. In the displayed basis of $M_{-1}$,
\[
 A=\begin{pmatrix}-2&1\\0&-4\end{pmatrix}.
\]
Thus the required raising map is integral and the required
Weyl map is a lattice isomorphism. The cokernel is
$\mathbb Z/8\mathbb Z$, because the matrix has determinant
$8$ and its entries have greatest common divisor $1$.
With the split lattice $\mathbb Zu\oplus\mathbb Zv$ and its
Weyl image, the same rational map instead has cokernel
$\mathbb Z/2\oplus\mathbb Z/4$. Both indices are $8$,
but the least inverse denominator changes from $4$ to $8$.
This proves that the product annihilator in AG-LEFSCH-INDEX
cannot always be replaced by the least common multiple under
its stated hypotheses. This example asserts the local
raising/Weyl lattice conditions only; it supplies no claim
of stability under all divided powers or of a geometric
realization.

The theta-to-Lefschetz gate for a total index is now smaller:
identify the actual integral theta/period lattice with the
chosen cohomological lattice, and verify that the normalized
raising map is integral and the Weyl (or correctly normalized
Fourier) operator identifies the two weight lattices.
Integral isotypic splitting is unnecessary for the index.
It matters for the full Smith group only through the explicit
finite defect sequence; even there, the possible exceptional
primes are restricted as above. This does not assert the
missing theta-lattice comparison or identify it with H8(iv).
The rational Arthur comparison and AG-ARTHUR-MIN supply
different parts of the construction, and A and M do not by
themselves provide these integral maps.

\subsection{The geometric input to the finite-field factoring interface}

\paragraph{AG-CYCLIC-DIV (recent source theorem, with its full geometric proof checked).}
The additional manuscript PF26, OpenAI,
\emph{Deterministic Polynomial Factorization over Prime Fields},
4 October 2026, Section 4, pp.~13--19, gives the following
geometric mechanism. Let $q\ne\operatorname{char}K$ be an odd
prime, let the finite field $K$ contain $\mu_q$, and let
$F\in K[X]$ be monic, square-free and split, of degree
$N\ge q$ divisible by $q$. On the smooth projective curve
$C:Y^q=F(X)$, the deck transformation $\sigma(Y)=\zeta_qY$
acts on its Jacobian. Put $\lambda=1-\sigma$, let $\Lambda$
be the degree-zero divisor lattice on the $q$ points at infinity,
and write $e=(\sigma-1)\infty_0$.
For $r\ge0$ and $m=1+(q-1)r$, every solution
\[
 \lambda^t d=[W],\qquad d\in J(\bar K),\quad
 W\in\Lambda,\quad 1\le t\le m,
\]
is rational over the extension of $K$ of degree $q^r$, and
has a rational divisor representative there. This is exactly
Lemma 4.3 and Lemma 4.4; their definitions and complete proofs,
including Lemmas 4.1--4.2, were independently read and checked.
For $r=v_q(N)$ the extension degree is at most $N$.
Corollary 4.5 supplies the required ramification-class chains,
and Proposition 4.6 proves that the associated norm labels
must separate two ramification points by step $m$.

The decisive information is integral: $\Lambda$ is a free
rank-one $\mathbb Z[\zeta_q]$ module and
$Ne\in\lambda^{m-1}\Lambda\setminus\lambda^m\Lambda$.
The source retains a Jacobian class together with its actual
infinity divisor. Its filtration proof bounds the field of
definition, while its final lattice obstruction forces
separation for every sequence of permitted choices. This
connects a finite \'etale algebra's idempotent problem to
Jacobians and rational divisors, as used by CX-ALG-027/028.
Here $\lambda$ is the cyclotomic operator $1-\sigma$, not
a $\lambda$-ring operation. The result supplies a specific
integral geometric mechanism; it does not identify an
automorphic theta lattice or an arbitrary inertia extension.
The complete factoring transfers retain PF26's algebraic
algorithm as an additional recent-manuscript input. The
review recorded here covers its entire Section 4 only, not
the norm-solving implementation or a formal certificate.

\subsection{Central Artin orders and an abstract cohomological functor}

\paragraph{AG-DENINGER: a classical existence theorem with an extra analytic comparison.}
Deninger, \emph{Is there a Birch and Swinnerton-Dyer conjecture for
Dedekind zeta functions?}, arXiv:2504.15767v3 (1 April 2026),
Theorems 2.15--2.16, pp.~10--14, constructs an abstract functor
on number fields with symplectic forms, positive compatible star
operators and Galois descent.  On a Galois field its representation
is the sum, once each, of the irreducible symplectic Artin
representations with root number $-1$.  All such functors are
isomorphic, with automorphism group
$\operatorname{Map}(\Pi_-,\{\pm1\})$, where $\Pi_-$ denotes
these finite-image representations of $G_{\mathbb Q}$.
This abstract theorem is unconditional.  The additional
identification
$\dim V_K=\operatorname{ord}_{s=1/2}\zeta_K(s)$ uses the
Serre hypothesis stated on pp.~1--2: the central Artin order
is one for precisely those symplectic representations with
root number $-1$, and zero for every other irreducible.
Neither A nor M supplies that multiplicity assertion.

The following extension retains arbitrary central multiplicities.
Its proof also supplies the invariant-space argument for fields
which are not Galois over $\mathbb Q$.  In all statements a
positive star structure means a nondegenerate complex alternating
form $B$ and an antilinear, equivariant operator $J$ such that
$J^2=-1$ and $h(v,w)=B(v,Jw)$ is positive definite hermitian,
linear in $v$.  Equivariance of $J$ is essential.

\begin{proposition}[AG-CENTRAL-STRUCT: the exact finite positive-star criterion]
Let $G$ be finite and let
$\Theta=\sum_{\pi\in\operatorname{Irr}(G)}m_\pi\chi_\pi$,
with $m_\pi\in\mathbb Z$.  Suppose
$m_{\bar\pi}=m_\pi$, and that $m_\pi$ is even whenever
$\pi$ is orthogonal.  The following are equivalent:
\begin{enumerate}
\item Every $m_\pi$ is nonnegative.
\item The matrix $(\Theta(gh^{-1}))_{g,h\in G}$ is positive
semidefinite.
\item There is a complex $G$-representation of character $\Theta$
with a positive star structure.
\end{enumerate}
When these conditions hold, the structured representation is
unique up to equivariant isomorphism.  Its group of equivariant
automorphisms preserving $B$ and $J$ is
\[
 \prod_{\pi\ {
m symplectic}}\mathrm O(m_\pi)
 \ \times\!
 \prod_{\pi\ {
m orthogonal}}\mathrm{Sp}(m_\pi/2)
 \ \times\!
 \prod_{\{\pi,\bar\pi\},\,\pi\not\simeq\bar\pi}
                              \mathrm U(m_\pi).          \tag{AG-central-groups}
\]
Here $\mathrm O(m)$ is the real orthogonal group and
$\mathrm{Sp}(r)$ is the compact group on a quaternionic space
of dimension $r$; factors of size zero are omitted.
\end{proposition}

\begin{proof}
The equivalence of the first two conditions is the finite
convolution calculation of MT-ART02: on the matrix coefficients
of $\pi$, the normalized convolution matrix has eigenvalue
$m_\pi/\dim\pi$.  The third condition implies the first because
multiplicities of an actual representation are nonnegative.
It remains to construct and classify the positive star structures.

Fix invariant positive hermitian forms on irreducible
representations, obtained by averaging.  An orthogonal irreducible
has an antiunitary equivariant real structure $C_\pi$ with
$C_\pi^2=1$; a symplectic irreducible has an antiunitary
equivariant structure $J_\pi$ with $J_\pi^2=-1$.
For example, in the symplectic case define $A$ by
$h(v,w)=B_\pi(v,Aw)$ for a chosen invariant alternating form.
It is antilinear and equivariant, and $A^2=\lambda I$ by
Schur's lemma.  For $v\ne0$,
\[
 0<h(Av,Av)=B_\pi(Av,A^2v)
             =-\lambda B_\pi(v,Av)=-\lambda h(v,v).
\]
Thus $\lambda$ is negative real.  Moreover
$h(Av,Aw)=-\lambda h(w,v)$, so
$J_\pi=(-\lambda)^{-1/2}A$ is antiunitary and has square
$-1$.  For a symmetric invariant form the same calculation
has a plus sign: $\lambda>0$, and
$C_\pi=\lambda^{-1/2}A$ is antiunitary with square $1$.

For orthogonal $\pi$, equip its multiplicity space
$M_\pi=\mathbb C^{m_\pi}$ with its standard antiunitary
quaternionic structure, which exists because $m_\pi$ is even,
and use $C_\pi\otimes J_{M_\pi}$.  For symplectic $\pi$,
equip $M_\pi$ with ordinary conjugation and use
$J_\pi\otimes C_{M_\pi}$.  Tensor products here are antilinear
on both factors and therefore define a well-defined antilinear
operator on the complex tensor product.  Both constructions
square to $-1$ and are antiunitary for the product hermitian form.

For a non-selfdual pair set $W=\pi\otimes\mathbb C^{m_\pi}$
and use $W\oplus\bar W$ with its orthogonal sum hermitian form.
The map
\[
             J(w,\bar v)=(-v,\bar w)
\]
is equivariant, antiunitary and has square $-1$.
Take the orthogonal direct sum of these constructions.
If $h$ is its hermitian form, put
$B(v,w)=-h(v,Jw)$.  This is complex bilinear and nondegenerate.
Antiunitarity and $J^2=-1$ give
$h(w,Jv)=-h(v,Jw)$, so $B$ is alternating, and
$B(v,Jw)=h(v,w)$.  This proves existence.

For classification, Schur's lemma makes the hermitian form on
each isotypic component a tensor product with a positive form
on the multiplicity space.  Equivariance makes $J$ exchange
dual isotypic components.  On a selfdual component it therefore
has the form just used, with an antiunitary real structure on
the multiplicity space of a symplectic irreducible and an
antiunitary quaternionic structure on that of an orthogonal
irreducible.  These structures have orthonormal real or
quaternionic bases.  In the latter case one chooses a unit
vector $e$, then $Je$; they are orthogonal, and their
orthogonal complement is $J$-stable, giving induction.
For a non-selfdual pair choose an orthonormal basis on one
multiplicity space and let $J$ determine the other.
This proves uniqueness up to isomorphism.
Preserving $B$ and $J$ is equivalent to being unitary and
commuting with $J$.  On the respective multiplicity spaces
these are precisely $\mathrm O(m)$, $\mathrm{Sp}(m/2)$ and
$\mathrm U(m)$, proving (AG-central-groups).
\end{proof}

\begin{proposition}[AG-CENTRAL-FUNCTOR: the remaining central gate is exact]
For every irreducible finite-image complex representation $\pi$
of $G_{\mathbb Q}$ put
\[
             m_\pi=\operatorname{ord}_{s=1/2}L(s,\pi).
\]
The following central, rather than global, holomorphy condition
is equivalent to existence of an abstract positive-star functor
with the specified central Artin characters:
\[
 \text{every finite-image Artin function is holomorphic at }1/2.
                                                        \tag{AG-Central-H}
\]
More precisely, under (AG-Central-H) there is a functor $K\mapsto V_K$
on number fields and embeddings, with isometric pullbacks $i^*$,
adjoints $(i^*)^\dagger$, and
\[
 (i^*)^\dagger i^*=1,\qquad
 i^*(i^*)^\dagger=\frac1{[L:K]}
                 \sum_{\sigma\in\operatorname{Gal}(L/K)}\sigma
 \quad(L/K\text{ Galois}).                              \tag{AG-central-descent}
\]
For $L/\mathbb Q$ Galois its character is
$\sum_\pi m_\pi\chi_\pi$, with $\pi$ ranging over
$\operatorname{Irr}\operatorname{Gal}(L/\mathbb Q)$; for every
number field $K$ one has
\[
             \dim V_K=\operatorname{ord}_{s=1/2}\zeta_K(s).
                                                        \tag{AG-central-dimension}
\]
Such functors with those characters and descent are unique up
to isomorphism.  Their common automorphism group is the product
in (AG-central-groups) over all finite-image irreducibles of
$G_{\mathbb Q}$.  Conversely the prescribed character at every
Galois field forces (AG-Central-H).  An equality of dimensions
alone is not claimed to force this condition.
\end{proposition}

\begin{proof}
Brauer meromorphy makes the orders $m_\pi$ integers.
Complex conjugation of the Artin Dirichlet coefficients gives
$m_{\bar\pi}=m_\pi$.  At $s=1/2$ all archimedean factors are
finite and nonzero.  For a selfdual $\pi$ its functional
equation therefore gives
\[
                        (-1)^{m_\pi}=W(\pi),             \tag{AG-central-parity}
\]
even if the order is negative.  The Fröhlich--Queyrut theorem
gives $W(\pi)=1$ for orthogonal $\pi$.  An exact primary
statement is Queyrut, \emph{Sur l'équation fonctionnelle des
séries $L$ d'Artin pour des caractères réels}, Mémoires SMF 37
(1974), theorem on p.~133: it applies to representations
realizable over $\mathbb R$, not to all real-valued characters.
Thus the parity hypotheses of AG-CENTRAL-STRUCT hold.
Condition (AG-Central-H) gives the remaining nonnegativity.

Choose one representative of each irreducible and form the
algebraic direct sum
\[
 \mathcal V=\bigoplus_\pi \pi^{\oplus m_\pi},\qquad
                V_K=\mathcal V^{G_K},
 \quad G_K=\operatorname{Gal}(\overline{\mathbb Q}/K),  \tag{AG-central-construction}
\]
equipped with the blockwise forms and star operators of
AG-CENTRAL-STRUCT.  Every vector has finite orbit.  For a
finite Galois extension $L/\mathbb Q$, its invariant subspace
contains only the finitely many irreducibles factoring through
$\operatorname{Gal}(L/\mathbb Q)$, hence is finite dimensional.
For arbitrary $K$ embed it in such an $L$; then $V_K\subset V_L$
is also finite dimensional.  The restriction of $h$ to $V_K$
is positive, $J$ preserves $V_K$, and the restriction of $B$
is nondegenerate because $B(v,Jv)=h(v,v)>0$ for $v\ne0$.

An embedding $\alpha:K\hookrightarrow L$ extends to
$\sigma\in G_{\mathbb Q}$; use the restriction of $\sigma$
to $V_K$.  Two extensions differ on the right by an element
of $G_K$, so this map is independent of the choice.  It is
isometric and respects $B,J$, and composition is functorial.
For an inclusion $K\subset L$ with $L/K$ Galois, the average
of $\operatorname{Gal}(L/K)$ on $V_L$ is the orthogonal
projector to $V_L^{\operatorname{Gal}(L/K)}=V_K$.
This proves (AG-central-descent), including the case that
$K/\mathbb Q$ is not Galois.  Unnormalized transfer is
$[L:K](i^*)^\dagger$, so its composite with pullback is
multiplication by the extension degree.

Choose a finite Galois $L/\mathbb Q$ containing $K$ and put
$H=\operatorname{Gal}(L/K)$.  Artin formalism for the
permutation representation $\operatorname{Ind}_H^G1$ gives
\[
 \operatorname{ord}_{1/2}\zeta_K
       =\sum_{\pi\in\operatorname{Irr}(G)}
                    m_\pi\dim\pi^H
       =\dim V_K.
\]
This proves the dimension formula.  Conversely an actual
representation with the prescribed character has nonnegative
irreducible multiplicities, which is precisely central
holomorphy.  Representations over other base number fields
are covered as well by induction to $G_{\mathbb Q}$.

Finally consider any functor with the stated character and
descent properties.  Take its directed union over finite
Galois extensions of $\mathbb Q$; isometric pullbacks are
injective.  This union carries a locally finite $G_{\mathbb Q}$
action, and the averaging identity recovers its $G_K$-invariants
as $V_K$.  On each finite quotient the prescribed character
gives exactly the multiplicities $m_\pi$.  The union therefore
decomposes into the blocks of (AG-central-construction).
The classification in AG-CENTRAL-STRUCT gives a blockwise
equivariant isometry to that construction and hence a natural
isomorphism of functors.  Natural automorphisms are precisely
the compatible equivariant block automorphisms, which gives
the asserted product of compact classical groups.
\end{proof}

Under Serre's extra multiplicity conjecture all the surviving
factors in (AG-central-groups) are $\mathrm O(1)=\{\pm1\}$,
recovering Deninger's result.  With higher multiplicities the
same central positivity condition allows larger automorphism
groups.  MT-ART02 identifies the exact finite positive matrix
which must hold at each Galois level.  M confines the possible
Artin divisor to the critical line, but it does not supply
(AG-Central-H) at its centre.  These are abstract representation
spaces attached to arithmetic curves $\operatorname{Spec}\mathcal O_K$;
no geometric cohomology, integral lattice, cycle map, regulator
or formula for the leading coefficient has been constructed.

\begin{proposition}[AG-CENTRAL-REAL: the real-coefficient obstruction survives higher order]
For a Galois field $L/\mathbb Q$ satisfying central Artin
holomorphy, its representation $V_L$ in AG-CENTRAL-FUNCTOR
has an equivariant real form if and only if every symplectic
irreducible has even central multiplicity.  In particular a
symplectic irreducible with root number $-1$ forbids such a
real form, regardless of whether its central order is one,
three or larger.  The same obstruction applies unconditionally
to Deninger's sign-selected functor whenever that negative-root
representation occurs.
\end{proposition}
\begin{proof}
A real form is equivalent to an equivariant antilinear operator
$C$ with $C^2=1$.  On a symplectic isotypic component write
$C=J_\pi\otimes A$ with $J_\pi^2=-1$ as above.  Then $A$
is antilinear on its multiplicity space and must satisfy
$A^2=-1$.  If $A=M\circ\mathrm{conj}$ on $\mathbb C^m$,
then $M\bar M=-I$, so
$|\det M|^2=(-1)^m$; this requires $m$ even.
Conversely the standard quaternionic $A$ exists in every even
dimension.  Orthogonal blocks always have a real form, and
nonselfdual pairs have one by exchanging the two conjugate
summands.  This proves the equivalence.  Formula
(AG-central-parity) makes $m_\pi$ odd when $W(\pi)=-1$.
For the sign-selected functor its multiplicity is already one.
\end{proof}

Thus the root-number class studied in MT-ART10--12 can obstruct
a functorial real or rational coefficient realization, as well
as force central vanishing.  This statement concerns the
equivariant coefficient structure; it does not rule out the
positive complex quaternionic realization just constructed.

\begin{proposition}[AG-CENTRAL-Q8: an unconditional arithmetic real-form obstruction]
For the explicit tame quaternionic field
\[
 K=\mathbb Q(\sqrt5,\sqrt{41}),\qquad
 N=K\left(\sqrt{\frac{5+\sqrt5}{2}
                         \frac{41+\sqrt{205}}{2}}\right),
\]
the actual central-order character
$\sum_\pi\operatorname{ord}_{s=1/2}L(s,\pi)\,\pi$ is the
character of a complex representation of $\operatorname{Gal}(N/\mathbb Q)$,
but that representation has no equivariant real form.
This assertion needs neither M, general central Artin holomorphy,
nor Serre's conjecture on the exact central multiplicities.
\end{proposition}
\begin{proof}
MT-ART10 records the classical Martinet example (1971,
Theorem IV.1, p.~406, Example 3, p.~407), and Fr\"ohlich's
root-number criterion: its group is $Q_8$ and the unique
degree-two symplectic character $\chi$ has $W(\chi)=-1$.
Every irreducible is either linear or this $\chi$.
The linear factors are Dirichlet functions (including $\zeta$),
all holomorphic at $1/2$.  If $C_4\subset Q_8$ and $\lambda$
is a faithful linear character, then
$\chi=\operatorname{Ind}_{C_4}^{Q_8}\lambda$; hence its
Artin function is an entire nontrivial Hecke function.
Thus all central orders are nonnegative integers, so their
sum is an actual character.  The functional equation with
$W(\chi)=-1$ makes its multiplicity positive and odd.
AG-CENTRAL-REAL now forbids a real form.  The field and its
root sign are the located classical input from MT-ART10;
this deduction does not claim a new computation of their values.
\end{proof}

This is a concrete obstruction for a proposed real-coefficient
central cohomology of the arithmetic curve
$\operatorname{Spec}\mathcal O_N$, despite positivity of the
complex central character.  The complex symplectic/star
structure of AG-CENTRAL-STRUCT remains available.

\subsection{A precise obstruction to the exponential half-Tate candidate}

\paragraph{AG-EXP-VANISH: the located classical obstruction.}
Fresán--Jossen's theorem, reproduced with its complete proof as
Deninger's Theorem 4.1, pp.~16--17, states
\[
 \operatorname{Ext}^1_{\mathcal M^{\exp}(K)}
             (\mathbb Q(0),\mathbb Q(1/2))=0
                       \qquad(K\subset\mathbb C).
\]
Here $\mathbb Q(1/2)$ is the dual of the exponential motive
of $(\mathbb A^1,x^2)$, in the specified exponential-motive
category.  This is a proved obstruction to that particular
candidate for the central cohomology, not a vanishing theorem
for all proposed arithmetic cohomology theories.
The following formulation isolates the mechanism and extends it
to any monodromy with no eigenvalue one.

\begin{proposition}[AG-EXP-SEPARATE: monodromy separation kills the extension]
Let $N\in\mathcal M^{\exp}(K)$ have perverse realization
$j_!\mathcal L[1]$ on $\mathbb A^1_{\mathbb C}$, where
$j:\mathbb G_m\hookrightarrow\mathbb A^1$, and let $T$
be the rational monodromy of the finite-rank local system
$\mathcal L$ about zero.  If $T-I$ is invertible, then
\[
 \operatorname{Ext}^1_{\mathcal M^{\exp}(K)}
                         (\mathbb Q(0),N)=0.
\]
In particular $T=-1$ recovers AG-EXP-VANISH.
\end{proposition}
\begin{proof}
Take an extension $0\to N\to M\to\mathbb Q(0)\to0$.
The perverse realization of $\mathbb Q(0)$ is
$j_!\mathbb Q[1]$.  Exactness of realization gives an
extension of these two perverse sheaves.  Its stalk complex
at zero vanishes, since those of both end terms vanish.
On $\mathbb G_m$ it is an extension of local systems, so in
a basis adapted to the extension its monodromy is
\[
                   \begin{pmatrix}T&b\\0&1\end{pmatrix}.
\]
Changing the lift of the last basis vector by
$u=-(T-I)^{-1}b$ makes this matrix block diagonal.
The vanishing stalk at zero identifies the perverse sheaf
with the extension by zero of its restriction.  Hence the
perverse extension splits.  Its largest trivial subobject
is exactly $j_!\mathbb Q[1]$, because $T-I$ is invertible
on the other summand.

Apply the theorem of the fixed part for exponential motives,
Fresán--Jossen, \emph{Exponential motives}, Theorem 6.5.1 in
the version cited as [8] in Deninger.  The exact statement
quoted and used in his proof of Theorem 4.1 says that the
largest trivial perverse subobject is realized
by a classical submotive $M_0\subset M$.  The composite
$M_0\to M\to\mathbb Q(0)$ realizes an isomorphism.
The realization is faithful and exact, so its kernel and
cokernel are zero; the composite is an isomorphism of motives.
Its inverse supplies a section of $M\to\mathbb Q(0)$,
proving that the extension splits.
\end{proof}

The fixed-part theorem is an explicit classical input here;
its complete source proof is not re-established by the matrix
calculation.  The exact monodromy hypothesis also matters:
when $1$ is an eigenvalue, $(T-I)^{-1}$ is unavailable and
this argument says nothing about the extension group.
The obstruction narrows the central cohomology construction
suggested by the trace/Fredholm programme of cote 1.  It does
not evaluate the distinct integral Kummer extensions of cote
149 or the general special-value $K$-groups in AG-EXT-K.

\subsection{The exact local exponential and determinant hypotheses}

The local regulator step in a theta or special-value construction
has a condition distinct from the global zero-location inputs.
We fix $G_p=\operatorname{Gal}(\overline{\mathbb Q}_p/\mathbb Q_p)$,
a finite extension $L/\mathbb Q_p$, and a de Rham
$L$-representation $V$ of $G_p$.  Put
\[
 D=D_{\mathrm{cris}}(V),\qquad
 t_V=D_{\mathrm{dR}}(V)/\operatorname{Fil}^0D_{\mathrm{dR}}(V).
\]
Frobenius on $D$ is $L$-linear here.  We use the convention
that $\mathbb Q_p(1)$ has Hodge--Tate weight $+1$,
crystalline Frobenius $p^{-1}$ and filtration jump $-1$.
The classical inputs are Bloch--Kato, \emph{L-functions and
Tamagawa numbers of motives}, Proposition 3.8 and its full
proof, pp.~354--359, Corollary 3.8.4, pp.~355--356,
and Definition 3.10, p.~359.  In particular
\[
\begin{split}
0\longrightarrow H^0(G_p,V)&\longrightarrow D^{\varphi=1}
 \longrightarrow t_V\xrightarrow{\exp_V}H_e^1(G_p,V)
 \longrightarrow0,\\
0\longrightarrow H^0(G_p,V)&\longrightarrow D
 \xrightarrow{\ d\mapsto((1-\varphi)d,\bar d)\ }
 D\oplus t_V\longrightarrow H_f^1(G_p,V)\longrightarrow0.
\end{split}
\tag{AG-local-exact}
\]
These are statements about rational local cohomology.  An
integral lattice or its Tamagawa index is additional data.

\begin{proposition}[AG-LOCAL-EXP: the exact exponential gate]
The kernel of $\exp_V:t_V\to H^1(G_p,V)$ is
\[
 \operatorname{im}\bigl(D^{\varphi=1}\longrightarrow t_V\bigr),
 \qquad
 H_f^1(G_p,V)/H_e^1(G_p,V)
       \simeq D/(1-\varphi)D.
\tag{AG-local-gate}
\]
Consequently the exponential is injective precisely when that
image is zero, and it is an isomorphism
$t_V\xrightarrow{\sim}H_f^1(G_p,V)$ precisely when
$D^{\varphi=1}=0$.
If $V\simeq V^*(1)$ and $H^0(G_p,V)=0$, then the following
are equivalent:
\[
 D^{\varphi=1}=0,\quad \exp_V\text{ is injective},\quad
 H_e^1=H_f^1,\quad H_g^1=H_f^1.
\]
Under these equivalent conditions the dual exponential induces
an isomorphism
\[
 H^1(G_p,V)/H_f^1(G_p,V)\xrightarrow{\sim}
            \operatorname{Fil}^0D_{\mathrm{dR}}(V).
\]
\end{proposition}
\begin{proof}
The first equality in (AG-local-gate) is exactness of the
first row of (AG-local-exact).  In the second row, quotient
the target $D\oplus t_V$ by the image of its second summand.
The resulting cokernel is $D/(1-\varphi)D$, proving the
second equality.  Since $1-\varphi$ is an endomorphism of
one finite-dimensional $L$-space, its cokernel vanishes if
and only if its kernel does.  If that kernel vanishes, both
injectivity and surjectivity of $\exp_V$ follow; conversely
surjectivity onto $H_f^1$ already forces it to vanish.
When $H^0=0$, the map $D^{\varphi=1}\to t_V$ is injective,
giving the additional injectivity equivalence.  Bloch--Kato
duality states that $H_e^1(V)$ and $H_g^1(V^*(1))$ are exact
annihilators, and $H_f^1(V)$ is the exact annihilator of
$H_f^1(V^*(1))$.  Selfduality and nondegeneracy therefore
give the last equivalence.  Finally the dual exponential is
adjoint to $\exp_{V^*(1)}$ under the perfect de Rham and
Tate pairings.  Dualizing the isomorphism onto $H_f^1$ gives
the asserted quotient isomorphism, with its stated target.
\end{proof}

\begin{proposition}[AG-LOCAL-EXCEPTION: generic selfduality does not suffice]
For every odd prime $p$, there is an irreducible crystalline
two-dimensional $\mathbb Q_p$-representation $V$ such that
\[
 \begin{gathered}
 \det V=\mathbb Q_p(1),\qquad
 H^0(G_p,V)=H^2(G_p,V)=0,\\
 \dim H_f^1(G_p,V)=1,\qquad
 \exp_V=0,\qquad\exp_V^*=0.
 \end{gathered}
\]
Its Hodge--Tate weights are $(2,-1)$.
\end{proposition}
\begin{proof}
Take $D=\mathbb Q_pe_1\oplus\mathbb Q_pe_2$, monodromy zero,
and
\[
 \varphi e_1=e_1,\qquad \varphi e_2=p^{-1}e_2,
 \qquad
 \operatorname{Fil}^iD=
 \begin{cases}
 D&i\le-2,\\
 \mathbb Q_p(e_1+e_2)&-1\le i\le1,\\
 0&i\ge2.
 \end{cases}
\tag{AG-local-example}
\]
The total Hodge and Newton numbers are both $-1$.
The only proper nonzero Frobenius-stable subspaces are
$\mathbb Q_pe_1$ and $\mathbb Q_pe_2$.  Each has induced
Hodge number $-2$, while its Newton number is respectively
$0$ or $-1$.  Thus $t_H\le t_N$ for every such subspace,
and $D$ is weakly admissible.  Colmez--Fontaine,
\emph{Construction des représentations $p$-adiques
semi-stables}, Theorem A (author version, p.~2), supplies
the crystalline representation $V$.  No proper nonzero
Frobenius-stable subspace has equal Hodge and Newton numbers,
so no such subspace is admissible.  The equivalence of
categories therefore also proves that $V$ is irreducible.

Its determinant has Frobenius $p^{-1}$ and its unique
filtration jump at $-1$, so it is $\mathbb Q_p(1)$.
The wedge pairing gives the required symplectic selfduality.
Moreover
\[
 H^0(G_p,V)=D^{\varphi=1}\cap\operatorname{Fil}^0D
   =\mathbb Q_pe_1\cap\mathbb Q_p(e_1+e_2)=0.
\]
Local duality gives $H^2=0$.  The map
$D^{\varphi=1}=\mathbb Q_pe_1\to
t_V=D/\mathbb Q_p(e_1+e_2)$ is an isomorphism.
The first row of (AG-local-exact) gives $\exp_V=0$ and
$H_e^1=0$; the second row gives $\dim H_f^1=\dim t_V=1$.
The local Euler characteristic gives $\dim H^1=2$.
Finally, by symplectic selfduality,
$\exp_{V^*(1)}=0$ as well.  Adjointness of exponential and
dual exponential under perfect pairings forces
$\exp_V^*=0$ on all of $H^1$.  Equivalently,
Bloch--Kato duality gives $H_g^1=(H_e^1)^\perp=H^1$.
\end{proof}

\paragraph{Exact scope of the recent-manuscript correction.}
Burungale--Kobayashi--Nakamura--Ota,
\emph{A local sign decomposition for symplectic self-dual
Galois representations of rank two}, arXiv:2508.17776v1,
allows $R=\mathbb Q_p$ in \S2.2(ii), p.~12.  Its
Definition 2.11, p.~16, defines genericity in this case
by $H^0=H^2=0$.  Thus (AG-local-example) satisfies exactly
those hypotheses, and the identity is a de Rham specialization
in the sense of \S5.1.1, p.~35.  Corollary 3.2, p.~19,
asserts nonvanishing of the ordinary dual exponential on a
basis of one of the signed lines.  Here it vanishes on the
whole cohomology, so that corollary is false in its stated
generality.  There is no purity or exclusion of Frobenius
fixed vectors in those definitions.

Separately, the displayed proof of Proposition 4.17,
pp.~28--29, uses $\exp^{-1}:H_f^1\to t_V$ as an
isomorphism.  AG-LOCAL-EXP gives its exact missing hypothesis.
This identifies a gap in that comparison as written; it does
not refute the main sign-decomposition theorem.
AG-LOCAL-SPLICE below repairs this precise determinant
comparison by a polarization-preserving specialization of
the full exact complex. It does not repair the separate
ordinary-exponential nonvanishing assertion. Applications for which the
crystalline period space vanishes, or which satisfy the
purity condition below, are not counterexamples.

\begin{proposition}[AG-LOCAL-DET: the finite determinant survives the exception]
Suppose $H^0(G_p,V)=0$.  With the ordering $D\oplus t_V$
as written in (AG-local-exact), that exact sequence gives
a canonical isomorphism of ordinary determinant lines
\[
 \delta_f:\det_L H_f^1(G_p,V)\xrightarrow{\sim}\det_L t_V.
\]
If $1-\varphi$ is invertible and $u$ is a top exterior
vector of $t_V$, its relation with the ordinary exponential is
\[
 \delta_f\bigl((\det\exp_V)(u)\bigr)
                      =\det(1-\varphi\mid D)\,u.
\tag{AG-local-det}
\]
The determinant isomorphism remains defined when the ordinary
exponential is zero.
\end{proposition}
\begin{proof}
The second row of (AG-local-exact) is now a short exact
sequence $0\to D\to D\oplus t_V\to H_f^1\to0$.
Use its exterior-power isomorphism, with the subspace basis
preceding lifted quotient vectors, and cancel the identical
source and target $\det D$ factors.  This defines $\delta_f$
without dividing by a Frobenius eigenvalue.  When $F=1-\varphi$
is invertible, choose a basis of $D$ and lift the exponential
of a basis of $t_V$ by vectors $(0,u_i)$.  The resulting
block matrix has diagonal blocks $F$ and $I$, proving
(AG-local-det).

In (AG-local-example), put $u=\bar e_1=-\bar e_2$ and
$a=1-p^{-1}$.  The image of $e_1$ in $D\oplus t_V$ is
$(0,u)$ and that of $e_2$ is $(ae_2,-u)$.
The quotient is generated by $h=[(e_1,0)]$.
In the ordered target basis $(e_1,e_2,u)$, the columns
consisting of those two image vectors and the lift of $h$
have determinant $-a$.  Hence
$\delta_f(h)=-a u\ne0$, although $\exp_V(u)=0$.
This also verifies directly that the two maps cannot be
interchanged at an exceptional period.
\end{proof}

\begin{proposition}[AG-LOCAL-PURITY: a geometric condition that removes the exception]
Assume all eigenvalues of crystalline Frobenius on $D$ are
algebraic numbers of a fixed Weil weight $w\ne0$, meaning
that every complex embedding has absolute value $p^{w/2}$.
Then $D^{\varphi=1}=0$, so $\exp_V:t_V\to H_f^1$ is an
isomorphism.  In particular this holds for
\[
 V=H^{2n+1}_{\text{ét}}
       (X_{\overline{\mathbb Q}_p},\mathbb Q_p)(n+1),
\]
when $X/\mathbb Q_p$ admits a smooth projective model over
$\mathbb Z_p$.  Crystalline comparison and the already-known
finite-field purity theorem give
$w=(2n+1)-2(n+1)=-1$.
\end{proposition}
\begin{proof}
The algebraic number $1$ has absolute value $1$ in every
complex embedding.  Since $p^{w/2}\ne1$, it cannot be an
eigenvalue.  Apply AG-LOCAL-EXP.  In the geometric case,
comparison identifies Frobenius eigenvalues with those of
the smooth projective reduction in degree $2n+1$, and the
Tate twist multiplies them by $p^{-(n+1)}$, giving the
displayed weight.  Exact classical locators are
Bhatt--Morrow--Scholze, \emph{Integral $p$-adic Hodge theory},
Theorem 1.1(i), author version p.~2, repeated as
Theorem 14.6(i), pp.~121--122; and Katz--Messing,
\emph{Some consequences of the Riemann hypothesis for varieties
over finite fields}, Theorem 1, pp.~74--75, and
Corollary 1, pp.~75--76.  If the reduction has several
geometric components, apply the latter theorem after a
finite residue extension to each component; eigenvalues
of a Frobenius power then have the claimed absolute
values, hence so do the original eigenvalues.
These are classical inputs, as in AG-PURITY.  No new purity
theorem or consequence of A or M is asserted.
\end{proof}

For motivic realizations this identifies the exact local
condition needed before turning a de Rham tangent vector
or a normalized theta derivative into a finite local
cohomology class by an invertible exponential.  At bad
reduction, or for an abstract local representation such
as (AG-local-example), purity of that precise crystalline
space must be checked separately.  Neither global GRH
input supplies it.  Even where the rational exponential
is invertible, the integral image, period normalization
and global cycle map remain the separate gates already
recorded in AG-THETA-FOURIER and AG-SHIMURA-PERIOD.

\subsection{Repairing the full local determinant comparison}

We retain the normalization $\varphi=p^{-1}$ on $L(1)$.
The classical potentially semistable theorem supplies the
Weil--Deligne representation of a de Rham representation.
In the unnormalized convention used here,
$N\varphi=p\varphi N$ and
$D_{\mathrm{cris}}$ is the inertia-invariant kernel of $N$.
Exact primary locators are Berger, \emph{Équations
différentielles $p$-adiques et $(\varphi,N)$-modules filtrés},
\S I.1, p.~16, for the conventions, and \S V.1, p.~31,
for de Rham $\Longleftrightarrow$ potentially semistable;
Nakamura, \emph{A generalization of Kato's local
$\varepsilon$-conjecture}, author version arXiv:1305.0880v3,
\S3.3, p.~35, constructs the associated Weil--Deligne object.
These substantial classical existence inputs are not new
consequences of A or M.

\begin{proposition}[AG-LOCAL-WD: the exceptional rank-two case is crystalline]
Let $V$ be a two-dimensional de Rham $L$-representation
of $G_p$ with $\det V=L(1)$.  If
$D_{\mathrm{cris}}(V)^{\varphi=1}\ne0$, then $V$ is
crystalline.  More precisely, if $V$ is not crystalline,
its crystalline period space is either zero or a line
on which $\varphi$ acts by $p^{-1}$ or $-p^{-1}$.
Consequently the ordinary exponential of every such
noncrystalline $V$ is an isomorphism $t_V\to H_f^1$.
\end{proposition}
\begin{proof}
Extend the coefficient field to split the finite inertia
representation on the two-dimensional Weil--Deligne
space $W$.  This does not change the relevant dimensions
or the presence of the eigenvalue one.  Inertia is
semisimple in characteristic zero, and its determinant
is trivial because the Weil--Deligne determinant of
$L(1)$ is unramified.  Thus $\dim W^I$ cannot be one:
a complementary inertia character would also have to be
trivial.  The only possibilities are zero and two.

If $W^I=0$, its subspace $(\ker N)^I$ is zero, hence
$D_{\mathrm{cris}}=0$.  If $W^I=W$, the representation
is semistable.  When $N=0$ it is crystalline.  Otherwise
$\ker N=\operatorname{im}N$ is a line.  If Frobenius
acts on that line by $\lambda$, the relation
$N\varphi=p\varphi N$ makes its action on the quotient
$p\lambda$.  Its determinant is $p^{-1}$, so
\[
                  p\lambda^2=p^{-1},\qquad
                  \lambda=\pm p^{-1}.
\]
Neither value is one.  This proves both assertions
about the crystalline space, and AG-LOCAL-EXP proves
the exponential assertion.  The line and its scalar
descend to $L$, since these two scalars already lie in
$\mathbb Q_p$.
\end{proof}

We now use the actual de Rham epsilon determinant,
including its graded ordering.  Let $p$ be odd and let
$V$ be symplectic selfdual of rank two, de Rham and
$H^0(G_p,V)=0$.  Write
\[
 \begin{gathered}
 H=H^1(G_p,V),\quad U=H_f^1(G_p,V),\quad E=D_{\mathrm{dR}}(V),\\
 F^0=\operatorname{Fil}^0E,\qquad t=E/F^0.
 \end{gathered}
\]
The symmetric Tate pairing $q$ on $H$ is perfect, and
$U$ is a one-dimensional Lagrangian.  The alternating
polarization $\beta$ on $E$ takes values in
$D_{\mathrm{dR}}(L(1))$.  Fix a basis $\zeta$ of
$L(1)$ and its Fontaine period $\tau$; the corresponding
de Rham basis is $\zeta/\tau$.  We identify the second
fundamental line with $L\zeta$ by the polarization.
Let $B_\epsilon:H\times H\to L\zeta$ be the alternating
pairing induced by the de Rham epsilon isomorphism.
Equivalently, define $w$ by
\[
                  B_\epsilon(x,y)=q(wx,y)\zeta.
\]
We use Nakamura's Weil--Deligne epsilon convention,
not the epsilon constant of the underlying Weil
representation alone.  His Remark 3.6, p.~38, records
their possible difference by
$\det(-\varphi\mid D_{\mathrm{st}}/D_{\mathrm{cris}})$.

\begin{proposition}[AG-LOCAL-SPLICE: the exceptional epsilon comparison survives]
With these conventions, even when
$D_{\mathrm{cris}}(V)^{\varphi=1}\ne0$, one has
\[
 B_\epsilon(x,y)
   =-\Gamma(V)\epsilon(W(V))q(x,y)\zeta
       \qquad(x\in U,\ y\in H).
\tag{AG-local-splice}
\]
In this exceptional case $V$ is crystalline and
$\epsilon(W(V))=1$.  Thus $U$ is exactly the
$-\Gamma(V)\epsilon(W(V))$ eigenspace of $w$.
Together with the inverse-exponential calculation on
the locus $D_{\mathrm{cris}}^{\varphi=1}=0$, this repairs
the stated conclusion of BKNO Proposition 4.17 without
restoring the false nonvanishing assertion in its
Corollary 3.2.
\end{proposition}
\begin{proof}
We first specify the determinant map that must be
specialized.  Put $V^*=V^\vee(1)$ and
$D^*=D_{\mathrm{cris}}(V^*)$.  The finite exact sequence
for $V$ and the Tate dual of the one for $V^*$ splice
to give
\[
\begin{aligned}
0\longrightarrow D&\longrightarrow D\oplus t
 \longrightarrow H\\
 &\longrightarrow (D^*)^\vee\oplus F^0
 \longrightarrow (D^*)^\vee\longrightarrow0.
\end{aligned}
\tag{AG-local-spliced-complex}
\]
The first map is $d\mapsto((1-\varphi)d,\bar d)$.
The last map uses the transpose of $1-\varphi^*$,
where $\varphi^*$ is the Frobenius on the actual Tate
dual.  The other maps use the finite local subgroup,
its exact annihilator, and the perfect de Rham pairing.
This is exactly Nakamura (22)--(24), pp.~36--37.
Apply his acyclic determinant trivialization (19),
cancel the two copies of each crystalline determinant,
and use
\[
        \det F^0\otimes\det t\longrightarrow\det E,
                  \qquad v\otimes u\longmapsto v\wedge u.
\]
The result is the canonical map
$\Theta:\det H\to\det E$ occurring in that definition.
Multiplying by $\Gamma(V)$ and then applying
$\theta_{\mathrm{dR}}$ gives $B_\epsilon$;
the latter map sends $\zeta/\tau$ to
$\epsilon(W(V))\zeta$ by Lemma 3.4, pp.~37--38.
We will specialize precisely this map, rather than
an unrelated determinant of $1-\varphi$.

By AG-LOCAL-WD, the exceptional representation is
crystalline.  Its determinant is $p^{-1}$ and it has
an eigenvalue one, so its two distinct Frobenius
eigenvalues are $1,p^{-1}$.  Choose the corresponding
basis of $D=E$.  Over $R=L[[z]]$ set
\[
 \varphi_z=
 \begin{pmatrix}1+z&0\\0&p^{-1}(1+z)^{-1}\end{pmatrix},
       \qquad \varphi_z^*=p^{-1}\varphi_z^{-t}.
\tag{AG-local-polarized-deformation}
\]
Keep $F^0$, the entire filtration and $\beta$ fixed.
Every alternating form in dimension two is multiplied
by the determinant of a matrix.  Consequently
$\beta(\varphi_zx,\varphi_zy)=p^{-1}\beta(x,y)$.
The same fixed polarization identifies $\varphi_z$
with $\varphi_z^*$ for every $z$.  This checks the
actual Tate twist and dual Frobenius; the two operators
$1-\varphi_z$ and $1-\varphi_z^*$ are conjugate and
have identical determinants.

The map
\[
 a_z:D_R\longrightarrow D_R\oplus t_R,
       \qquad d\longmapsto((1-\varphi_z)d,\bar d)
\]
is a split injection with free rank-one cokernel $U_R$.
Indeed $a_0$ is injective by $H^0=0$, so one of its
maximal minors is nonzero in $L$ and hence a unit in
$R$.  The same statement holds for the actual dual
map, and the fixed polarization identifies its
cokernel with $U_R$.  Form
$H_R=U_R\oplus U_R^\vee$ with its hyperbolic symmetric
pairing.  At $z=0$ identify it with $(H,q)$ preserving
$U$: any complement can be made isotropic by replacing
a vector $c$ by
$c-q(c,c)h/(2q(h,c))$, where $h$ spans $U$.
Using the fixed de Rham pairing and this Tate pairing,
splice the finite sequence over $R$ with its actual
dual, in exactly the order of
(AG-local-spliced-complex).  The sequence is exact and
all its modules are finite free.  The determinant
construction therefore gives a regular isomorphism
\[
                   \Theta_R:\det H_R\xrightarrow{\sim}\det E_R
\]
which specializes to the original $\Theta$.  This uses
only finite free linear algebra and base change of
determinants; no family of Galois representations is
being asserted.

For completeness, compute that map over the fraction
field $K=L((z))$.  Both $A=1-\varphi_z$ and
$A^*=1-\varphi_z^*$ are invertible there.  The finite
cokernel has the identification
\[
 [(d,u)]\longmapsto u-\overline{A^{-1}d},
\]
which makes the formal exponential $t_K\to U_K$ an
isomorphism.  In the target basis $D_K\oplus t_K$,
the columns of $a_z$ followed by the lift $(0,u)$ of
an exponential basis have determinant $\det A$.
In the dual sequence, put a basis of its subspace
first and lifts of a basis of its quotient last.
Solving the quotient map uses $(A^{*t})^{-1}$.
Moving its rank-one subspace past the two crystalline
coordinates has sign $(-1)^2=1$, so its corresponding
determinant is $(\det A^*)^{-1}$.
Their ratio is one, by the conjugacy just proved.
Moreover each crystalline determinant has graded
degree two.  Moving or cancelling these factors
introduces no further Koszul sign.  After these
cancellations the map is precisely
\[
 \begin{gathered}
 \det H_K\longrightarrow U_K\otimes(H_K/U_K)
 \xrightarrow{\ \exp^{-1}\otimes\exp^*\ }
 t_K\otimes F^0_K,\\
 t_K\otimes F^0_K\longrightarrow\det E_K,
       \qquad u\otimes v\longmapsto v\wedge u.
 \end{gathered}
\tag{AG-local-generic-splice}
\]
This verifies both the Euler-factor cancellation and
the remaining order; in particular its last arrow
is $-u\wedge v$ in the opposite order.

By construction the formal dual exponential is the
adjoint of the exponential for the actual polarized
dual finite sequence.  Hence
$q(\exp(u),y)=\beta(u,\exp^*(y))/(\zeta/\tau)$.
Equation (AG-local-generic-splice) consequently gives
\[
            \Theta_K(x\wedge y)
                  =-q(x,y)\,\zeta/\tau
                    \qquad(x\in U_K).
\]
Both sides come from regular $R$-linear maps
$U_R\otimes H_R\to\det E_R$.  Their equality over
$K$ implies equality over the domain $R$, and then
at $z=0$.  For the actual crystalline representation,
the Weil--Deligne representation is unramified with
$N=0$, so its epsilon constant is one by Nakamura
\S3.3, p.~35, property (4).  The fixed filtration
gives the fixed factor
$\Gamma(V)=(-1)^{k-1}$ for weights $(k,1-k)$.
Multiplication by this factor and $\theta_{\mathrm{dR}}$
proves (AG-local-splice).

Nondegeneracy of $q$ gives $wx=-\Gamma(V)x$ on $U$.
Since $B_\epsilon$ is alternating, $w$ is skew-adjoint
for the symmetric form $q$, so it has trace zero.
Its other eigenvalue is therefore $\Gamma(V)$;
the two eigenvalues are distinct.  Thus the indicated
eigenspace is exactly $U$.  On the complementary locus
$D^{\varphi=1}=0$, AG-LOCAL-EXP makes both maps used
in BKNO Proposition 4.17 actual isomorphisms.  The
displayed calculation in its complete proof,
pp.~28--29, then gives the same formula with the
stated Weil--Deligne epsilon factor.  The new argument
above fills the exceptional case where that proof's
inverse exponential is unavailable.
\end{proof}

This repairs the rational de Rham determinant comparison
for the local signed line.  It does not make the ordinary
exponential or dual exponential nonzero in
AG-LOCAL-EXCEPTION, and it does not supply an integral
theta-period comparison or a global arithmetic cycle.
For a noncrystalline rank-two CM realization with
determinant $L(1)$, AG-LOCAL-WD already removes this
particular exponential obstruction.  The universal
epsilon-isomorphism and its global applications remain
the separate results cited by BKNO; the proof here
concerns their precise de Rham determinant specialization.

\subsection{Effective integral denominators for the Frobenius grading}

Let $X/\mathbb F_q$ be smooth, projective and geometrically
connected of dimension $d$.  Use $\ell$-adic cohomology for
$\ell\ne p$, crystalline cohomology after inverting $p$,
or a rational cohomology satisfying Katz--Messing's stated
Poincaré duality, weak Lefschetz and zeta-function hypotheses.
Let $F$ act by Frobenius pullback and set
\[
 G_i(T)=\det(T-F\mid H^i(X)),\qquad b_i=\deg G_i,
 \qquad P(T)=\prod_{i=0}^{2d}G_i(T).
\]
These monic polynomials lie in $\mathbb Z[T]$, and the
roots of $G_i$ have every complex absolute value $q^{i/2}$.
Thus distinct nonconstant $G_i$ are relatively prime.
Katz--Messing, Theorem 1 and Theorem 2(1), pp.~74--77,
show that the same $G_i$ work in crystalline cohomology
and that the rational grading projectors are algebraic
polynomials in Frobenius.  The following formulation gives
an explicit denominator and its exact integral meaning.

We use the correspondence action
$C_*(x)=p_{2*}(p_1^*x\cup\operatorname{cl}(C))$.
Accordingly $\Phi$ denotes the transpose graph of the
$q$-power Frobenius morphism, so its action is $F$;
its powers are again integral graph correspondences.
All assertions of idempotence below concern their actions
on cohomology, equivalently the specified homological
quotient of the correspondence ring.

\begin{proposition}[AG-FROB-PROJECTOR: an explicit resultant denominator]
For an index $i$ with $b_i>0$, put
\[
 Q_i=\prod_{j\ne i}G_j,\qquad
 R_i=\operatorname{Res}(G_i,Q_i)\in\mathbb Z\setminus\{0\}.
\]
There is an integral cycle $C_i$, an integer linear
combination of transpose graphs of Frobenius iterates and the diagonal,
whose action on $H^j(X)$ is $R_i\delta_{ij}$.
Hence $\pi_i=C_i/R_i$ is a rational algebraic K\"unneth
projector.  Its displayed denominator satisfies
\[
 0<|R_i|\le
 \prod_{j\ne i}
       (q^{i/2}+q^{j/2})^{b_i b_j}.
\tag{AG-frob-resultant}
\]
One common denominator for all these projectors is
$\mathcal D=\prod_{i<j}|\operatorname{Res}(G_i,G_j)|$,
where factors involving $G_i=1$ are one.
\end{proposition}
\begin{proof}
The adjugate of the integral Sylvester matrix gives polynomials
$A_i,B_i\in\mathbb Z[T]$ satisfying
\[
                  A_iG_i+B_iQ_i=R_i.
\]
Indeed the Sylvester map sends the coefficient vectors of
polynomials of degrees less than $\deg Q_i$ and $\deg G_i$
to those of $A_iG_i+B_iQ_i$; multiplying its adjugate by
the coefficient vector of the constant polynomial produces
the displayed identity, up to the harmless common sign
fixed by the convention for the resultant.
Set $C_i=B_i(\Phi)Q_i(\Phi)$.  Cayley--Hamilton gives
$G_j(F)=0$ on $H^j$.  For $j\ne i$ the factor $G_j$
in $Q_i$ makes the action zero.  For $j=i$ the Bezout
identity makes it $R_iI$.  This proves the cycle assertion
and the cohomological projector identities.

Writing $\alpha$ and $\beta$ for the roots of $G_i$ and
$G_j$ in a common complex algebraic closure gives
\[
 |\operatorname{Res}(G_i,G_j)|
   =\prod_{\alpha,\beta}|\alpha-\beta|
   \le(q^{i/2}+q^{j/2})^{b_i b_j}.
\]
Multiplicativity of resultants proves (AG-frob-resultant).
Each $|R_i|$ divides $\mathcal D$, proving the common
denominator statement.  Zero cohomology groups simply
contribute the polynomial one.
\end{proof}

\begin{proposition}[AG-FROB-ORDER: the exact arithmetic order behind the certificate]
The natural injection of finite free abelian groups
\[
 A=\mathbb Z[T]/(P)\longrightarrow
 B=\prod_i\mathbb Z[T]/(G_i)
\]
has cokernel of order $\mathcal D$.  Let $e_i\in B$
be one in its $i$th factor and zero elsewhere.  The least
positive integer $a_i$ for which $a_ie_i$ belongs to $A$
is exactly the order of the class of $e_i$ in $B/A$;
it divides $|R_i|$.  It is computable by the Smith normal
form of the displayed integral map.  This is the exact
denominator in the abstract Frobenius polynomial order;
an action on a particular cohomology representation can
have further relations and a smaller denominator.
\end{proposition}
\begin{proof}
For two relatively prime monic polynomials $G,Q$ there
is an exact sequence
\[
 0\longrightarrow\mathbb Z[T]/(GQ)
 \longrightarrow\mathbb Z[T]/(G)\oplus\mathbb Z[T]/(Q)
 \xrightarrow{\ (a,b)\mapsto a-b\ }
 \mathbb Z[T]/(G,Q)\longrightarrow0.
\tag{AG-frob-crt}
\]
For injectivity, a polynomial divisible by both $G$ and
$Q$ is divisible by $GQ$ over $\mathbb Q[T]$, and monic
division makes the quotient integral.  If
$a-b=Gu+Qv$, the polynomial $a-Gu=b+Qv$ lifts the
pair, proving exactness in the middle.  The last map is
surjective.  Its target is the cokernel of multiplication
by $Q$ on the free group $\mathbb Z[T]/(G)$.
The determinant of that multiplication map is
$\operatorname{Res}(G,Q)$, so the target has its absolute
value as order.  Induct on the number of factors and use
the multiplicativity of the resultant.  The indices
multiply to give $\mathcal D$.

The assertion about $a_i$ is the definition of the order
of an element in the finite quotient.  The Bezout identity
in AG-FROB-PROJECTOR proves $R_ie_i\in A$ and hence
$a_i\mid |R_i|$.  Choosing the monomial bases in all
factors gives an explicit integral matrix for $A\to B$;
its Smith normal form computes the order of each $e_i$
class.  This also explains the last qualification: a
representation annihilated by $P$ need not be faithful
as an $A$-module.
\end{proof}

\begin{proposition}[AG-FROB-COLLISION: exact bad primes for Frobenius-only projectors]
Fix a prime $\ell\ne p$ and put
$\Lambda_j=H^j_{\text{ét}}(X_{\overline{\mathbb F}_q},
\mathbb Z_\ell)/\mathrm{torsion}$.
A polynomial $f\in\mathbb Z_\ell[T]$ acting as the identity
on $\Lambda_i$ and as zero on every $\Lambda_j$, $j\ne i$,
exists if and only if $\ell\nmid R_i$.
This exact obstruction concerns the algebra generated by
Frobenius, and need not obstruct an integral algebraic
K\"unneth projector supplied by other cycles.
\end{proposition}
\begin{proof}
If $R_i$ is a unit, the Bezout polynomial
$f=B_iQ_i/R_i$ is integral over $\mathbb Z_\ell$.
The characteristic polynomial of $F$ on the free lattice
$\Lambda_j$ is $G_j$, because it becomes that polynomial
after tensoring with $\mathbb Q_\ell$.  Cayley--Hamilton
therefore proves the required identities integrally.

Conversely, $\ell\mid R_i$ means that the reductions of
$G_i$ and some $G_j$, $j\ne i$, have a common root
$\alpha\in\overline{\mathbb F}_\ell$.
On both nonzero spaces
$\Lambda_i\otimes\overline{\mathbb F}_\ell$ and
$\Lambda_j\otimes\overline{\mathbb F}_\ell$, that root
is an eigenvalue of the reduced Frobenius.  Applying a
putative $f$ to eigenvectors forces simultaneously
$\bar f(\alpha)=1$ and $\bar f(\alpha)=0$, a contradiction.
No semisimplicity assumption is needed.

For the distinction in the last sentence, take
$X=\mathbb P^1_{\mathbb F_q}$.  Here
$G_0=T-1$, $G_2=T-q$, and
\[
 \pi_0=\frac{q-\Phi}{q-1},\qquad
 \pi_2=\frac{\Phi-1}{q-1}.
\]
The least denominator in $\mathbb Z[\Phi]$ is exactly
$q-1$: evaluation of any integer polynomial at $q$ and
$1$ gives values congruent modulo $q-1$, and the displayed
formulas attain that bound.  Yet, for a rational point
$x\in\mathbb P^1(\mathbb F_q)$, the two integral cycles
$[x\times\mathbb P^1]$ and $[\mathbb P^1\times x]$
are precisely the two K\"unneth projectors in our convention.
Thus every prime dividing $q-1$ is a genuine obstruction
within the Frobenius polynomial algebra, even though these
integral cycle projectors already exist.
\end{proof}

For cote 16, this separates three pieces of the problem:
finite-field purity supplies the rational algebraic grading;
the exact resultants control its Frobenius-generated integral
order; and an integral inverse Lefschetz correspondence still
requires its own cycle construction.  Neither the displayed
projectors nor their denominators prove idempotence modulo
rational equivalence in the full Chow ring.  The latter
would require control of the homologically trivial error,
for example an appropriate nilpotence theorem.  The
corresponding analytic inputs A and M do not enter these
already-pure finite-field deductions.  The next subsection
supplies the exact nilpotence lift and, in the Kimura-finite
sector, proves the stronger rational Chow identity directly.

\subsection{Chow projectors: nilpotence, parity and exact denominator primes}

Keep $G_i,P,R_i,C_i,\Phi$ from AG-FROB-PROJECTOR, omitting
the indices with $G_i=1$, and put
$\mathcal D=\prod_{i<j}|\operatorname{Res}(G_i,G_j)|$.
All products of correspondences in this subsection mean
composition.  The following elementary correction states
precisely which ring must satisfy the nilpotence hypothesis.

\begin{proposition}[AG-FROB-NIL: an integral polynomial correction]
Let $A$ be a unital algebra over
$R=\mathbb Q$ or $R=\mathbb Z[1/\mathcal D]$, containing
the specified element $\Phi$.  Suppose that, in this ring,
\[
                         P(\Phi)^N=0\qquad(N\ge1).
\tag{AG-chow-nil}
\]
Put $e_i=C_i/R_i$ and define
\[
 s_N(x)=\sum_{j=N}^{2N-1}
             {2N-1\choose j}x^j(1-x)^{2N-1-j},
                   \qquad \widetilde e_i=s_N(e_i).
\tag{AG-chow-correction}
\]
Then the $\widetilde e_i$ are pairwise orthogonal
idempotents summing to one, and
\[
                 G_i(\Phi)^N\widetilde e_i=0.
\tag{AG-chow-primary}
\]
They are the unique idempotents in the commutative
algebra $R[\Phi]$ congruent to the respective $e_i$
modulo its nilpotent ideal $(P(\Phi))$.
As rational combinations of integral graph cycles,
$\widetilde e_i$ have denominator dividing
$R_i^{\,2N-1}$; one common denominator divides
$\mathcal D^{\,2N-1}$.

For $A=\operatorname{CH}^d(X\times X)_{\mathbb Q}$,
the assertions are rational Chow identities.  They
are identities in
$\operatorname{CH}^d(X\times X)_{\mathbb Z[1/\mathcal D]}$
if (AG-chow-nil) is known in that localized ring.
Knowing it only after tensoring with $\mathbb Q$
does not assert those integral identities.
\end{proposition}
\begin{proof}
The binomial theorem gives
\[
 s_N(x)+s_N(1-x)=1,\qquad
 x^N\mid s_N(x),\qquad (1-x)^N\mid1-s_N(x)
                         \quad\hbox{in }\mathbb Z[x].
\]
Thus $x^N(1-x)^N$ divides $s_N(x)(1-s_N(x))$.
The degree is exactly $2N-1$: its leading coefficient
is $(-1)^{N-1}{2N-2\choose N-1}\ne0$, by the
alternating binomial-sum identity obtained from
Pascal's rule.  This also covers $N=1$.

Write $e_i=E_i(\Phi)$, where
$E_i(T)=B_i(T)Q_i(T)/R_i$.  The Bezout construction
gives $E_i\equiv\delta_{ij}\pmod{G_j}$.
Consequently
\[
 s_N(E_i(T))\equiv\delta_{ij}\pmod{G_j(T)^N}.
\]
The polynomials $G_j^N$ are pairwise comaximal over
$R$, because their pairwise resultants are units.
The Chinese remainder theorem therefore makes
\[
 \begin{gathered}
 s_N(E_i)^2-s_N(E_i),\quad
 s_N(E_i)s_N(E_j)\ (i\ne j),\\
 \sum_i s_N(E_i)-1,
 \quad G_i^N s_N(E_i)
 \end{gathered}
\]
all divisible by $P^N$.  Evaluation at $\Phi$ and
(AG-chow-nil) prove every stated identity.

The same congruences show that $s_N(E_i)-E_i$ is
divisible by $P$, so these idempotents lift the
claimed classes modulo $(P(\Phi))$.

For uniqueness, any two commuting idempotents $a,b$
with $a-b$ in a nilpotent ideal satisfy
$a(1-b)=(a-b)(1-b)$ in that ideal.  The left side
is itself idempotent, hence is zero.  The same
argument gives $b(1-a)=0$, so $a=b$.
The elements in question commute because they are
polynomials in $\Phi$.  This argument also explains
why commuting idempotent products that are nilpotent
must vanish; it does not assume a semisimple ring.

Finally (AG-chow-correction) can be written as
\[
 \widetilde e_i=
 \frac{1}{R_i^{\,2N-1}}
 \sum_{j=N}^{2N-1}{2N-1\choose j}
                  C_i^j(R_i-C_i)^{2N-1-j}.
\]
Every numerator is an integral polynomial in graph
correspondences, proving the denominator claim.
These calculations take place in the chosen ring
$A$, which proves the last distinction as well.
\end{proof}

There is a sharper answer for Kimura-finite Chow motives
over the finite base field itself.  We use the following
located classical tensor facts from André,
\emph{Motifs de dimension finie}, Bourbaki no.~929,
Astérisque 299 (2005): Theorem 2.8, p.~124, gives
Kimura finiteness for motives of abelian type;
Corollaries 3.19--3.22, p.~132, identify the even/odd
dimension and give categorical Cayley--Hamilton;
\S4.2, p.~138, states that Frobenius is a tensor
automorphism of the identity on motives over
$\mathbb F_q$ for every adequate equivalence.
The coefficient field is $\mathbb Q$; the geometric
base field is allowed to have positive characteristic.

\begin{proposition}[AG-FROB-CHOW: exact Cayley--Hamilton in a finite-dimensional motive]
If $h(X)$ is Kimura finite over $\mathbb F_q$, then
\[
                     P(\Phi)=0
       \quad\text{in }\operatorname{CH}^d(X\times X)_{\mathbb Q}.
\tag{AG-chow-cayley}
\]
Consequently the original $e_i=C_i/R_i$, with no
nilpotence correction, are pairwise orthogonal
rational Chow projectors summing to the diagonal,
and $G_i(\Phi)e_i=0$ in that ring.  They are
represented by cycles whose denominators divide
$R_i$.  This applies in particular when $h(X)$ is
of abelian type over $\mathbb F_q$.

For every prime $\ell\ne p$, the rational degree-$i$ Chow
projector can be represented by a polynomial in
$\Phi$ with coefficients in $\mathbb Z_{(\ell)}$
if and only if $\ell\nmid R_i$.  Thus the set of
necessary denominator primes away from $p$ is exact
inside the Frobenius polynomial algebra.  These
rational Chow identities do not by themselves give
idempotence in the localized integral Chow ring.
\end{proposition}
\begin{proof}
Choose a Kimura decomposition defined over the base
field,
\[
                   h(X)=M_+\oplus M_-.
\]
An exterior power of $M_+$ vanishes and a symmetric
power of $M_-$ vanishes.  Realization in super vector
spaces shows that $H(M_+)$ is purely even and
$H(M_-)$ purely odd: a nonzero odd vector has
nonzero exterior powers of all orders in the super
convention, and a nonzero even vector has nonzero
symmetric powers of all orders.  Thus these two
realizations are exactly $H^{\mathrm{even}}(X)$ and
$H^{\mathrm{odd}}(X)$.  Frobenius is a natural
endomorphism of the identity category, so it
commutes with the two splitting idempotents and
preserves $M_+$ and $M_-$.

We make the characteristic-polynomial normalization
explicit.  For an even object $M$ of dimension $m$
and an endomorphism $f$, categorical Cayley--Hamilton
uses the monic polynomial
\[
 c_f(T)=\sum_{j=0}^m(-1)^j
                 \operatorname{tr}(\wedge^j f)T^{m-j}.
\tag{AG-categorical-char}
\]
It annihilates $f$.  In terms of the reciprocal
polynomial $\det(1-tf)$ used in the cited discussion,
this is $T^m\det(1-T^{-1}f)$.
One can check this normalization directly: partially
trace the vanishing antisymmetrizer on $M^{\otimes(m+1)}$
after applying $1\otimes f^{\otimes m}$.  Group
permutations by the length $j+1$ of the cycle through
the first position.  There are
${m\choose j}j!$ choices for that cycle, and summing
the remaining signed permutations contributes
$(m-j)!\operatorname{tr}(\wedge^{m-j}f)$.
Division by $(m+1)!$ gives
\[
 0=\frac1{m+1}\sum_{j=0}^m(-1)^j
                  \operatorname{tr}(\wedge^{m-j}f)f^j,
\]
which is $(-1)^mc_f(f)/(m+1)=0$.
The cited dimension result ensures
$\wedge^{m+1}M=0$ for this $m$; an even object of
dimension zero is zero.  This is the usual classical
tensor proof, included to fix the reciprocal convention.

Apply (AG-categorical-char) to $\Phi|M_+$.
Trace compatibility with realization identifies its
polynomial with
\[
                         P_+(T)=\prod_{i\ {
m even}}G_i(T).
\]
For the odd part, pass to the category of formally
$\mathbb Z/2$-graded objects in the same rational
tensor category.  The original category embeds
fully faithfully as the degree-zero objects.
Let $\Pi$ be the degree-one unit, whose self-braiding
is $-1$ and whose square is the unit.
Then $M_-\otimes\Pi$ is even, because
\[
 \wedge^j(M_-\otimes\Pi)
              \simeq \operatorname{Sym}^j(M_-)\otimes\Pi^j.
\]
Give $\Pi$ the identity endomorphism.  Tensoring by
it changes the categorical trace sign of $M_-$.
The realization of $M_-$ was purely odd, so the
new trace and all exterior coefficients are exactly
the ordinary traces and exterior coefficients of
$H^{\mathrm{odd}}(X)$.  Hence the monic polynomial
for $\Phi|M_-\otimes1_\Pi$ is
\[
                         P_-(T)=\prod_{i\ {
m odd}}G_i(T).
\]
Cayley--Hamilton makes it zero.  Tensoring back by
the invertible $\Pi$ and using the faithful embedding
gives $P_-(\Phi)|M_-=0$ in the original Chow category.
Likewise $P_+(\Phi)|M_+=0$.  Their product
$P=P_+P_-$ therefore annihilates all of $h(X)$,
proving (AG-chow-cayley).

Apply AG-FROB-NIL with $N=1$, for which $s_1(x)=x$.
This proves the Chow projector and denominator
claims.  If $\ell\nmid R_i$, the displayed polynomial
$C_i/R_i$ has coefficients in $\mathbb Z_{(\ell)}$.
If $\ell\mid R_i$ and $\ell\ne p$, any polynomial
with such coefficients would act integrally on the
free $\ell$-adic lattices and give the same degree
projector.  AG-FROB-COLLISION rules this out by the
common reduced Frobenius eigenvalue.  This proves
the exact prime-set assertion.  The tensor proof
was in the category with rational coefficients;
it has not proved that its zero correspondence
remains zero before torsion is removed.
\end{proof}

\begin{proposition}[AG-FROB-CYCLES: cycle degrees and a conditional integral annihilator]
Under the rational hypothesis of AG-FROB-NIL,
\[
      \widetilde e_i\operatorname{CH}^r(X)_{\mathbb Q}=0
                         \qquad(i\ne2r).
\tag{AG-chow-degrees}
\]
This holds in particular for the original $e_i$ of
AG-FROB-CHOW.  It places all rational codimension-$r$
cycles in the degree-$2r$ Chow summand, but does not
assert that the cycle map on that summand is injective
or surjective.

If (AG-chow-nil) holds in the localized integral
Chow correspondence ring, then, for $0\le r\le d$
and $i\ne2r$, the integral image
$\widetilde e_i\operatorname{CH}^r(X)_{\mathbb Z[1/\mathcal D]}$
is annihilated by the nonzero integer
$G_i(q^r)^N$.  Moreover
\[
       0<|G_i(q^r)^N|
                      \le(q^r+q^{i/2})^{Nb_i}.
\tag{AG-chow-annihilator}
\]
\end{proposition}
\begin{proof}
The $q$-power Frobenius pulls a codimension-$r$
cycle back to $q^r$ times itself.  This can be seen
integrally at the generic point of each irreducible
cycle: the smooth ambient variety has a regular
local ring of dimension $r$ there, and Frobenius
replaces a regular parameter ideal by its $q$th
powers, whose colength is $q^r$.  Frobenius is
finite flat on a smooth variety over the perfect
base field (étale locally this is the free
monomial-basis map on affine space), so this is its
flat pullback on Chow groups.  Our transpose-graph
convention is precisely that pullback.

Apply (AG-chow-primary) to
$\widetilde e_i\operatorname{CH}^r$.
Frobenius acts there by $q^r$, so the relation gives
$G_i(q^r)^N\widetilde e_i z=0$ for every cycle $z$.
All complex absolute values of roots of $G_i$ are
$q^{i/2}$; if $i\ne2r$, none can equal $q^r$.
Thus $G_i(q^r)$ is a nonzero integer.  Dividing
over $\mathbb Q$ proves (AG-chow-degrees), and
multiplying the bounds
$|q^r-\alpha|\le q^r+q^{i/2}$ over its roots proves
(AG-chow-annihilator).  When the relation was assumed
only rationally, this argument gives the rational
vanishing only; no integral annihilator is then claimed.
\end{proof}

For clarity, denominators and equalities are separate
issues.  A cycle in
$\operatorname{CH}^d(X\times X)_{\mathbb Z[1/\mathcal D]}$
can become zero over $\mathbb Q$ while remaining
nonzero torsion over that localization.  If the map
from this localized Chow group to its rationalization
is injective, the rational projector identities
above do hold integrally; alternatively the explicit
localized nilpotence hypothesis of AG-FROB-NIL suffices.
Without either input the rational tensor theorem alone
does not decide the integral identities.  Algebraically,
$\mathbb Z[t]/(\ell t,t^2-t)$ has a nonzero torsion
idempotent $t$ invisible after tensoring with $\mathbb Q$;
this illustrates the issue and is not offered as a
geometric Chow-ring example.

Thus in the finite-field Artin--Tate and abelian/CM
motivic sector, classical parity and Frobenius already
promote the grading to rational Chow correspondences
with the explicit resultants above.  This is a genuine
refinement of the cote 16 grading subproblem, without
a new GRH premise.  The torsion-sensitive integral
correspondence ring, the inverse Lefschetz cycle and
the general non-finite-dimensional motive retain their
own hypotheses; neither A nor M supplies them here.

\begin{proposition}[AG-FROB-Z: an effective obstruction to one coefficient localization]
For every odd prime $\ell$, the adopted input A supplies
a rational prime $p$ with
\[
                p\equiv1\pmod\ell,
                     \qquad p<2(\ell\log\ell)^2,
\tag{AG-frob-Z-bound}
\]
such that neither cohomological degree projector of
$\mathbb P^1_{\mathbb F_p}$ is a polynomial in
Frobenius with coefficients in $\mathbb Z_\ell$.
Nevertheless both projectors are restrictions of
integral product-cycle projectors on
$\mathbb P^1_{\mathbb Z}$.

More generally, for any two finite sets of rational
primes $S,T$, there is a prime $p\notin T$ where
these degree projectors cannot be represented by
polynomials in Frobenius over $\mathbb Z[S^{-1}]$.
An explicit witness is provided by (AG-frob-Z-bound)
for any odd prime $\ell\notin S$ greater than every
element of $T$.  Thus shrinking the arithmetic base
at finitely many primes and inverting finitely many
coefficient primes do not make this Frobenius-only
grading construction integral at every remaining
reduction.
\end{proposition}
\begin{proof}
Bach--Sorenson, \emph{Explicit bounds for primes in
residue classes}, Theorem 5.3, printed pp.~1731--1732,
gives a prime $p\equiv m\pmod q$ with
$p<2(q\log q)^2$ under ERH, for the coprime residue
classes in question.  Apply it with $m=1,q=\ell$.
The analytic functions in this cyclotomic case are
the primitive Dirichlet factors, so A supplies the
required hypothesis.  This is the residue-class
specialization of the same explicit-prime source
used in CX-GEO-004, with its sharper Theorem 5.3
constant.  The classical published numerical
estimate is an input here; its computational
verification is not being re-created.

Frobenius on $H^0$ and $H^2$ of $\mathbb P^1$ has
eigenvalues $1,p$.  A polynomial producing one of
the two degree projectors would have values $1,0$
at these two arguments, in one order or the other.
But every $f\in\mathbb Z_\ell[T]$ satisfies
$f(p)\equiv f(1)\pmod\ell$ when $p\equiv1\pmod\ell$.
This is a contradiction, also the special case
$R_i=p-1$ of AG-FROB-COLLISION.  Here $p>\ell$,
so the distinct-coefficient-prime requirement
$\ell\ne p$ is automatic.

The section at infinity of $\mathbb P^1_{\mathbb Z}$
gives the two integral relative product cycles
\[
 [\{\infty\}\times\mathbb P^1],\qquad
 [\mathbb P^1\times\{\infty\}].
\]
Their restrictions to every fiber are the two
integral projectors already computed in
AG-FROB-COLLISION.  Thus the obstruction belongs
to the polynomial algebra generated by the
individual fiber's Frobenius, not to the availability
of integral algebraic cycles.

Finally choose $\ell\notin S$ greater than every
prime in $T$.  The prime from (AG-frob-Z-bound)
satisfies $p>\ell$, so it is outside $T$.
The coefficient ring $\mathbb Z[S^{-1}]$ embeds
in $\mathbb Z_\ell$, and any proposed polynomial
there would contradict the argument above.
This proves the precise two-localization statement.
Existence of such obstructing fibers without an
explicit bound already follows from Dirichlet's
theorem; A supplies the displayed effective witness.
\end{proof}


\subsection{Actual inverse Lefschetz cycles on products of elliptic curves}
\label{ag:lefschetz-geometric}

The inherited statement H8 (run 20261112), as used in the Hodge
sketch Route~4, identifies the exterior-algebra map with a subset
inclusion matrix.  RETURN\_DP\_MULT, DP15 and Q4e, already check
its three-dimensional $J-I$ obstruction at cote~16, scan~14,
section~7.  We retain that attribution.  The following formulation
adds an actual cycle realization and the exact least denominator;
it does not identify the separate H8(iv) order or an automorphic
lifting lattice.

\begin{proposition}[AG-LEFSCH-SMITH: the full geometric lattice]
Let $g\ge1$ and $X=\prod_{i=1}^{g}E_i$ be a product of elliptic curves over a
subfield $F\subset\mathbb C$, with their origins $o_i$.  Put
$\theta=\sum_i\operatorname{pr}_i^*[o_i]$, and let $L$ denote
cup product with $\theta$.  Fix $0\le k\le g$ and put $a=g-k$.
The map $L^a:H^k(X(\mathbb C),\mathbb Z)\longrightarrow
H^{2g-k}(X(\mathbb C),\mathbb Z)$ is injective, and its finite
cokernel is the following direct sum of cyclic groups:
\begin{equation}
 \begin{gathered}
 \displaystyle\bigoplus_{r=0}^{\lfloor k/2\rfloor}\ 
 \bigoplus_{i=0}^{r}
 \left(\mathbb Z\Big/
 a!\binom{r+a-i}{r-i}\mathbb Z\right)^{m_{r,i}},\\
 m_{r,i}=\binom{g}{k-2r}2^{k-2r}
 \left[\binom{g-k+2r}{i}-\binom{g-k+2r}{i-1}\right].
 \end{gathered}
 \tag{AG-lefschetz-coker}
\end{equation}
Here $\binom n{-1}=0$, and a cyclic group of order one is zero.
The displayed list is a diagonal form; it need not be in
invariant-factor order.  In particular, the exact least positive
integer clearing the cohomological inverse is
\begin{equation}
 D_{g,k}=a!\,d_{g,k},\qquad
 d_{g,k}=\operatorname{lcm}_{0\le j\le\lfloor k/2\rfloor}
                         \binom{a+j}{j}.
 \tag{AG-lefschetz-denominator}
\end{equation}
For the integral divided-power map $L^{[a]}=L^a/a!$, remove
$a!$ from every cyclic order in (AG-lefschetz-coker); its exact
inverse denominator is $d_{g,k}$.
\end{proposition}
\begin{proof}
Integral K\"unneth gives a direct sum indexed by the subset $S$
of factors contributing $H^1$.  Write $s=|S|$ and $k=s+2r$.
The $H^1$ tensor has rank $2^s$; each remaining factor contributes
its generator in degree zero or two.  If $n=g-s=2r+a$, the
source basis in this block is indexed by the $r$-subsets $J$
of an $n$-set, and the target by its $(r+a)$-subsets $K$.
All raising operators have even degree.  Their squares vanish,
so the matrix is
\[
 a!\bigl(\mathbf 1_{J\subseteq K}\bigr)_{K,J}
       =a!\,W_{r,r+a}(2r+a)^{\mathsf t}.
\]
Wilson's integral diagonal-form theorem applies because
$r\le r+a\le(2r+a)-r$.  Its diagonal entries are
$\binom{r+a-i}{r-i}$ with multiplicity
$\binom ni-\binom n{i-1}$, for $0\le i\le r$.
A primary proof is Ghorbani--Khosrovshahi--Maysoori--Mohammad-Noori,
\emph{Inclusion Matrices and Chains}, Corollary~3, author p.~10;
Theorem~1 and its complete induction occupy pp.~8--9,
with the unimodular equivalence defined on p.~7.  Thus this
is equivalence over $\mathbb Z$, not just a rational eigenvalue
calculation.  Multiplication by $a!$ preserves the same two
unimodular changes of basis.  There are $\binom gs$ subsets
$S$ and $2^s$ copies for each, proving the asserted direct sum
and injectivity.

For an injective map of free lattices of the same rank,
$N(L^a)^{-1}$ is integral if and only if $N$ annihilates its
cokernel.  Its least such $N$ is therefore the least common
multiple of the diagonal entries.  Replacing $i$ by $r-j$
gives exactly the binomials $\binom{a+j}{j}$ with
$0\le j\le r$.  The block $r=\lfloor k/2\rfloor$ occurs
with $s=0$ or $1$ and nonzero $H^1$ tensor, so every entry
needed for (AG-lefschetz-denominator) really occurs.  Finally,
the same subset expansion gives an integral matrix for
$L^{[a]}$, and removes precisely the common factorial.
\end{proof}

\begin{proposition}[AG-LEFSCH-CYCLES: the sharp denominator is attained by cycles]
In AG-LEFSCH-SMITH there are integral Chow projectors
$\Pi_j\in\mathrm{CH}^{g}(X\times_F X)$ for every degree,
and a correspondence
\[
 C_k\in\mathrm{CH}^{g-a}(X\times_F X)
                   \otimes\mathbb Z[1/D_{g,k}]
\]
whose coefficients can in fact all be taken in
$D_{g,k}^{-1}\mathbb Z$, such that, in the Chow
correspondence ring with coefficients $\mathbb Z[1/D_{g,k}]$,
\begin{align*}
 C_k&=\Pi_k C_k\Pi_{2g-k},\\
 C_k L^a\Pi_k&=\Pi_k,&
 L^a C_k\Pi_{2g-k}&=\Pi_{2g-k}.
\end{align*}
Here $L=\Delta_{X*}\theta$ denotes the algebraic correspondence.
There is likewise an inverse to $L^{[a]}$ with coefficients
in $d_{g,k}^{-1}\mathbb Z$.  These denominator bounds are sharp
even if arbitrary algebraic correspondences are allowed:
there is an integral correspondence inducing
$N(L^a)^{-1}$ on $H^{2g-k}$ if and only if $D_{g,k}\mid N$;
replace $D_{g,k}$ by $d_{g,k}$ for the divided-power map.
No A, M, Hodge conjecture, or Chow-nilpotence hypothesis is used.
\end{proposition}
\begin{proof}
Use the action convention
$Z_*\alpha=(p_2)_*(p_1^*\alpha\cdot Z)$.
On a pointed elliptic curve $(E,o)$ the following are integral
cycles, with the indicated correspondence degrees:
\[
 \begin{gathered}
 \pi_0=[o\times E],\quad \pi_2=[E\times o],\quad
 \pi_1=[\Delta_E]-\pi_0-\pi_2,\\
 U=[o\times o],\qquad V=[E\times E].
 \end{gathered}
\]
The $\pi_j$ have degree zero, $U$ degree $1$, and $V$ degree
$-1$.  Direct composition of product cycles gives
\[
 VU=\pi_0,\qquad UV=\pi_2,\qquad U^2=V^2=0.
\]
Moreover $\pi_0,\pi_1,\pi_2$ are orthogonal idempotents
summing to the diagonal; $U$ and $V$ kill the $\pi_1$ summand.
For example $VU$ is the pushforward of $[o,o,E]$ from
$E^3$, and $UV$ that of $[E,o,o]$.  The other relations follow
by the same product intersections, using $[o]^2=0$ in the
Chow ring of a curve.  Consequently
$\pi_0,\pi_2,U,V$ are exact matrix units on the two Tate
summands, already modulo rational equivalence with integral
coefficients.  These are identities over the field $F$.

Take external products of the three projectors and sum all
patterns of total degree $j$ to obtain $\Pi_j$.
Fix a subset $S$ carrying $\pi_1$, and set
$N_S=\bigotimes_{i\in S}(E_i,\pi_1)$.
The other factors are copies of the two Tate summands.
External products of $U,V,\pi_0,\pi_2$ give matrix units
between every two patterns $J,K$ on those factors, tensored
with the identity on $N_S$.  This needs no algebraic basis
or splitting of the $H^1$ of any elliptic curve.
Composition of these cycles is exactly composition of their
integer matrices.

On this actual Chow summand, $L$ is the sum of the raising
matrix units $U_i$.  Indeed $\Delta_{E_i*}[o_i]=[o_i\times o_i]$,
and $U_i$ kills the $\pi_1$ factor.  The $U_i$ commute and
square to zero, so the matrix of $L^a$ is the very same
$a!W_{r,r+a}(2r+a)^{\mathsf t}$ used above, now as a
matrix of Chow-motive morphisms.  Invert this square matrix
over $\mathbb Q$.  Its inverse entries belong to
$D_{g,k}^{-1}\mathbb Z$ by its integral diagonal form.
Replace each entry by the corresponding product-cycle
matrix unit and sum over $S$.  A unit taking a pattern of
size $r+a$ to one of size $r$ has degree $-a$, so the result
has codimension $g-a$ on $X\times X$.
The exact matrix-unit relations prove all three displayed
Chow identities over $\mathbb Z[1/D_{g,k}]$ itself:
they arise from identities between integer matrices after
clearing denominators, so no hidden Chow torsion is discarded.
In particular this construction does not
lift merely homological identities by an unspecified
nilpotence theorem.  The divided-power construction is the
same, with the factorial removed; its raising correspondence
is already integral by the subset expansion.

Multiplying the constructed correspondence by $D_{g,k}$
produces an integral cycle.  Conversely, an integral
correspondence acts integrally on the torsion-free Betti
cohomology.  If its action on $H^{2g-k}$ is
$N(L^a)^{-1}$, the lattice criterion in AG-LEFSCH-SMITH
forces $D_{g,k}\mid N$.  This proves the asserted exact
optimality among all cycles, not only the product cycles.
\end{proof}

For $k=2$, the divided-power cokernel is exactly
$\mathbb Z/(g-1)$, recovering H8 and DP15 rather than
reclassifying them as new findings.  For $(g,k)=(5,4)$,
(AG-lefschetz-coker) gives
$\operatorname{coker} L\simeq(\mathbb Z/2)^{43}\oplus
\mathbb Z/6$, and the least denominator is $6$.
The product-intersection construction and its localized Chow
identities also work over any field, by the very same matrix-unit
relations; the global sharpness assertion above uses the stated
Betti lattice.  Thus one coefficient denominator works for this
construction on every good reduction of a product over a number
field.  The conclusion applies in particular to products
of CM elliptic curves used as test objects in the Hodge
sketch.  It closes the cycle-realization and denominator
comparison for this product polarization.  General
polarizations, nonsplit abelian varieties, the integral
order in H8(iv), and the specific theta-lift lattice still
require their own comparisons.


\subsection{Integral Fourier transforms: the exact cycle gate and a CM torsion obstruction}
\label{ag:fourier-cycle-gate}

HodgeSketch Route~4 cites Beckmann--de Gaay Fortman from memory,
and compares the theta-lift divisibilities with algebraic divided
powers and Fourier transforms.  We now locate the precise theorem
and distinguish its cohomological conclusion from identities of
integral Chow correspondences.  The latter distinction already
matters on the CM curve used in AG-C149-ELLIPTIC.

\begin{proposition}[AG-FOURIER-MIN: the classical minimal-class gate]
Let $(A,\theta)$ be a principally polarized complex abelian variety
of dimension $g\ge2$, and identify $A$ with its dual by the
polarization.  Write $\mathcal A^{2*}(A)$ for the image of the
integral cycle-class map, and put
$\gamma_c=\theta^c/c!\in H^{2c}(A,\mathbb Z)$.
The following assertions are equivalent:
\begin{enumerate}
\item $\gamma_{g-1}$ is the class of an integral one-cycle.
\item There is an integral, possibly inhomogeneous, cycle
$\Gamma\in\mathrm{CH}(A\times\widehat A)$ with
$\operatorname{cl}(\Gamma)=\operatorname{ch}(\mathcal P_A)$;
thus it induces the full integral cohomological Fourier transform.
\item Every integral Hodge class of degree $2g-2$ on $A$ is
algebraic.
\item For each positive-degree algebraic integral cohomology
class $x$ and every $n\ge1$, the integral class $x^n/n!$ is
algebraic.
\end{enumerate}
These are Beckmann--de Gaay Fortman, Theorem~1.1,
with the constructive proof in their Lemmas~3.4--3.5,
Theorem~3.8, Propositions~3.11--3.12 and section~4.1
(author pp.~8--13).  They imply the existence of an integral
Fourier transform \emph{up to homology}.  A weak integral
Fourier transform agreeing with the rational Chow transform
requires the stronger lift of
$\Theta^{g-1}/(g-1)!\in\mathrm{CH}_1(A)_{\mathbb Q}$ to
$\mathrm{CH}_1(A)$, where $\Theta$ is the symmetric polarization
class.  That stronger equivalence is their Theorem~3.8;
one must not replace this Chow class by its cohomology class.
\end{proposition}
\begin{proof}
Here is the construction, with its exact integral input.
The Pontryagin ring of an abelian variety has canonical divided
powers on its positive-dimensional cycles by Moonen--Polishchuk,
Theorem~1.6 and Corollary~1.7 (author pp.~6--7).
Suppose $z\in\mathrm{CH}_1(A)$ has class $\gamma_{g-1}$.
Let $j_1,j_2:A\to A\times A$ be the two coordinate inclusions,
and let $\Delta$ be the diagonal.  The integral cycle
\[
 R=(-1)^{g+1}(j_{1*}z+j_{2*}z-\Delta_*z)
\]
has cohomology class
$\rho=c_1(\mathcal P_A)^{2g-1}/(2g-1)!$ by
Beckmann--de Gaay Fortman's Lemma~3.5.  Define
\[
 \Gamma=(-1)^g\sum_{n=0}^{2g}\bigl((-1)^gR\bigr)^{[n]},
\]
using those Pontryagin divided powers; the zeroth term is the
origin.  Their Lemma~3.4 gives
$\operatorname{cl}(\Gamma)=\operatorname{ch}(\mathcal P_A)$.
Indeed $n!\operatorname{cl}(R^{[n]})=
\operatorname{cl}(R)^{\star n}$, and integral cohomology of
$A\times A$ is torsion-free, so these divided classes agree
with the rational exponential identity.  Terms of dimension
larger than $2g$ vanish.  This proves $1\Rightarrow2$ without
assuming that $z_{\mathbb Q}=\Theta^{g-1}/(g-1)!$ in Chow.

The cohomological Fourier transform is an integral isomorphism
$H^2(A,\mathbb Z)\to H^{2g-2}(\widehat A,\mathbb Z)$,
with the appropriate Tate twist.  Lefschetz $(1,1)$ and the
correspondence $\Gamma$ therefore give $2\Rightarrow3$.
Assertion~3 implies~1 because $\gamma_{g-1}$ is an integral
Hodge class.  The transpose Fourier correspondence, followed
by inversion and its usual sign, gives an inverse on
cohomology and preserves algebraic classes integrally.
Consequently $(-1)^g\mathcal F$ identifies the intersection
and Pontryagin rings on their algebraic cohomology images.
It sends positive codimension to positive dimension.  The
Pontryagin divided powers just cited transport to the unique
operations $x\mapsto x^n/n!$ in the torsion-free cohomology
ring.  Thus $2\Rightarrow4$.  Taking $x=\theta$ and
$n=g-1$ gives $4\Rightarrow1$.

For the last, stronger statement use the same construction
with a lift of $\Theta^{g-1}/(g-1)!$ in rational Chow.
Now the exponential identity holds in rational Chow, so
$\Gamma_{\mathbb Q}=\operatorname{ch}(\mathcal P_A)$.
Conversely a weak integral transform sends a symmetric
polarization divisor to an integral cycle whose rational
class is $(-1)^{g-1}\Theta^{g-1}/(g-1)!$.
This is exactly the additional condition in Theorem~3.8.
The cohomological construction alone supplied no such
rational-Chow equality.
\end{proof}

\begin{proposition}[AG-DP-ALG: the exact algebraic denominator of a divided power]
For a principally polarized complex abelian variety define
\[
 \kappa_c(A)=\min\{N\ge1:N\gamma_c\in\mathcal A^{2c}(A)\},
 \qquad 0\le c\le g.
\]
This integer exists and divides $c!$.
There is an integral correspondence in
$\mathrm{CH}^{g+c}(A\times A)$ inducing
$N L^{[c]}:x\mapsto N\gamma_c\cup x$ on \emph{all}
cohomology if and only if $\kappa_c(A)\mid N$.
Thus the exact algebraic denominator of this whole
raising operator is $\kappa_c(A)$, even though its
cohomological matrix is integral.
\end{proposition}
\begin{proof}
The integers $N$ with $N\gamma_c$ algebraic form a nonzero
ideal of $\mathbb Z$: it contains $c!$, since
$c!\gamma_c=\theta^c$ is algebraic.  Its positive generator
is $\kappa_c(A)$.  If $C$ is the stipulated integral
correspondence, apply it to $1=\operatorname{cl}([A])$.
The result is the class of the integral cycle $(p_2)_*C$,
and equals $N\gamma_c$.  This proves necessity.
Conversely choose an integral cycle $Z$ with class
$N\gamma_c$.  The integral correspondence $\Delta_*Z$
acts by cup product with that class, by the projection
formula.  This proves sufficiency.  The statement concerns
the operator on all degrees; it does not infer an obstruction
to an inverse restricted to a particular pair of degrees.
\end{proof}

\paragraph{AG-DP-EGS: a proved obstruction to the blanket geometric extension.}
We import the precise primary theorem of Engel--de Gaay
Fortman--Schreieder, \emph{Matroids and the integral Hodge
conjecture for abelian varieties}, arXiv:2507.15704v3
(27 March 2026; manuscript dated 31 March): Theorem~1.1
on p.~2, with the final proof in section~8.3, pp.~79--80.
For a very general principally polarized complex abelian
variety of dimension $g\ge4$, every integral algebraic
class in codimension $2\le c\le g-1$ is an even multiple
of $\gamma_c$.  Their Corollary~1.2, on p.~2 with proof
on p.~80, says the image is exactly $2\mathbb Z\gamma_c$
when $g=4$ or $5$.  This source theorem is a cited input:
the final reductions and their stated antecedents were read,
but the full 85-page degeneration/combinatorial proof and
its computational certificates were not independently audited.

It follows that $\kappa_c(A)$ is even in this range.
In particular
\[
 \kappa_2(A)=2\quad(g\ge4),\qquad
 \kappa_c(A)=2\quad(g\in\{4,5\},\ 2\le c\le g-1).
\]
The first equality also uses $\kappa_2(A)\mid2$;
the second is precisely their exact cycle-image result.
AG-DP-ALG therefore gives the exact factor two obstruction
to algebraic divided-power raising in these cases.
AG-FOURIER-MIN rules out an integral Fourier transform even
up to homology on these very general varieties.  In particular
it also rules out a weak integral Chow Fourier transform.
These are formal consequences with their full proofs just
given; they do not reprove the imported theorem.

Every principally polarized complex abelian variety has the
same integral polarized exterior cohomology model
$\bigwedge^*\mathbb Z^{2g}$: a symplectic basis writes
$\theta=\sum_i u_i\wedge v_i$.  Accordingly the Wilson
matrices and all the integral groups in AG-LEFSCH-SMITH
also apply to its cohomology, whether or not the pairs
come from elliptic factors.  Yet $\kappa_2=1$ on a product
of elliptic curves, where the subset cycles are integral,
and $\kappa_2=2$ on a very general variety of the same
dimension $g\ge4$.  This is a geometric obstruction to
transporting integral cohomological divided powers to
algebraic correspondences solely from that lattice.
It addresses the blanket extension of HodgeSketch Route~4,
not H8's correct exterior-algebra calculation.  The
very-general complex statement supplies no counterexample
on the CM locus or over $\overline{\mathbb Q}$.

\begin{proposition}[AG-FOURIER-ISOGENY: isogeny degrees annihilate the obstruction]
Let $A$ be a complex abelian variety of dimension $g$,
and suppose there is an isogeny
$f:A\times B\to J$ of degree $d$, where $J$ is a product
of Jacobians of smooth projective curves.  Then
\[
 d\,\mathcal Z_1(A)=0,\qquad
 \mathcal Z_1(A)=\frac{H^{2g-2}(A,\mathbb Z)\cap H^{g-1,g-1}(A)}
                         {\operatorname{cl}\mathrm{CH}_1(A)}.
\]
For principally polarized $A$, in particular
$\kappa_{g-1}(A)\mid d$.  Several such isogenies, with
possibly different $B$ and $J$, give the greatest common
divisor of their degrees as an annihilator.  In particular,
if that greatest common divisor is one and $A$ is principally
polarized, all the equivalent conditions of AG-FOURIER-MIN
hold.  This isolates the finite-prime degree gate in attempts to obtain the
integral Fourier property by adjoining Jacobian factors
and passing through isogenies.
\end{proposition}
\begin{proof}
Products of Jacobians satisfy the integral Hodge conjecture
for one-cycles by Beckmann--de Gaay Fortman, Theorem~1.2
and Corollary~4.1, whose proof uses the actual Abel--Jacobi
curve class and the minimal-class criterion.
Let $i:A\hookrightarrow A\times B$ use the origin of $B$,
and let $q:A\times B\to A$ be projection.  For an integral
Hodge class $\alpha$ of degree $2g-2$, the class
$f_*i_*\alpha$ is an integral one-cycle Hodge class on $J$.
Choose an integral cycle $Z$ representing it.
The cycle $q_*f^*Z$ is integral, since the isogeny is finite
flat.  On integral cohomology
$f^*f_*\beta=d\beta$: the left side is the sum of
translations by the kernel, and translations act trivially
on cohomology.  Also $q_*i_*=1$.  Hence
$\operatorname{cl}(q_*f^*Z)=d\alpha$.
The asserted divisibility and greatest-common-divisor
claim follow, the latter by an integer Bezout combination.
For the very general varieties in AG-DP-EGS, all these
degrees must be even.  This recovers the precisely scoped
obstruction in Engel--de Gaay Fortman--Schreieder,
Corollary~1.5, p.~3, whose proof is on p.~81; it does not
exclude the separate product/CM families.
\end{proof}

\begin{proposition}[AG-FOURIER-CM2: the full motivic identities need factor two]
Let $E:y^2=x^3-x$ over $\mathbb Q$, with origin $o=\infty$.
The modified diagonal
\[
 \Delta_e=\Delta_{123}-\Delta_{12o}-\Delta_{1o3}-\Delta_{o23}
                +\Delta_{1oo}+\Delta_{o2o}+\Delta_{oo3}
       \in\mathrm{CH}_1(E^3)
\]
has exact order two, even after base change to $\mathbb C$.
Here $\Delta_{12o}$ denotes the image $t\mapsto(t,t,o)$,
and the other terms use the indicated coordinates.

There is a weak integral Fourier correspondence $F_0$
for the principal polarization, with the exact square
identity $F_0^2=-[-1]^*$ in integral Chow correspondences.
However, no choice of integral Fourier correspondence can
satisfy both that square identity and
\[
 F\circ m_* = \Delta^*\circ(F\otimes F)
\]
exactly while inducing the usual rational Chow Fourier
transform.  The least positive factor $d$ for which
both identities can hold after multiplication by $d$
is $2$, over each of $\mathbb Q$, $\overline{\mathbb Q}$
and $\mathbb C$.  The canonical $F_0$ already realizes
this factor over $\mathbb Q$; its product-compatibility
defect has exact order two in the group of integral
Chow correspondences.
\end{proposition}
\begin{proof}
This is an explicit CM specialization of Moonen--Polishchuk,
Proposition~2.5, Remark~2.6 and Theorem~3.11
(author pp.~10--11 and 17--19).  We give the full cycle
detection because it records information invisible in
cohomology.

Use the integral projectors $\pi_0,\pi_1,\pi_2$ from
AG-LEFSCH-CYCLES.  The double cover $x:E\to\mathbb P^1$
is branched at $o$.  Pulling back the diagonal relation
on $\mathbb P^1\times\mathbb P^1$ gives
\[
 [\Delta_E]+[\Gamma_{[-1]}]=2\pi_0+2\pi_2.
\]
Thus $[-1]$ acts as $-1$ on the $\pi_1$ motive, integrally.
Applying $\pi_1^{\otimes3}$ to the small diagonal gives
exactly $\Delta_e$.  To check this directly, each
$\pi_2$ collapses its coordinate to $o$ and gives the
seven displayed terms; collapsing all three coordinates
pushes a curve to a point and gives zero.  A $\pi_0$
in one coordinate gives the corresponding coordinate
axis.  Applying $1-\pi_2$ in another coordinate kills
that axis, while two different $\pi_0$'s already give
zero.  Hence there are no additional terms.
Simultaneous inversion fixes the small diagonal and
commutes with these projectors, but acts as $(-1)^3$
on their tensor product.  Consequently $2\Delta_e=0$
over $\mathbb Q$.

For nonvanishing work over $\mathbb Q(i)$ or $\mathbb C$.
Put $P=(0,0)$ and let $\alpha(x,y)=(-x,iy)$.
Define $f:E^2\to E^3$ by $f(a,b)=(a,\alpha(a),b)$,
and set $T=(\operatorname{pr}_2)_*f^!$ on these
codimension-two cycles.  The map $f$ is a regular
immersion, the zero fiber of the smooth homomorphism
$(x,y,z)\mapsto y-\alpha(x)$; its normal line bundle
is trivial.  The kernel of $1-\alpha$ consists of
$o,P$, each with multiplicity one.  The seven
intersections therefore give
\begin{align*}
 T(\Delta_{123})&=[o]+[P],&T(\Delta_{12o})&=2[o],\\
 T(\Delta_{1o3})&=[o],&T(\Delta_{o23})&=[o],\\
 T(\Delta_{1oo})&=[o],&T(\Delta_{o2o})&=[o],\\
 T(\Delta_{oo3})&=0.
\end{align*}
For the final term the curve is contained in the
immersion, so the refined intersection is the first
Chern class of the trivial excess normal line and
is zero.  The others are the simple kernels of
$1-\alpha$, $\alpha$, or the identity just described.
It follows that $T(\Delta_e)=[P]-[o]$.
This is a nonzero element of exact order two in
$\mathrm{Pic}^0(E)\simeq E$: $P\ne o$ and
$\operatorname{div}(x)=2[P]-2[o]$.
Thus $\Delta_e$ is nonzero even over $\mathbb C$;
together with the preceding identity this proves
its exact order, also over $\mathbb Q$.

The weak transform and its square can be seen directly.
With $U=[o\times o]$ and $V=[E\times E]$, put
\[
 F_0=V-U-\pi_1\in\mathrm{CH}(E^2).
\]
The normalized Poincare bundle under the principal
polarization has $c_1(\mathcal P_E)=-\pi_1$;
the same product intersections give
$c_1(\mathcal P_E)^2=-2U$.
Thus $F_{0,\mathbb Q}=\operatorname{ch}(\mathcal P_E)$.
The exact matrix-unit relations from AG-LEFSCH-CYCLES
yield
\[
 F_0^2=-\pi_0+\pi_1-\pi_2=-[-1]^*.
\]
To compute the product defect over $\mathbb Q$ itself,
put $\epsilon=U+\pi_0+\pi_2-[\Delta_E]-V$ and let
$j_{12},j_{23}:E^2\to E^3$ be the coordinate inclusions.
Let $\delta(x,y)=(x,y,x)$.  We claim the exact identity
\begin{equation}
 j_{12*}\epsilon\star j_{23*}\epsilon-\delta_*\epsilon
       =\Delta_e.                                      \label{ag:eq:cm-cube-defect}
\end{equation}
Here $\star$ is Pontryagin product on $E^3$.
For completeness, the calculation is integral and
componentwise.  Write $A_i$ for the three coordinate
axes, $D_{ij}$ for the two-coordinate diagonals, and
$S_{ij}$ for the coordinate surfaces.  Set
$\alpha=A_1+A_2-D_{12}$ and
$\beta=A_2+A_3-D_{23}$.  Then
\[
 j_{12*}\epsilon=[o,o,o]+\alpha-S_{12},\qquad
 j_{23*}\epsilon=[o,o,o]+\beta-S_{23}.
\]
The zero-dimensional components in \eqref{ag:eq:cm-cube-defect}
cancel, and the one-dimensional difference is
\[
 \alpha+\beta-(D_{13}+A_2-\Delta_{123})=\Delta_e.
\]
To check dimension two, for an integral linear form
$u$ in the three coordinates write $K_u=u^*[o]$.
The only nonzero terms in $\alpha\star\beta$ after
cancellation are
\[
 \alpha\star\beta=K_y-K_{y-z}-K_{x-y}+K_{x-y+z}.
\]
For example $A_1\star A_3=K_y$,
$A_1\star D_{23}=K_{y-z}$,
$D_{12}\star A_3=K_{x-y}$, and
$D_{12}\star D_{23}=K_{x-y+z}$; each map onto
the displayed surface has degree one.  The other
nonzero terms $S_{12},S_{23}$ cancel in pairs.
The dimension-two difference is therefore
\[
 K_{x-y+z}-K_{x-y}-K_{y-z}+K_{x-z}+K_y-K_x-K_z=0.
\]
This is an equality in $\mathrm{Pic}(E^3)$ over
$\mathbb Q$: apply the theorem of the cube to the
symmetric line bundle $\mathcal O_E(o)$ and the
three maps $x,-y,z$, and use
$K_{x+z}+K_{x-z}=2K_x+2K_z$ from the theorem of the
square.  Thus no division by two has been made.
In dimension three,
$\alpha\star S_{23}=0=S_{12}\star\beta$:
in each expression two degree-one maps onto $E^3$
cancel, while the remaining term has a smaller image
and pushes forward to zero.  The dimension-four
term is zero by dimension.  This proves
\eqref{ag:eq:cm-cube-defect} over $\mathbb Q$.

Here is how this exact cycle identity measures the
motivic identity, including its sign.  For abelian
varieties $X,Y$ and $\eta\in\mathrm{CH}(X\times Y)$,
write
\[
 G_\eta=j_2^*\circ(\eta\star-)\circ j_{1*}.
\]
Moonen--Polishchuk Lemma~3.6, whose proof is valid
over an arbitrary field, identifies $G_\eta$ with
the integral correspondence $([-1]_X\times1_Y)^*\eta$;
thus $\eta\mapsto G_\eta$ is an isomorphism of
abelian groups.  Direct substitution gives
$F_0=-G_\epsilon$.  Permuting the second and third
coordinates in \eqref{ag:eq:cm-cube-defect}, using the
symmetry of $\epsilon$ and $\Delta_e$, gives
\[
 j_{13*}\epsilon\star j_{23*}\epsilon
       -(\Delta\times1)_*\epsilon=\Delta_e.
\]
The projection formula for Pontryagin product then
computes the entire dual product defect as
\begin{equation}
 B:=F_0\circ\Delta^*+
       m_*\circ(F_0\otimes F_0)=G_{\Delta_e},       \label{ag:eq:cm-dual-defect}
\end{equation}
where on the right $X=E^2$ and $Y=E$.
Indeed the first term on the left is
$-j_3^*\circ(((\Delta\times1)_*\epsilon)\star-)
\circ j_{12*}$, and the second is the same expression
with positive sign and kernel
$j_{13*}\epsilon\star j_{23*}\epsilon$.
Consequently $B$ has exact order two, over $\mathbb Q$
and after the stated base changes.

Finally put $I=[-1]^*$ and
$P=F_0\circ m_*-\Delta^*\circ(F_0\otimes F_0)$.
The already proved equality $F_0^2=-I$ makes $F_0$
invertible integrally, with inverse $-IF_0$.
Inversion commutes with $F_0$ and is compatible
with the multiplication and diagonal maps, so
\[
 P=I F_0\,B\,(F_0\otimes F_0).
\]
The identity follows by substituting
$B=F_0\Delta^*+m_*(F_0\otimes F_0)$ and
$F_0^2=-I$.  The outer compositions are integral
isomorphisms.  Thus $P$ too has exact order two:
$F_0$ satisfies the required product identity
up to factor two already over $\mathbb Q$.

To exclude another integral lift with factor one,
we use the precise classical criterion:
Moonen--Polishchuk, Definition~3.4 and Theorem~3.11,
say that, for an elliptic curve over an algebraically
closed field and any isomorphism to its dual, a motivic
integral Fourier transform up to factor $d$ exists
if and only if $d\Delta_e=0$.  The two identities
in the statement are exactly their conditions (ii)
and (iii), together with their rational-Chow
condition (i).  Since $\Delta_e$ has exact order two,
this proves that no other integral lift or isomorphism
to the dual can make $d=1$ work over $\mathbb C$.
Such a lift over $\overline{\mathbb Q}$ or $\mathbb Q$
would give one over $\mathbb C$ by base change, so
it is impossible there as well.  The explicit
calculation above provides factor two over $\mathbb Q$,
without extending the necessity theorem beyond its
stated field hypothesis.  This proves all the asserted
minimal factors; the general criterion remains
properly credited to its source.
\end{proof}

The last result does not contradict AG-LEFSCH-CYCLES:
those inverse-map identities were proved using an exact
matrix subring, while the Fourier assertion also requires
compatibility with the multiplication correspondence.
Nor does it obstruct the Fourier--Mukai derived equivalence.
It locates a nonzero integral Chow torsion class that must
be retained if HodgeSketch Route~4 or cote~16 is interpreted
as a strict motivic Fourier comparison.  Rational motives,
integral Betti cohomology, and the associated $L$-functions
do not detect this class; the relevant integral cohomology
is torsion-free.  The zero-location inputs A and M do not
enter this calculation of the exact compatibility factor.

\subsection{Archive and run consequences, with the original scope retained}
\label{ag:archive}

\paragraph{AG-C1.} Cote 1/79--80 and 111--112 gives the trace/Fredholm
framework, and SA23 geo-H-lidskii/geo-H-genus is its reported modern
reconstruction. A yields the earlier zeta positivity realization;
M yields AG-DET for the stated Hecke functions. AG-MOMENT and AG-HANKEL
give further consequences in that framework. MT-M29 now derives the
AG-DET operator from the canonical image-power grading of the locally
finite part of the actual adelic summation quotient, retaining the
zero multiplicities. Its positivity still uses A or M; an independent
geometric positivity or GGP-period comparison remains a separate
construction.

\paragraph{AG-C16.} Cote 16, scans 14--18, no.~7, concerns the integral
Lefschetz correspondence ring. Scans 44--50 contain the Coates letter's
conditions A, C and index positivity; the attached Hodge sketch uses a scan
numbering one higher than the GCT\_NSRH intake. These are algebraicity and
integrality problems. Neither A nor M constructs the required
correspondence or bounds its denominator. The sketch's explicit route is
the rank comparison on $X\times X$, not a claim that zero location
reverses Grothendieck's implication to the Weil conjectures.
AG-FROB-PROJECTOR/ORDER/COLLISION now give exact denominators
for the grading inside the Frobenius algebra; AG-FROB-CHOW
promotes it to rational Chow projectors for Kimura-finite
motives, and AG-FROB-NIL/CYCLES state the precise integral
nilpotence and torsion conditions. Those constructions do not by themselves give inverse Lefschetz.
AG-LEFSCH-SMITH/CYCLES now construct actual inverse
correspondences, with the exact least denominator, for
products of pointed elliptic curves with the product
polarization.  AG-FOURIER-MIN now locates the exact
minimal-class gate for an integral cohomological Fourier
kernel, while AG-DP-ALG isolates algebraicity of each
whole divided-power operator.  The primary theorem
imported in AG-DP-EGS disproves the blanket integral
extension on very general complex principally polarized
abelian varieties, without asserting a counterexample on
the CM locus or over $\overline{\mathbb Q}$.
AG-FOURIER-CM2 gives a separate arithmetic obstruction:
on $y^2=x^3-x$ a full integral motivic Fourier product
identity requires exactly factor two.  Neither obstruction
invalidates the constructed product inverse Lefschetz maps,
and the particular theta-chain comparison remains open.

\paragraph{AG-C34-74 and AG-C34-64.} Original PDF p.~74, manuscript p.~73, (2.5.5),
has now been visually checked: its multiplicative $1$ denotes the trivial
invariant group, equivalently additive $0$. The finite-field curve theorem
is the stated input, so inherited Q04 is unchanged. The original
definitions and full proof route for Theorem 1.1 have also been read.
AG-C34-TYPE identifies the literal obstruction in Q03, and
AG-C34-REPAIR proves the coherent intrinsic correction. AG-C34-INDEX
gives its finite Frobenius/Jacobian index. AG-C34-SECTION also proves
the formulation using a specified section and solvable $H$;
AG-C34-NONSOLV gives its abstract group-theoretic boundary.
No assertion of new purity or
of a consequence of A/M is made.

\paragraph{AG-C44.} Cote 44 Theorems 4--5 and the Serre letter, at the
locators in AG-SIGN, concern polarization and heights. AG-SIGN corrects
the central-derivative bridge but leaves cycle existence and the
height formula explicit. The supplied Hodge sketch also cites a separate
2026 CM-Hodge theorem for its abelian positivity consequences; those
claims are not consequences of either A or M and are carried with
their original dependency rather than independently certified here.

\paragraph{AG-C149 and AG-GT.} ORDER\_tonight, item 31, points to cote
149 scans 79--85, LM6 A.44--A.45, integral bridge B3, and the
irregular-prime/$\ell=2$ questions. GRH does not remove the integral
class-group and generator-relation gates in AG-LM6-IRREG/NORM.
For an actual punctured normal bundle, AG-C149-COMP now gives the
classical comparison under the explicit $\pi_2=0$ hypothesis, and
AG-C149-ELLIPTIC separates its global class from every section's
obstruction. AG-C149-NORMAL/CUSP now supply the actual normalizer
comparison in a family of four-punctured projective lines; every
arithmetic section even lifts through the full profinite normalizer.
A different specified family still requires its own comparison map.
The semi-graded Euler information in AG-EXT explains why the endpoint
zero theorems do not evaluate that extension data. In SA22 run 1,
A6-node-character (AG-NODE) is the quadratic Artin character times the cyclotomic
character; A gives effective detection of that finite quadratic sign.
M additionally covers finite nonabelian splitting fields in effective
Frobenius tests such as AG-FROB. These improvements do not reverse the
recorded refutation AG-INERTIA of run 2 A:inertia-maximality: finite-index inertia
subgroups have the same weights. Nor do they identify the profinite GT
kernel from the cyclotomic quotient. The existing supplied claim about
rational Deligne--Drinfeld freeness (CC-GT03, external 2026 input) is
carried separately from the full profinite Ihara question R-GTHAT.

\paragraph{AG-PURITY and AG-WM.} Finite-field purity and the usual
Grothendieck--Lefschetz trace formula are already classical. The full
geometric weight--monodromy statement remains represented in these
materials by BR-WM-02's finite exterior-power trace criterion; its
remaining inequalities are geometric hypotheses, not consequences of
global Dirichlet or Hecke zero locations. Likewise AG-STANDARD records that the supplied archive
questions about standard conjectures require algebraic cycles or
correspondences. No formal logical independence from A or M is asserted;
the precise point is that the displayed deductions do not supply those
inputs.
